\documentclass[11pt]{ctexart}
\usepackage[a4paper,margin=1in]{geometry}
\usepackage{graphicx}
\usepackage{xcolor}
\usepackage{hyperref}
\definecolor{refgray}{gray}{0.45}
\title{\textbf{市场最新观点汇总}\\[0.4em]{\large 美联储鹰派转向与地缘缓和预期交织，全球资产定价锚点切换，市场聚焦AI结构性需求与供给侧约束分化}}
\author{KC桌面 自动生成}
\date{260619}
\begin{document}
\maketitle
\tableofcontents
\newpage
\section{一页摘要}
\begin{itemize}
  \item 美联储6月FOMC鹰派意外超预期，点阵图暗示加息，但新主席Warsh的鸽派措辞与‘小数点左侧’通胀观对冲了部分鹰派信号，市场正从‘数点’转向‘听词’，政策不确定性溢价上升。
  \item 全球资金流入创历史纪录，美国股票单周流入1192亿美元，科技股集中度触及历史泡沫区间，BofA牛熊指标进入极端看多区域，市场定价核心从经济增长切换为政治周期与收益率曲线的博弈。
  \item 中国内需数据走弱但非失速，政策定力为‘微调而非转向’，房地产量稳价跌、消费盈利预期下修尚未触底，结构性分化加剧，选股逻辑需从‘买赛道’切换为‘买护城河’。
  \item AI需求从训练向推理迁移，定制芯片(ASIC)占比预计2027年超GPU，半导体设备市场增速大幅上调，洁净室瓶颈缓解，存储器投资权重跃升至约50\%，日本股市盈利驱动范式确立。
  \item 大宗商品定价逻辑从需求预期切换为供给约束，铝市场面临中东减产与亚洲增产的时间错配，铜和黄金受地缘风险折价过度，霍尔木兹海峡复航或成战术催化剂。
\end{itemize}
\section{宏观与利率：美联储沟通范式切换与全球央行分化}
\textbf{美联储正经历格林斯潘时代以来最彻底的沟通方式变革，前瞻指引系统性削减，市场需将更多权重放在每一个事件日上，点阵图参考价值下降，Warsh的措辞成为更重要的政策风向标。与此同时，全球央行言行一致性显著分化，ECB用行动强化鹰派可信度，BOE则削弱其信号威慑力。}

\begin{itemize}
  \item 6月FOMC点阵图中位数显示2026年利率升至3.75\%，比市场预期更鹰派，但Warsh将点阵图形容为‘铅笔画的’，并强调‘一半同事认为利率应维持或更低’，这种鹰派数据与鸽派叙事的组合削弱了点阵图作为政策信号的市场权重。
  \item Warsh提出‘关注小数点左侧’的通胀观，暗示2.0\%-2.9\%的通胀区间均符合2\%目标，这一表述可能系统性降低加息门槛，若被正式或非正式采用，将重塑长端利率和美元的定价逻辑。
  \item JPM的NLP模型将6月FOMC政策声明评为一年来最鹰派，声明完全重写，不再提供前瞻指引，转而采用‘委员会将实现价格稳定’的宣言式表述，回归格林斯潘时代的含蓄与不可预测性。
  \item 从SOFR期货期权推断的隐含概率分布显示，不仅整体向鹰派方向移动，且‘双边尾部均增强’，意味着市场认为极端情景（更鹰派或更鸽派）的概率都在上升，政策不确定性被重新定价。
  \item NOM编制的ECBspeak指数升至20.7，为2022年3月以来最高，但这次欧洲央行在鹰派沟通后迅速跟进加息，言行一致推动NOM上调加息预测，而BOE鹰派指数仅4.4且按兵不动，央行沟通本身正成为独立的政策工具。
  \item BofA牛熊指标升至9.2，进入‘极端看多’区间，自2002年以来触发17次，平均回撤2-3\%，最大回撤15-20\%，命中率约60\%，当前资金流向的集中度已触及历史泡沫的‘天花板’。
\end{itemize}
\begin{figure}[htbp]
\centering
\includegraphics[width=0.88\linewidth]{figures/fig_060.jpg}
\caption{Figure 1 - Figure 1: The 2026 dots show 12.5bp of rate hikes, and 6 of the 18 dots projected two or more hikes Distribution of projected midpoint of fed funds target range for 2026, March 2026 vs. June 2026; number of participants}
\end{figure}
\begin{figure}[htbp]
\centering
\includegraphics[width=0.88\linewidth]{figures/fig_061.jpg}
\caption{Figure 2 - Figure 2: Both the statement and the press conference were scored hawkishly Fed Hawk-Dove Score\textbackslash{}* of FOMC statements and prepared press conference remarks; Index}
\end{figure}
\begin{figure}[htbp]
\centering
\includegraphics[width=0.88\linewidth]{figures/fig_062.jpg}
\caption{Figure 2 - Figure 3: Post-FOMC, the implied distribution for Dec '26 shifted hawkishly, and both tails strengthened, indicating higher policy uncertainty Probabilities of different fed funds outcomes for Dec '26 calculated from the risk-neut}
\end{figure}
\begin{figure}[htbp]
\centering
\includegraphics[width=0.88\linewidth]{figures/fig_063.jpg}
\caption{Figure 4 - Figure 4: Inflation markets shifted hawkishly with rising real yields and falling breakevens Daily change from 6/16 to 6/17 in 5-, 10-, and 30-year breakevens and real yields, bps}
\end{figure}
\begin{figure}[htbp]
\centering
\includegraphics[width=0.88\linewidth]{figures/fig_064.jpg}
\caption{Figure 5 - Figure 5: The WAM of the Fed's SOMA remains significantly longer than that of the overall Treasury market}
\end{figure}
{\small\color{refgray}\textbf{References:} [R001] GS：外汇市场的真正拐点不是地缘冲突缓和，而是美联储的鹰派姿态切换; [R003] NOM：央行“叫得凶”和“咬得狠”正在分岔，市场低估了这种分化; [R006] NOM：鲍威尔之后，美联储正经历一场更深刻的权力重构; [R009] JPM：市场低估的不只是加息幅度，更是货币政策沟通机制的历史性转折; [R031] 美国银行：市场正在定价的不是繁荣，而是繁荣之后的“谁先离场”}

\section{AI与算力：从训练到推理的架构迁移与定制芯片崛起}
\textbf{AI基础设施投资正从训练主导转向推理驱动，定制芯片(ASIC)凭借系统级TCO优势，预计到2027年出货量占比将首次超过GPU，半导体设备市场增速大幅上调，洁净室瓶颈缓解为2027年后的需求释放铺平道路。}

\begin{itemize}
  \item JPM将2026年WFE市场增速预测从21\%大幅上调至28\%，2027年从18\%上调至29\%，并首次给出2028年16\%的增长指引，核心驱动力是AI需求从训练向推理迁移，内存芯片投资权重从个位数跃升至约50\%。
  \item 全球前四大美国云服务商2026年资本开支增速从63\%上调至80\%，2027年从40\%上调至50\%，2026年同比增量将超过2500亿美元，为有史以来最大年度增幅，北美规划电力容量已超175GW。
  \item 到2027年，AI加速器出货量中定制芯片(ASIC/XPU)占比将从2025年的32\%跃升至53\%，首次超过GPU，Google TPUv7在每美元算力上达GPU的2.7倍，Broadcom占据80-85\%高端ASIC市场份额。
  \item 洁净室产能瓶颈正在缓解，通过工厂空间的创造性利用和洁净室扩建的正式启动，此前市场对2026年设备交付的担忧正在消散，从2027年起需求将迎来真正的释放窗口。
  \item 台积电CoWoS产能规划从2026年底11.5万片/月扩至2028年22万片，3D SoIC产能从2027年底4万片/月扩至2028年6.5万片，客户需求从NVIDIA扩展到AMD、谷歌、博通等。
  \item BARC预计2026-2027年全球光学网络市场维持两位数增长，电信与云双引擎驱动，城域市场ZR/ZR+可插拔模块是核心推手，2027年市占率预计达51\%，Ciena在云市场市占率预计达78\%。
\end{itemize}
\begin{figure}[htbp]
\centering
\includegraphics[width=0.88\linewidth]{figures/fig_101.jpg}
\caption{Figure 1 - Figure 1: Shipment value of semiconductors and SPE (YoY)}
\end{figure}
\begin{figure}[htbp]
\centering
\includegraphics[width=0.88\linewidth]{figures/fig_102.jpg}
\caption{Figure 2 - Figure 2: SPE shipment trend by region}
\end{figure}
\begin{figure}[htbp]
\centering
\includegraphics[width=0.88\linewidth]{figures/fig_103.jpg}
\caption{Figure 3 - Figure 3: Global SPE shipment trend excluding China}
\end{figure}
\begin{figure}[htbp]
\centering
\includegraphics[width=0.88\linewidth]{figures/fig_104.jpg}
\caption{Figure 4 - Figure 4: Tokyo Electron SPE sales}
\end{figure}
\begin{figure}[htbp]
\centering
\includegraphics[width=0.88\linewidth]{figures/fig_161.jpg}
\caption{FIGURE 1 - BCI, US expect DD growth in the cloud vertical in CY26-CY27, as AI-related cloud build-outs broaden and continue to accelerate. FIGURE 1. Optical Market Expected to Enjoy DD Growth Trends, Driven by Telco Recovery and AI}
\end{figure}
\begin{figure}[htbp]
\centering
\includegraphics[width=0.88\linewidth]{figures/fig_162.jpg}
\caption{FIGURE 2 - - We currently do not forecast meaningful coherent Optical inside the data center in CY26 (potentially could see 2H26), and believe that to be more of a CY27+ opportunity. - We continue to expect the total Optical market to have a LSD CAGR from CY22-27, taking}
\end{figure}
\begin{figure}[htbp]
\centering
\includegraphics[width=0.88\linewidth]{figures/fig_163.jpg}
\caption{FIGURE 3 - - We expect Sub to see +DD recovery in CY27 (on easy comps), after multiple down years through CY26. Sub is MSD-HSD\% of the total market. FIGURE 3. Despite Cloud Growth, Telco Still Majority of the Market}
\end{figure}
{\small\color{refgray}\textbf{References:} [R013] JPM：半导体设备市场的下一个结构性拐点不在需求，而在供给侧的产能释放节奏; [R014] JPM：AI定制芯片正在从“备选方案”变成“主导逻辑”; [R024] BARC：光学市场真正的拐点不在于AI需求本身，而在于供给结构正在经历不可逆的重塑}

\section{能源与大宗商品：供给约束下的时间错配与地缘风险折价}
\textbf{大宗商品定价逻辑正从需求预期驱动切换为供给约束驱动，铝市场面临中东减产与亚洲增产的时间错配，铜和黄金因地缘风险折价过度，霍尔木兹海峡复航或成战术催化剂，而地震勘探行业正进入有纪律的多年度上行周期。}

\begin{itemize}
  \item GS大幅下调中东铝产量预测，2026年和2027年分别下调约66万吨和100万吨，即使海峡重新开放，受损电解槽修复需约12个月，供给缺口将延续至2027年；同时上调印尼产量至170万吨和290万吨，中国产量至4560万吨和4630万吨。
  \item GS将2026年Q3和2027年LME铝价预测分别上调至3300美元/吨和2950美元/吨，但仍低于远期曲线，核心逻辑是亚洲供给增量最终覆盖中东损失，2026年全球短缺72万吨，2027年转为过剩59万吨。
  \item JPM认为市场对美伊局势缓和的潜在影响定价不足，自冲突以来紫金矿业H股下跌26\%，紫金黄金下跌49\%，而恒生指数仅下跌8\%，铜和黄金基本面未实质性恶化，下跌更多来自情绪层面的风险折价。
  \item 全球地震勘探行业正经历从‘资本纪律’到‘资源重置’的范式转换，EAGE 2026参会人数超6000人创历史新高，主要油企储量寿命从2012年的约55年骤降至2026年预计的约20年，前沿勘探重回议程。
  \item AI技术对地震勘探行业非但未削弱核心企业护城河，反而让高质量成像数据成为更稀缺的资产，AI无法解决基于物理的波动方程成像问题，HPC算力是结构性护城河。
  \item 中国钢铁市场长材表观消费周环比增长5.1\%，扁平材增长2.0\%，电炉开工率提升2个百分点至64.4\%，供给端电炉产能利用率回升驱动的再定价正在发生，但终端采购节奏仍偏谨慎。
\end{itemize}
\begin{figure}[htbp]
\centering
\includegraphics[width=0.88\linewidth]{figures/fig_170.jpg}
\caption{Exhibit 12 - Exhibit 1: We Raise our 2027 Price Forecast to \textbackslash{}\$2,950 But Our Forecast Stays Well Below the Forwards LME Aluminium Price}
\end{figure}
\begin{figure}[htbp]
\centering
\includegraphics[width=0.88\linewidth]{figures/fig_171.jpg}
\caption{Exhibit 2 - Exhibit 2: We Expect the Global Aluminium Market to Shift into a Smaller Surplus in 2027 Than Previously Forecast on Lower Middle East Supply Global Aluminium Market Balance}
\end{figure}
\begin{figure}[htbp]
\centering
\includegraphics[width=0.88\linewidth]{figures/fig_172.jpg}
\caption{Exhibit 3 - Exhibit 3: We Expect a Larger Q3 Deficit Before the Market Returns to Surplus in Q4 Primary Aluminium Market Balance}
\end{figure}
\begin{figure}[htbp]
\centering
\includegraphics[width=0.88\linewidth]{figures/fig_173.jpg}
\caption{Exhibit 4 - Exhibit 4: Aluminium Production Margins Are Approaching the 2022 Highs}
\end{figure}
\begin{figure}[htbp]
\centering
\includegraphics[width=0.88\linewidth]{figures/fig_174.jpg}
\caption{Exhibit 5 - Exhibit 5: Low Inventory Cover Supports Higher Smelter Margins Aluminium Inventories as Days of Consumption}
\end{figure}
\begin{figure}[htbp]
\centering
\includegraphics[width=0.88\linewidth]{figures/fig_175.jpg}
\caption{Exhibit 6 - Exhibit 6: Two Supply Shocks: Middle East Output Falls, Indonesia Supply Catching Up Aluminium Production Share}
\end{figure}
\begin{figure}[htbp]
\centering
\includegraphics[width=0.88\linewidth]{figures/fig_085.jpg}
\caption{Figure 1 - Figure 1: Stock Price Performance since US-Iran Conflict}
\end{figure}
\begin{figure}[htbp]
\centering
\includegraphics[width=0.88\linewidth]{figures/fig_077.jpg}
\caption{Exhibit 1 - Exhibit 1: While still early cycle, the change in tone suggests that exploration is moving back up the agenda for the majors, supporting a recovery in seismic demand over time}
\end{figure}
\begin{figure}[htbp]
\centering
\includegraphics[width=0.88\linewidth]{figures/fig_078.jpg}
\caption{Exhibit 2 - Exhibit 2: Industry reserve life was broadly flat at 10.2 years in 2024 Reserve life (year) per FAS69 data}
\end{figure}
\begin{figure}[htbp]
\centering
\includegraphics[width=0.88\linewidth]{figures/fig_079.jpg}
\caption{Exhibit 3 - Exhibit 3: Top Projects oil reserve life has fallen 35 years since 2012 Top Projects reserve life, by year of report and breakeven}
\end{figure}
\begin{figure}[htbp]
\centering
\includegraphics[width=0.88\linewidth]{figures/fig_160.jpg}
\caption{Exhibit 1 - Exhibit 2: Weekly steel demand}
\end{figure}
{\small\color{refgray}\textbf{References:} [R011] GS：地震勘探行业正在进入一个被低估的、有纪律的多年度上行周期; [R012] JPM：市场低估了供给侧的再定价，而非需求端的边际变化; [R023] 摩根斯坦利：中国钢铁市场的真实复苏信号，被周度数据的噪音掩盖了; [R033] GS：铝市场真正被低估的不是地缘溢价，而是两股供给冲击的时间错配}

\section{中国地产与消费：量稳价跌与盈利预期重置}
\textbf{中国房地产市场交易量回暖但价格持续探底，以价换量仍是主旋律；消费板块估值已跌至十年均值下方1.5个标准差，但80\%覆盖公司的盈利预测仍在被下调，成本端压力将在2026年下半年集中释放，盈利预期重置尚未触底。}

\begin{itemize}
  \item 30城新房成交量小幅回升且高于去年同期，16城二手房成交量持平于高位，但70城二手房价格指数5月继续下降，多个独立数据源显示价格已显著低于2025年甚至逼近2019年基准，量稳价跌背离反映积压需求释放而非信心修复。
  \item 5月社零同比-0.6\%，为2022年12月以来首次转负，耐用消费品是主要拖累项，汽车同比-16\%、家电-16\%、家具-9\%，房地产后周期消费传导机制断裂，以旧换新政策红利快速衰减。
  \item JPM指出消费板块盈利下修尚未触底，必需消费和可选消费2026年一致预期EPS分别较年初下调9\%和7\%，成本压力的充分定价可能要到2026年底才会显现，CPI-PPI剪刀差将扩大至约2.9个百分点。
  \item MS认为国内需求走弱但非失速，固定资产投资减速部分源于地方政府化解隐债的主动调控，零售同比转负受去年高基数影响，两年复合增速5月略好于4月，政策将‘微调而非转向’。
  \item 外卖补贴新规的核心机制不是设定价格上限，而是通过七天事前披露和资金来源报告让激进的补贴行为变得更慢、更透明，降低了美团的下行尾部风险，但尚未转化为正面盈利催化剂。
  \item 能量饮料市场正从品类渗透红利期转向存量博弈期，品类整体投影销售额增速从+15\%降至+10\%，购买频次出现2025年11月以来首次同比下滑(-2\%)，品牌分化加剧，Monster表现最抗跌。
\end{itemize}
\begin{figure}[htbp]
\centering
\includegraphics[width=0.88\linewidth]{figures/fig_007.jpg}
\caption{Exhibit 1 - Exhibit 1: 30-city daily property transaction volume in the primary market edged up over the last week and remained above year-ago level}
\end{figure}
\begin{figure}[htbp]
\centering
\includegraphics[width=0.88\linewidth]{figures/fig_008.jpg}
\caption{Exhibit 2 - Exhibit 2: 16-city daily property transaction volume in the secondary market was roughly flat over the last week but remained above year-ago level}
\end{figure}
\begin{figure}[htbp]
\centering
\includegraphics[width=0.88\linewidth]{figures/fig_009.jpg}
\caption{Exhibit 3 - Exhibit 3: NBS 70-city secondary home prices continued to decline in May}
\end{figure}
\begin{figure}[htbp]
\centering
\includegraphics[width=0.88\linewidth]{figures/fig_039.jpg}
\caption{Figure 1 - Figure 1: China retail sales YoY trend}
\end{figure}
\begin{figure}[htbp]
\centering
\includegraphics[width=0.88\linewidth]{figures/fig_040.jpg}
\caption{Figure 2 - Figure 2: China retail sales YoY trend – key categories}
\end{figure}
\begin{figure}[htbp]
\centering
\includegraphics[width=0.88\linewidth]{figures/fig_041.jpg}
\caption{Figure 3 - Figure 3: Retail sales yoy growth ranked by category – May 2026}
\end{figure}
\begin{figure}[htbp]
\centering
\includegraphics[width=0.88\linewidth]{figures/fig_141.jpg}
\caption{Exhibit 1 - Exhibit 1: Retail sales YoY weakened due to a high base and a softening job market}
\end{figure}
\begin{figure}[htbp]
\centering
\includegraphics[width=0.88\linewidth]{figures/fig_142.jpg}
\caption{Exhibit 2 - Exhibit 2: Firm pace of resolving hidden debt after a relatively slow start}
\end{figure}
\begin{figure}[htbp]
\centering
\includegraphics[width=0.88\linewidth]{figures/fig_143.jpg}
\caption{Exhibit 3 - Exhibit 3: Policy bank bond issuance has seen a net contraction YTD}
\end{figure}
\begin{figure}[htbp]
\centering
\includegraphics[width=0.88\linewidth]{figures/fig_144.jpg}
\caption{Exhibit 4 - Exhibit 4: Consumer confidence has dipped again after a period of modest improvement}
\end{figure}
{\small\color{refgray}\textbf{References:} [R002] GS：中国经济的真实温度，藏在每周高频数据的边际变化里; [R007] JPM：消费股真正的底部不在宏观数据里，而在企业盈利预期的重新校准; [R015] 摩根斯坦利：市场低估的不是下行风险，而是政策“微调”而非“转向”的定力; [R018] JPM：外卖补贴新规的真正价值不在上限，而在让烧钱竞争变得“无法隐身”; [R022] JPM：中国消费板块的真正底部不在估值，而在盈利预期的重置; [R026] JPM：能量饮料市场真正进入“买量竞争”阶段，品牌分化才刚刚开始}

\section{区域市场：日本盈利驱动与韩国K型复苏闭合}
\textbf{日本股市上涨核心驱动力从估值修复切换为盈利驱动，AI半导体超级周期叠加企业盈利能力结构性改善，TOPIX EPS预计连续两年两位数增长；韩国经济正从K型复苏转向全面反弹，科技出口量价齐升与消费超预期回暖形成共振。}

\begin{itemize}
  \item JPM将日经225年底目标价上调至75,000点，TOPIX目标价上调至4,400点，TOPIX每股收益预计2026财年同比增长11\%，2027财年再增长12\%，日经225中AI相关股票市值占比已接近50\%。
  \item 日本股市估值（PE）从高位回落到16倍，过热担忧减轻，半导体和半导体设备企业利润超预期增长将EPS预期从198日元拉升至234日元，盈利在撑盘而非纯情绪驱动。
  \item 韩国2026年GDP增速预计从2025年的1.1\%跃升至2.8\%，半导体出口占韩国总出口比重已达40\%，2026年5月日均出口同比增速达60.9\%创周期新高，DRAM和HBM的bit出货量预计维持两位数增长。
  \item 韩国消费复苏速度快于预期，财政支持、财富效应和工资增长共同推动，服务业生产增速预计在2026年初达到4\%，韩国央行预计从2026年7月开始在一年内加息4次，终端利率达3.50\%。
  \item 东盟电信板块超额收益的核心变量是竞争格局的结构性改善而非5G技术升级，泰国和新加坡以双寡头格局领先，TRUE与DTAC合并后竞争保持理性，ARPU止跌回升。
  \item 日本制药板块加权市盈率仅15.8倍，较美国同行折价约27\%，Daiichi Sankyo年初至今下跌25\%，但MS认为其ADC技术平台被全球顶级药企认可，管线价值被市场低估。
\end{itemize}
\begin{figure}[htbp]
\centering
\includegraphics[width=0.88\linewidth]{figures/fig_051.jpg}
\caption{Figure 7 - Figure 1: Target Prices for Japanese Equities}
\end{figure}
\begin{figure}[htbp]
\centering
\includegraphics[width=0.88\linewidth]{figures/fig_052.jpg}
\caption{Figure 3 - Figure 3: TOPIX EPS Forecast}
\end{figure}
\begin{figure}[htbp]
\centering
\includegraphics[width=0.88\linewidth]{figures/fig_053.jpg}
\caption{Figure 5 - Figure 5: TOPIX ROE Forecast}
\end{figure}
\begin{figure}[htbp]
\centering
\includegraphics[width=0.88\linewidth]{figures/fig_054.jpg}
\caption{Figure 2 - Figure 2: Japanese Equities Outlook to End-2026 (TOPIX, Nikkei 225) Note: Orange shading indicates the end-2026 target price. Figure 4: TOPIX P/E Forecast}
\end{figure}
\begin{figure}[htbp]
\centering
\includegraphics[width=0.88\linewidth]{figures/fig_055.jpg}
\caption{Figure 6 - Figure 6: 2026 Year-to-Date Performance}
\end{figure}
\begin{figure}[htbp]
\centering
\includegraphics[width=0.88\linewidth]{figures/fig_056.jpg}
\caption{Figure 7 - Figure 7: Sector Allocation Figure 8: EPS Estimates by TOPIX 17 Sectors}
\end{figure}
{\small\color{refgray}\textbf{References:} [R008] JPM：日本股市上涨的核心驱动力不是流动性，而是盈利能力的结构性跃迁; [R027] 摩根斯坦利：东盟电信股真正的回报来源不是5G，而是竞争格局的结构性改善; [R034] 摩根斯坦利：韩国K型复苏的裂口正在闭合，但市场低估了财富效应与政策节奏的共振; [R035] 摩根斯坦利：日本制药板块的错位定价正在为深度价值投资者创造窗口}

\section{中国医药与机器人：全球定价权再校准与‘大脑’路线未定}
\textbf{中国医药资产被低估的不是基本面而是全球定价权的再校准，临床数据和BD活动依然强劲，但政策不确定性压制短期可见度；机器人行业硬件成本优势已建立，但通往‘物理AI’的技术路线尚未收敛，远期市场空间线性外推可能建立在未选定路径上。}

\begin{itemize}
  \item GS参加全球医疗健康大会后判断，中国医药股年内平均下跌11\%但基本面强劲，市场真正担心的是中美BD交易监管审查和国内医疗反腐两个‘看不见’的政策变量，而非对行业长期价值的否定。
  \item Gilead透露约50\%的优先BD机会来自或关联中国，Pfizer认为到2030年中国biotech可能具备全球开发能力，中国创新不再是‘低成本替代’而是真正的创新来源，ADC和双抗领域出海逻辑清晰。
  \item 宇树科技2025年人形机器人出货5500+台全球第一，营收同比增长335\%至17亿元，毛利率约60\%，电机、减速器、激光雷达全部自研，外购部件仅占成本14-18\%，成本优势构成护城河。
  \item NOM调研指出行业当前最硬的约束条件已从‘能不能动’转向‘能不能想’，世界模型技术路线尚未收敛，VLA和世界模型两条路线还在竞争，硬件厂商可能被AI巨头逐步挤压。
  \item 梅卡曼德在大型3D抓取领域连续六年国内市场份额第一，核心壁垒在于‘眼-脑’层面的数据飞轮：部署越多数据越多，模型越强部署门槛越低，但灵巧手尚未形成规模化收入。
  \item 深度机器人（DEEP Robotics）2025年首次实现全年盈利，营收同比增长227\%至3.37亿元，毛利率52.83\%，电力巡检和消防场景从试点走向批量部署，但价格和性能仍是规模化复制的硬约束。
\end{itemize}
\begin{figure}[htbp]
\centering
\includegraphics[width=0.88\linewidth]{figures/fig_027.jpg}
\caption{Exhibit 1 - Exhibit 1: China-sourced BD deals surged to over 50 each quarter since 4Q25 Number of in-license/out-license BD and M\&A for China Biotech}
\end{figure}
\begin{figure}[htbp]
\centering
\includegraphics[width=0.88\linewidth]{figures/fig_028.jpg}
\caption{Exhibit 2 - Exhibit 2: Total deal sizes continue upward trend Total deal value of out-license BD and company acquisition activities for China Biotech (FY2022-26)}
\end{figure}
\begin{figure}[htbp]
\centering
\includegraphics[width=0.88\linewidth]{figures/fig_029.jpg}
\caption{Exhibit 3 - Exhibit 3: Therapeutic areas are diversifying in China out-license deals... Number of out-license BD deals for China Pharma/Biotech (FY2022—2026YTD), by therapeutic areas Unit: \# of BD deals}
\end{figure}
\begin{figure}[htbp]
\centering
\includegraphics[width=0.88\linewidth]{figures/fig_030.jpg}
\caption{Exhibit 4 - Exhibit 4: ...as well as in modalities where China has built visible clinical and platform depth Number of out-license BD deals for China Pharma/Biotech (FY2022—2026YTD), by modality}
\end{figure}
\begin{figure}[htbp]
\centering
\includegraphics[width=0.88\linewidth]{figures/fig_031.jpg}
\caption{Exhibit 5 - Exhibit 5: Number of China out-license BD deals has reached about a quarter of total number of deals globally China go-global as \% of global BD deals}
\end{figure}
\begin{figure}[htbp]
\centering
\includegraphics[width=0.88\linewidth]{figures/fig_032.jpg}
\caption{Exhibit 6 - Exhibit 6: China out-license BD deals have gained a significant proportion of total global deal size China go-global deal size as \% of global BD deal size}
\end{figure}
{\small\color{refgray}\textbf{References:} [R004] GS：中国医药资产被低估的不是基本面，而是全球定价权的再校准; [R028] NOM：四足机器人已越过盈亏线，但真正考验在于“大脑”尚未定型; [R029] NOM：宇树科技IPO揭示的真正悬念，不在硬件成本，而在“大脑”路线的未定; [R030] NOM：机器人视觉的竞争壁垒不在“手”，而在“眼”和“脑”的数据飞轮}

\section{半导体与功率器件：供给驱动的涨价周期与定制化浪潮}
\textbf{功率半导体涨价周期是供给驱动型而非需求驱动型，全球主要公司资本开支连续两年下降，供给侧理性收敛使定价权向供应商倾斜；半导体设备市场受益于AI需求结构拓宽，存储器投资权重飙升，洁净室瓶颈缓解为2027年后释放需求。}

\begin{itemize}
  \item 全球功率分立器件营收2026年4月同比增长16\%，MS判断本轮涨价周期是供给驱动型，全球主要功率器件公司资本开支连续两年下降，2026年仅温和回升约11\%，供给端扩张力度远弱于以往复苏周期。
  \item 中国功率器件市场未来三年新增产能有限，士兰微、华大半导体等企业正将资源转向AI电源管理芯片等更高附加值领域，而非继续扩大IGBT和MOSFET传统产能，只要需求不急剧恶化，涨价趋势有望延续至2026年下半年。
  \item 工业和汽车合计占功率器件需求的72\%，工业自动化需求强劲，头部企业1Q26收入同比增长21\%，中国电动车批发增速从4月零增长回升到5月的7\%，但光伏需求疲软，新增装机同比下降51\%。
  \item JPM将2026年WFE市场增速从21\%上调至28\%，2027年从18\%上调至29\%，AI需求从训练向推理迁移使存储器投资权重从个位数跃升至约50\%，三大存储器行业总资本开支预估从3000亿上调至4500亿美元。
  \item 洁净室产能瓶颈通过工厂空间创造性利用和正式扩建正在被打破，从2027年起需求将迎来真正释放窗口，台积电CoWoS产能规划从2026年底11.5万片/月扩至2028年22万片。
  \item 全球IT服务订单下滑的真正推手是地缘政治摩擦而非AI颠覆，Accenture 3Q新订单同比下降3\%，管理服务订单骤降15\%而咨询订单加速增长13\%，AI业务正从概念验证走向商业部署。
\end{itemize}
\begin{figure}[htbp]
\centering
\includegraphics[width=0.88\linewidth]{figures/fig_149.jpg}
\caption{Exhibit 1 - Exhibit 1: Total discrete revenue yoy growth turned positive since 4Q25}
\end{figure}
\begin{figure}[htbp]
\centering
\includegraphics[width=0.88\linewidth]{figures/fig_150.jpg}
\caption{Exhibit 2 - Exhibit 2: Auto and industrial accounted for 70\% of end demand in discrete semiconductors (2025)}
\end{figure}
\begin{figure}[htbp]
\centering
\includegraphics[width=0.88\linewidth]{figures/fig_151.jpg}
\caption{Exhibit 3 - Exhibit 3: China EV wholesale growth is turning incrementally positive}
\end{figure}
\begin{figure}[htbp]
\centering
\includegraphics[width=0.88\linewidth]{figures/fig_152.jpg}
\caption{Exhibit 4 - Exhibit 4: Industrial automation companies' revenue grew robustly at 21\% YoY in 1Q26}
\end{figure}
\begin{figure}[htbp]
\centering
\includegraphics[width=0.88\linewidth]{figures/fig_153.jpg}
\caption{Exhibit 5 - Exhibit 5: Global leading power semi companies' capex declined for 2 years}
\end{figure}
\begin{figure}[htbp]
\centering
\includegraphics[width=0.88\linewidth]{figures/fig_154.jpg}
\caption{Exhibit 6 - Exhibit 6: Capacity additions likely to slow in 2026-28 (kwpm - 12" equiv)}
\end{figure}
\begin{figure}[htbp]
\centering
\includegraphics[width=0.88\linewidth]{figures/fig_155.jpg}
\caption{Exhibit 7 - Exhibit 7: MOSFET: Prices up 1\% Y/Y in Apr 2026}
\end{figure}
\begin{figure}[htbp]
\centering
\includegraphics[width=0.88\linewidth]{figures/fig_156.jpg}
\caption{Exhibit 8 - Exhibit 8: MOSFET: Shipments up 23\% Y/Y in Apr 2026}
\end{figure}
\begin{figure}[htbp]
\centering
\includegraphics[width=0.88\linewidth]{figures/fig_157.jpg}
\caption{Exhibit 9 - Exhibit 9: IGBT: Prices up 1\% Y/Y in Apr 2026}
\end{figure}
\begin{figure}[htbp]
\centering
\includegraphics[width=0.88\linewidth]{figures/fig_158.jpg}
\caption{Exhibit 10 - Exhibit 10: IGBT: Shipments increased 17\% Y/Y in Apr 2026}
\end{figure}
\begin{figure}[htbp]
\centering
\includegraphics[width=0.88\linewidth]{figures/fig_166.jpg}
\caption{Exhibit 4 - Exhibit 2: Both sales and profit growth rates slow slightly Accenture sales growth (LC basis) and operating profits/adjusted operating profits}
\end{figure}
\begin{figure}[htbp]
\centering
\includegraphics[width=0.88\linewidth]{figures/fig_167.jpg}
\caption{Exhibit 1 - Exhibit 2: Both sales and profit growth rates slow slightly Accenture sales growth (LC basis) and operating profits/adjusted operating profits}
\end{figure}
\begin{figure}[htbp]
\centering
\includegraphics[width=0.88\linewidth]{figures/fig_168.jpg}
\caption{Exhibit 3 - Exhibit 3: Headcount rises slightly qoq Accenture headcount and attrition rate}
\end{figure}
\begin{figure}[htbp]
\centering
\includegraphics[width=0.88\linewidth]{figures/fig_169.jpg}
\caption{Exhibit 4 - Exhibit 4: 3Q order growth turns negative at \$-3\textbackslash{}\%\$ yoy Accenture orders (USD basis)}
\end{figure}
{\small\color{refgray}\textbf{References:} [R013] JPM：半导体设备市场的下一个结构性拐点不在需求，而在供给侧的产能释放节奏; [R017] 摩根斯坦利：功率半导体的涨价周期并非需求驱动，而是供给侧的理性收敛; [R025] GS：全球IT服务订单的转折点不是AI，而是地缘政治}

\section{外汇与资本流动：人民币定价机制微调与跨境框架重构}
\textbf{人民币中间价模型内部结构变化比方向更重要，央行在主动压缩贬值空间而非放任贬值；中国资本流动监管并非单向收紧，而是构建‘前门敞开、后门收紧’的治理框架，对金融机构意味着截然不同的机遇与风险。}

\begin{itemize}
  \item NOM模型显示官方定盘价系统性地高于模型预测值约600个基点，央行在中间价设定中持续加入比模型预期更强的‘抗贬值’力量，但操作方向与市场预期相反——它在抑制升值而非阻止贬值。
  \item NOM模型最新预测USD/CNY中间价为6.7745，较前一交易日下降385个基点，考虑逆周期因子后预测值为6.7965，欧元贡献最大正向变动（35.5 pips），隔夜美元走弱是主要驱动力。
  \item JPM系统梳理过去12个月资本流动监管更新，放松方面包括跨国公司本外币一体化资金池升级为全国性框架、境内企业对外放款额度提升且改为备案制、QDII额度审批重启；收紧方面包括个人境外投资纳入监管、跨境证券交易执法加强。
  \item 中行最受益于外汇存款利差走阔，外汇存款占集团负债16\%（同业平均4\%），离岸贷款限制取消后成为NIM正面推动力；渣打和HSBC受益于跨国公司现金管理放松和人民币国际化加速。
  \item MS认为陆家嘴论坛最重要的信号不是资本账户开放速度，而是货币政策传导机制升级，央行将隔夜逆回购利率走廊收窄至+/-25个基点，旨在改善政策利率向市场利率的传导效率，而非暗示降息。
  \item 亚洲出口领先指标NELI在7月攀升至121，为2006年以来最高水平，领先实际出口数据约三个月，且上升正在变得更加广泛，科技周期仍是主力但非科技因素（如霍尔木兹海峡重新开放）正在接力。
\end{itemize}
\begin{figure}[htbp]
\centering
\includegraphics[width=0.88\linewidth]{figures/fig_033.jpg}
\caption{Figure 1 - Figure 1: Resident outflows from China have picked up in recent years, with unrecorded flows accounting for a major share}
\end{figure}
\begin{figure}[htbp]
\centering
\includegraphics[width=0.88\linewidth]{figures/fig_034.jpg}
\caption{Figure 2 - Figure 2: The synchronized strength in CNY FX and Chinese equities has created return synergies for domestic investors}
\end{figure}
\begin{figure}[htbp]
\centering
\includegraphics[width=0.88\linewidth]{figures/fig_035.jpg}
\caption{Figure 3 - Figure 3: Resident outflows through southbound Bond Connect have picked up since last year's policy relaxation}
\end{figure}
\begin{figure}[htbp]
\centering
\includegraphics[width=0.88\linewidth]{figures/fig_036.jpg}
\caption{Figure 4 - Figure 4: The PBOC has resumed the expansion of the QDII quota}
\end{figure}
\begin{figure}[htbp]
\centering
\includegraphics[width=0.88\linewidth]{figures/fig_037.jpg}
\caption{Figure 5 - Figure 5: Dim Sum and Panda bond issuance is growing}
\end{figure}
\begin{figure}[htbp]
\centering
\includegraphics[width=0.88\linewidth]{figures/fig_038.jpg}
\caption{Figure 6 - Figure 6: FX deposits as \% of total liabilities by bank}
\end{figure}
{\small\color{refgray}\textbf{References:} [R005] JPM：市场误读了资本流动监管，真正的信号是框架重构而非收紧; [R019] 摩根斯坦利：2026陆家嘴论坛的真正信号不是开放，而是金融业增长逻辑的重新定价; [R020] NOM：人民币中间价模型正在发出一个被市场忽视的定价信号; [R021] NOM：亚洲出口的超级周期并非只有AI，一个被低估的转折点正在浮现; [R032] 摩根斯坦利：市场误读了陆家嘴论坛的信号，真正重要的是利率传导机制的重塑}

\section{风险偏好与资金流向：极端集中与泡沫特征重现}
\textbf{全球资金流入创历史纪录，美国股票单周流入1192亿美元，科技股集中度触及历史泡沫区间，BofA牛熊指标进入极端看多区域，市场定价核心从经济增长切换为政治周期与收益率曲线的博弈，谁先离场成为关键问题。}

\begin{itemize}
  \item 过去一周全球股票基金净流入1264亿美元创历史纪录，美国股票单周流入1192亿美元同样刷新历史，科技股单周流入192亿美元也创纪录，按年化推算2026年全年美国股票基金流入将达7390亿美元。
  \item AI相关十大股票占标普500指数权重已达39\%，接近甚至超过历史上几次著名泡沫的集中度峰值：19世纪80年代铁路股、1972年‘漂亮50’、1999年TMT泡沫，每一次集中度达到这个水平都对应泡沫尾声。
  \item BofA牛熊指标升至9.2进入‘极端看多’区间，自2002年以来触发17次，全球股市平均回撤2-3\%，最大回撤15-20\%，命中率约60\%，当前市场定价的核心变量已从经济增长切换为政治周期与收益率曲线的博弈。
  \item 欧洲股票连续10周流出，中国股票连续12周流出，全球资金在‘押注美国+科技’上极度集中，市场正在定价的不是繁荣而是繁荣之后的‘谁先离场’。
  \item 美国大选成为真正变量，若共和党在11月丢掉参议院，市场预期美元跌、利率跌、股票跌，华尔街对特朗普满意度处于历史高位但选民通胀满意度仅28\%，政治周期可能率先打破泡沫。
  \item 瑞士手表出口5月同比微增0.4\%，调整日历效应后增长10.4\%，但美国市场基数效应是主要推手而非需求强劲，LTM对美出口仍跌19\%，中东地缘风险成为奢侈品消费新隐忧。
\end{itemize}
\begin{figure}[htbp]
\centering
\includegraphics[width=0.88\linewidth]{figures/fig_074.jpg}
\caption{EXHIBIT 1 - EXHIBIT 1: Exports of Swiss Wristwatches watches by price categories Swiss watch exports by price - Total market (1'000 units)}
\end{figure}
\begin{figure}[htbp]
\centering
\includegraphics[width=0.88\linewidth]{figures/fig_075.jpg}
\caption{EXHIBIT 2 - EXHIBIT 2: Swiss watch exports by material Swiss watch exports by material - Total market (CHF Mil)}
\end{figure}
\begin{figure}[htbp]
\centering
\includegraphics[width=0.88\linewidth]{figures/fig_076.jpg}
\caption{EXHIBIT 3 - EXHIBIT 3: Swiss watch exports by country in absolute numbers Swiss watch exports by country - Key markets only (CHF Mil)}
\end{figure}
{\small\color{refgray}\textbf{References:} [R010] Bernstein：瑞士手表出口反弹只是表象，真正的复苏信号尚未到来; [R031] 美国银行：市场正在定价的不是繁荣，而是繁荣之后的“谁先离场”}

\section{结语}
综合来看，全球市场正经历多重定价锚点的切换：美联储沟通范式变革、AI从训练到推理的架构迁移、大宗商品从需求到供给的约束切换、以及中国从总量到结构的增长逻辑重塑。这些变化共同指向一个核心结论——投资者需要从宏观叙事下沉到微观结构，从板块整体重估转向公司层面的全球竞争力验证，在分化中寻找不对称的风险收益机会。
\newpage
\section{报告覆盖清单}
\begin{itemize}
  \item [R001] 宏观与利率：美联储沟通范式切换与全球央行分化 - GS：外汇市场的真正拐点不是地缘冲突缓和，而是美联储的鹰派姿态切换
  \item [R002] 中国地产与消费：量稳价跌与盈利预期重置 - GS：中国经济的真实温度，藏在每周高频数据的边际变化里
  \item [R003] 宏观与利率：美联储沟通范式切换与全球央行分化 - NOM：央行“叫得凶”和“咬得狠”正在分岔，市场低估了这种分化
  \item [R004] 中国医药与机器人：全球定价权再校准与‘大脑’路线未定 - GS：中国医药资产被低估的不是基本面，而是全球定价权的再校准
  \item [R005] 外汇与资本流动：人民币定价机制微调与跨境框架重构 - JPM：市场误读了资本流动监管，真正的信号是框架重构而非收紧
  \item [R006] 宏观与利率：美联储沟通范式切换与全球央行分化 - NOM：鲍威尔之后，美联储正经历一场更深刻的权力重构
  \item [R007] 中国地产与消费：量稳价跌与盈利预期重置 - JPM：消费股真正的底部不在宏观数据里，而在企业盈利预期的重新校准
  \item [R008] 区域市场：日本盈利驱动与韩国K型复苏闭合 - JPM：日本股市上涨的核心驱动力不是流动性，而是盈利能力的结构性跃迁
  \item [R009] 宏观与利率：美联储沟通范式切换与全球央行分化 - JPM：市场低估的不只是加息幅度，更是货币政策沟通机制的历史性转折
  \item [R010] 风险偏好与资金流向：极端集中与泡沫特征重现 - Bernstein：瑞士手表出口反弹只是表象，真正的复苏信号尚未到来
  \item [R011] 能源与大宗商品：供给约束下的时间错配与地缘风险折价 - GS：地震勘探行业正在进入一个被低估的、有纪律的多年度上行周期
  \item [R012] 能源与大宗商品：供给约束下的时间错配与地缘风险折价 - JPM：市场低估了供给侧的再定价，而非需求端的边际变化
  \item [R013] AI与算力：从训练到推理的架构迁移与定制芯片崛起 / 半导体与功率器件：供给驱动的涨价周期与定制化浪潮 - JPM：半导体设备市场的下一个结构性拐点不在需求，而在供给侧的产能释放节奏
  \item [R014] AI与算力：从训练到推理的架构迁移与定制芯片崛起 - JPM：AI定制芯片正在从“备选方案”变成“主导逻辑”
  \item [R015] 中国地产与消费：量稳价跌与盈利预期重置 - 摩根斯坦利：市场低估的不是下行风险，而是政策“微调”而非“转向”的定力
  \item [R016] 未被主题摘要引用 - 摩根斯坦利：中国电动车市场的份额洗牌远未结束，真正的分化才刚刚开始
  \item [R017] 半导体与功率器件：供给驱动的涨价周期与定制化浪潮 - 摩根斯坦利：功率半导体的涨价周期并非需求驱动，而是供给侧的理性收敛
  \item [R018] 中国地产与消费：量稳价跌与盈利预期重置 - JPM：外卖补贴新规的真正价值不在上限，而在让烧钱竞争变得“无法隐身”
  \item [R019] 外汇与资本流动：人民币定价机制微调与跨境框架重构 - 摩根斯坦利：2026陆家嘴论坛的真正信号不是开放，而是金融业增长逻辑的重新定价
  \item [R020] 外汇与资本流动：人民币定价机制微调与跨境框架重构 - NOM：人民币中间价模型正在发出一个被市场忽视的定价信号
  \item [R021] 外汇与资本流动：人民币定价机制微调与跨境框架重构 - NOM：亚洲出口的超级周期并非只有AI，一个被低估的转折点正在浮现
  \item [R022] 中国地产与消费：量稳价跌与盈利预期重置 - JPM：中国消费板块的真正底部不在估值，而在盈利预期的重置
  \item [R023] 能源与大宗商品：供给约束下的时间错配与地缘风险折价 - 摩根斯坦利：中国钢铁市场的真实复苏信号，被周度数据的噪音掩盖了
  \item [R024] AI与算力：从训练到推理的架构迁移与定制芯片崛起 - BARC：光学市场真正的拐点不在于AI需求本身，而在于供给结构正在经历不可逆的重塑
  \item [R025] 半导体与功率器件：供给驱动的涨价周期与定制化浪潮 - GS：全球IT服务订单的转折点不是AI，而是地缘政治
  \item [R026] 中国地产与消费：量稳价跌与盈利预期重置 - JPM：能量饮料市场真正进入“买量竞争”阶段，品牌分化才刚刚开始
  \item [R027] 区域市场：日本盈利驱动与韩国K型复苏闭合 - 摩根斯坦利：东盟电信股真正的回报来源不是5G，而是竞争格局的结构性改善
  \item [R028] 中国医药与机器人：全球定价权再校准与‘大脑’路线未定 - NOM：四足机器人已越过盈亏线，但真正考验在于“大脑”尚未定型
  \item [R029] 中国医药与机器人：全球定价权再校准与‘大脑’路线未定 - NOM：宇树科技IPO揭示的真正悬念，不在硬件成本，而在“大脑”路线的未定
  \item [R030] 中国医药与机器人：全球定价权再校准与‘大脑’路线未定 - NOM：机器人视觉的竞争壁垒不在“手”，而在“眼”和“脑”的数据飞轮
  \item [R031] 宏观与利率：美联储沟通范式切换与全球央行分化 / 风险偏好与资金流向：极端集中与泡沫特征重现 - 美国银行：市场正在定价的不是繁荣，而是繁荣之后的“谁先离场”
  \item [R032] 外汇与资本流动：人民币定价机制微调与跨境框架重构 - 摩根斯坦利：市场误读了陆家嘴论坛的信号，真正重要的是利率传导机制的重塑
  \item [R033] 能源与大宗商品：供给约束下的时间错配与地缘风险折价 - GS：铝市场真正被低估的不是地缘溢价，而是两股供给冲击的时间错配
  \item [R034] 区域市场：日本盈利驱动与韩国K型复苏闭合 - 摩根斯坦利：韩国K型复苏的裂口正在闭合，但市场低估了财富效应与政策节奏的共振
  \item [R035] 区域市场：日本盈利驱动与韩国K型复苏闭合 - 摩根斯坦利：日本制药板块的错位定价正在为深度价值投资者创造窗口
\end{itemize}
\section{逐篇报告摘录}
\subsection{[R001] GS：外汇市场的真正拐点不是地缘冲突缓和，而是美联储的鹰派姿态切换}
{\small\color{refgray}归类：宏观与利率：美联储沟通范式切换与全球央行分化}

GS：外汇市场的真正拐点不是地缘冲突缓和，而是美联储的鹰派姿态切换

过去一周，外汇市场经历了一场典型的“假突破”。地缘冲突缓和的消息一度推动新兴市场货币大幅反弹，投资者似乎看到了风险偏好回归的曙光。然而，随后美联储的鹰派表态几乎完全抹去了这些涨幅。市场参与者很容易将这理解为一次普通的“利好出尽”，但GS最新的外汇策略报告提出了一个更值得深思的判断：市场正在经历的不是情绪波动，而是定价锚点的根本性切换。

这份报告的核心洞察在于，美联储的沟通方式发生了超出预期的变化。过去几年，市场习惯了美联储相对透明、可预测的政策指引框架。但这次会议传递出的信号，无论是鹰派程度还是沟通方式，都让市场措手不及。更重要的是，这种变化并非孤立事件，而是与AI驱动的美国需求韧性、全球央行沟通范式演变等结构性力量交织在一起。这意味着，外汇市场的定价逻辑正在被重塑，而许多投资者可能仍在使用旧地图。

以下是我们从这份GS报告中提炼出的五个关键判断，以及它们对全球资产配置的含义。

1. 美联储的鹰派意外比地缘风险缓和更能决定美元走向，因为市场从未真正为战争定价

报告首先回应了一个最直观的问题：为什么地缘冲突缓和这种典型的“美元利空”，会被美联储的鹰派表态轻易压制？GS的分析给出了一个反直觉的答案——市场从来就没有为地缘风险进行过显著定价。

报告指出，市场一直存在“冲突终将解决”的基线预期，因此从未在价格中计入大量的风险溢价。相比之下，美联储鹰派立场的意外程度要大得多。从历史经验看，利率差异对美元的影响比油价更大、更持续。因此，即使地缘风险缓和暂时削弱了美元的避险吸引力，美联储更鹰派的立场也足以提供一个更强大的美元支撑因素。

这并不意味着油价对美元不再重要。报告特别指出，油价下跌恰恰是美联储开始重新评估政策立场的最初催化剂。但问题在于，自年初以来，AI驱动的美国需求强劲增长，已经成为推高美国增长、通胀和中性利率定价的主要因素。因此，美联储的鹰派立场虽然看似由能源价格触发，但其背后是更坚实的经济基本面支撑。这种“新”的美元支撑因素，实际上与过去一段时间持续支撑美元需求的旧因素一脉相承。

对于投资者而言，这意味着短期内美元强势的逻辑可能比表面看起来更坚固。单纯押注地缘风险缓和带来的美元走弱，可能低估了利率差异这一更核心的定价力量。

2. 新兴市场套利交易并未终结，但融资货币的选择决定了收益的天壤之别

报告对新兴市场套利交易的观点，是全文中最具操作性的部分。GS并未因美联储鹰派而全面看空新兴市场货币，而是提出了一个更精细的判断：套利交易仍能产生正回报，但前提是选对融资货币。

报告区分了两种市场环境：在“股市涨、利率跌”的背景下，用美元融资的套利表现更好；但在“股市涨、利率涨”的背景下（即当前的环境），用日元或欧元等低收益货币融资，反而能带来更高的正回报。这一判断基于历史数据，尤其是 年加息周期中，高收益货币与低收益货币之间的显著分化。

具体到新兴市场内部，报告建议投资者关注几个结构性差异。墨西哥比索因其对美元周期的敏感性，在美元相对强势的环境中表现突出，但即将到来的USMCA谈判是主要风险。在拉美市场，报告建议将部分巴西雷亚尔头寸转向哥伦比亚比索，前提是哥伦比亚国内政治局势保持稳定。而巴西国内政治噪音上升，加上央行可能转向鸽派，将削弱雷亚尔的套利吸引力。

这意味着，当前新兴市场套利交易的核心挑战不是方向问题，而是结构问题。投资者需要放弃“一刀切”的新兴市场配置思路，转向更精细的货币对选择和对冲策略。

日元是这份报告中逻辑最清晰的货币之一。报告指出，即使日本央行在6月加息并继续释放加息信号，其加息步伐仍不足以抵消日元的基本面弱势。美联储加息概率上升，进一步推高了美元兑日元汇率。

报告中的GSBEER模型显示，日元在美联储鹰派冲

[... middle omitted ...]

条件：美国衰退风险上升或日本央行转向更鹰派。目前来看，这两个条件都不具备。

对于持有日元敞口的投资者，这份报告提供了一个清晰的观察框架：不要期待日本央行能靠渐进式加息扭转日元趋势，也不要高估干预的效果。日元的弱势可能会持续，直到美国经济出现实质性放缓信号。

4. 瑞士法郎的“软干预”只是象征性表态，利差劣势才是真正的定价力量

瑞士央行在本周维持利率不变，符合市场预期。但报告关注的重点是瑞士央行在声明中对汇率的措辞调整。瑞士央行将此前“有更强的意愿干预以应对瑞郎升值”的表述，微调为“如有必要，有更强的意愿干预”。这一变化被市场解读为干预立场的软化。

然而，GS认为这种措辞调整基本上是象征性的。报告指出，在众多主要央行正在加息或考虑加息的背景下，瑞士央行的按兵不动使其处于全球央行光谱中偏鸽派的一端。这种鸽派分化本身就应对瑞郎构成压力。瑞士央行行长也承认，与主要央行的利差扩大正在对瑞郎形成下行压力。GS因此判断，措辞调整不足以改变这一动态，瑞郎短期内将继续表现不佳。

\subsection{[R002] GS：中国经济的真实温度，藏在每周高频数据的边际变化里}
{\small\color{refgray}归类：中国地产与消费：量稳价跌与盈利预期重置}

当市场还在为月度PMI的读数争论不休时，一份来自GS的每周经济追踪报告揭示了一个更微妙的现实：中国经济正在经历的不是简单的“复苏”或“放缓”，而是一个由多重力量拉扯、板块剧烈分化的“再均衡”过程。这份6月18日发布的报告，通过四大类高频指标——消费与出行、生产与投资、其他宏观活动、市场与政策——拼凑出一幅比任何单一数据都更立体的图景。

核心判断是：经济活动的“量”在边际上企稳，但“价”的信号依然疲弱，且结构性分化正在加剧。这意味着，对于产业决策者和投资者而言，整体宏观数字的参考价值正在下降，真正重要的是捕捉不同板块内部边际变化的拐点。这份报告的价值不在于给出一个“好”或“坏”的结论，而在于提供了每周更新、可交叉验证的观察框架，让读者有能力在官方数据发布之前，提前感知水温的变化。

1. 房地产市场出现了一个值得警惕的背离：交易量回暖未能阻止价格继续探底

报告中最引人注目的信号来自房地产市场。30城新房成交量在最近一周小幅上升，并且高于去年同期水平；16城二手房成交量也基本持平于高位。从“量”的角度看，政策松绑后的脉冲效应似乎仍在持续，市场并没有出现断崖式下跌。

然而，量稳的背后是价的持续疲软。国家统计局70城二手房价格指数在5月继续下降。更令人担忧的是，报告追踪的多个独立数据源——中原、贝壳、诸葛、国信达——所构建的房价指数均显示，当前价格水平已显著低于2025年，甚至逼近或低于2019年基准。这意味着，以价换量仍然是市场的主旋律。卖方为了促成交易，正在不断下调预期价格。

这个背离的含义是深远的。交易量的企稳，更多反映的是积压需求的释放和业主的主动降价，而非市场信心的根本性修复。只要价格下跌预期没有扭转，潜在购房者就会继续持币观望，等待更低的“底部”。房地产市场的真正企稳，需要看到“量价齐稳”甚至“量稳价升”的信号，而目前我们只看到了前半段。对于持有资产或从事相关行业的人来说，量稳价跌的阶段可能是最消耗耐心的——既看不到崩盘的恐慌，也看不到反转的希望。

GS这份报告从5月起改为周度发布，一个明确的原因是为了“更紧密地追踪能源价格上涨对经济活动的供给冲击”。报告中的数据显示，6月18日，国内汽油和柴油价格分别下调了515元/吨和495元/吨，这是在前期国际油价（Brent）大幅攀升后的首次调降。但更值得关注的是，硫酸、甲醇、聚丙烯等基础化工品的价格在过去一周继续上行，其中硫酸价格已较年初翻倍有余。

这意味着，上游能源成本的上涨，正在通过产业链条逐步向下游传导。但传导并非同步且均匀的。化肥、聚丙烯等与农业、包装、汽车相关的化工品价格涨幅温和，而硫酸这种工业基础原料的价格则出现了剧烈波动。这种非线性的传导，会给不同行业带来截然不同的成本压力。

对于制造业企业而言，当前最需要警惕的不是某一周的价格波动，而是成本压力的累积效应。如果能源价格维持高位，而下游终端需求无法同步提价，那么中游制造环节的利润将面临持续的挤压。这份报告提供的数据，恰好可以用来判断这种挤压是正在加剧还是边际缓和。例如，通过对比化工品价格指数与钢价、煤耗等生产指标，可以初步评估工业企业的“成本-产出”剪刀差是否在扩大。

3. 消费场景恢复的“天花板”效应已现，结构升级而非总量扩张成为新主线

报告中的出行数据提供了关于消费的另一个视角。国内客运航班量在过去一周有所上升，但依然低于去年同期水平。同时，航班取消率虽然下降，但仍处于高位。交通拥堵指数则小幅上升。将这些信号综合起来，可以看到一个矛盾的现象：人们在出行，但出行的效率和意愿似乎没有完全恢复。

更值得玩味的是Morning Consult消费者信心指数。该指数在经历了2025年初的低谷后持续回升，到2026年7月已高于2024年同期水平。然而，这一信心改善并未完全转化为更强的消费支

[... middle omitted ...]

相对稳健但缺乏惊喜的图景。钢铁需求和产量在最近一周均小幅回升，但绝对量仍低于疫情前的2019年水平。沿海省份的日均煤炭消费量则保持在高于去年同期的水平，显示电力需求依然旺盛。

最引人注目的是地方政府专项债券的发行进度。截至报告发布时，2026年累计发行规模已达1.67万亿元人民币，节奏明显快于2025年同期。这直接支撑了基建投资的资金来源。同时，官方港口集装箱吞吐量和20个主要港口离港船舶货运量均保持在高位，且高于去年同期水平，表明出口活动依然强劲。

这几组数据合在一起，形成了一个重要的判断：在房地产投资持续疲软的背景下，基建和出口正在充当经济稳定器的角色。但这里有两个关键问题报告没有完全回答。第一，专项债的加速发行能否持续？如果后续发行节奏放缓，基建投资的支撑力将减弱。第二，出口的韧性更多是价格因素（出口商品涨价）还是数量因素（全球需求仍旺）？如果是前者，那么对国内就业和制造业利润的拉动效应可能有限。这两个问题的答案，将决定当前的经济韧性是可持续的，还是暂时的。

\subsection{[R003] NOM：央行“叫得凶”和“咬得狠”正在分岔，市场低估了这种分化}
{\small\color{refgray}归类：宏观与利率：美联储沟通范式切换与全球央行分化}

NOM：央行“叫得凶”和“咬得狠”正在分岔，市场低估了这种分化

理解全球央行当前的政策节奏，不能只看加息还是降息。NOM最新一份研报提供了一个更精细的观察框架：用彭博的央行沟通情绪指数，去衡量央行“说了什么”，再与其实际行动对比，看“做了什么”。结论是清晰的——三大央行的言行一致性正在显著分化，而这种分化背后是各自经济基本面和政策约束的根本差异。对于资产定价而言，这意味着一刀切的“全球央行转向”叙事可能已经失效，投资者需要重新校准对不同市场的利率风险定价。

这份报告最值得关注的判断不是某个具体的利率预测，而是它揭示了一个正在发生的结构性变化：央行沟通本身正在成为独立的政策工具，且其可信度和有效性正在因央行的“兑现记录”而出现分化。欧洲央行正在用行动强化其鹰派沟通的可信度，英国央行则在用“按兵不动”削弱其鹰派信号的威慑力，而美联储的新任主席则可能让这类沟通指数本身失去意义。这三条线索，每一条都指向不同的资产配置含义。

1. 欧洲央行正在用行动证明“鹰派不是空话”，这是当前最值得追踪的定价信号

NOM编制的ECBspeak指数已经升至20.7，这是2022年3月以来的最高水平。但更关键的是，与上次加息周期不同，这次欧洲央行在鹰派沟通升温后迅速跟进——在指数达到这一水平时，欧洲央行已经完成了25个基点的加息。而在2022年3月，指数达到类似水平后，欧洲央行等了四个月才启动首次加息。

这种“言行一致”的转变，直接推动了NOM欧洲经济团队上调加息预测：在9月和12月各加25个基点，并延续到2027年3月。这并非简单的预测调整，而是对央行行为模式变化的确认——欧洲央行正在从“落后于曲线”转向“主动管理预期”。

这一变化的底层逻辑在于，欧元区的通胀粘性比预期更强，尤其是第二轮通胀效应的担忧仍在主导决策者的思维。NOM报告虽然没有完全展开，但这里需要指出的是，欧洲央行的这一行为模式变化，意味着欧元区利率的“上限”可能高于市场当前定价。如果市场仍在按“欧洲央行很快会转向宽松”来定价，那么这种定价偏差本身就构成了交易机会。

2. 英国央行正在用“按兵不动”证明其鹰派沟通的“空转”，政策可信度面临侵蚀

与欧洲央行形成鲜明对比的是英国央行。BOEspeak指数目前为4.4，虽然较此前有所上升，但英国央行在昨天的会议上以7比2的投票结果维持利率不变。NOM欧洲经济团队已经将今年的加息预测从一次25个基点调整为“零加息”，并预计2027年下半年将进行两次25个基点的降息。

这里的关键比较是，在上次加息周期中，BOEspeak指数达到3.7时，英国央行就启动了首次加息。而这次指数已经升至4.4，央行却选择按兵不动。NOM报告给出的解释是，英国的政策利率（3.75\%）起点远高于欧洲央行（2.25\%），且被认为已经处于中性利率之上。同时，英国劳动力市场正在趋势性走弱，通胀涨幅也低于预期。

但更深层的含义在于，英国央行的“鹰派沟通”正在失去市场约束力。如果市场开始相信“英国央行只会叫不会咬”，那么通胀预期锚定的难度就会上升。这反过来可能迫使英国央行在未来采取更激进的行动来重建信誉。对于英镑和英国国债的投资者而言，这意味着一个两难局面：如果相信英国央行的沟通，则利率应上行；如果不相信，则通胀风险溢价需要重新定价。

3. 美联储的“沟通工具”可能正在失效，新主席的风格让指数本身失去意义

Fedspeak指数目前为7.7，从历史标准看，这是一个非常温和的上升。在上个周期，Fedspeak指数在2021年5月就达到了7.7，但美联储直到10个月后才开始加息，事实证明这一反应严重落后于曲线。NOM报告暗示，当前类似的风险正在积聚：通胀上升、金融条件宽松，叠加财政主导和美国政治周期，可能导致历史重演。

但报告提出的一个更具颠覆性的观

[... middle omitted ...]

率的下降。在这种环境下，利率波动率可能系统性上升，而依赖于“央行看跌期权”的资产定价逻辑需要重新审视。

4. 报告尚未完全回答的关键问题：能源价格下跌能否成为改变欧洲央行行为的变量？

NOM报告承认，如果近期能源价格的下跌能够持续，欧洲央行可能会降低对通胀的担忧，从而软化其鹰派立场。但报告也明确指出，“目前几乎没有这种迹象”。

这里存在一个尚未被充分讨论的二阶效应：能源价格下跌本身可能改变欧洲央行的通胀预期锚定机制，但前提是这种下跌被视为结构性而非周期性。如果市场认为能源价格下跌是需求疲软的结果（而非供给改善），那么欧洲央行反而可能将其视为经济走弱的信号，从而在鹰派沟通与实际政策之间产生新的张力。

换句话说，能源价格下跌对欧洲央行的影响方向并非单向的“宽松信号”。它是通过需求渠道还是供给渠道传导，决定了最终的政策含义。NOM报告没有展开这一讨论，而这恰恰是当前市场定价中最容易产生分歧的环节。

5. 一个可供读者自行验证的观察框架：用“言行差距”追踪央行可信度

\subsection{[R004] GS：中国医药资产被低估的不是基本面，而是全球定价权的再校准}
{\small\color{refgray}归类：中国医药与机器人：全球定价权再校准与‘大脑’路线未定}

GS：中国医药资产被低估的不是基本面，而是全球定价权的再校准

当一家全球顶级投行在参加了2026年6月的全球医疗健康大会、与30多家机构投资者交流后，得出了一个看似矛盾的结论——中国医药股年内平均下跌11\%，但基本面（临床数据、BD活动）依然强劲——这通常意味着市场正在犯一个经典错误：用短期情绪定价长期资产。

GS这份研报的核心判断值得每一个关注中国医药产业的决策者认真对待：当前中国医药资产的风险回报正在向有利方向倾斜，但驱动下一阶段行情的逻辑，将从“板块整体重估”转向“公司层面的全球竞争力验证”。这不是一个简单的抄底信号，而是一个结构性分化的起点。

1. 市场真正担心的不是基本面，而是两个“看不见”的政策变量

投资者对中国医药板块的谨慎态度，并非源于对临床数据或研发管线的否定。GS在与投资者的交流中发现，真正压制情绪的是两个政策层面的不确定性。

第一个是中国与美国之间关于BD交易的监管审查。虽然目前尚未出现实质性的阻断案例，但双边监管趋严的预期已经影响了投资者的风险偏好。GS的调研显示，投资者的讨论更多集中在中国监管一侧，因为政策透明度较低。一个值得关注的细节是：多数投资者认为，如果中国监管审查聚焦于平台级技术输出（而非单一资产），结构性变化的概率不大。这意味着，近期的BD交易趋势将成为判断政策影响的关键观测窗口。

第二个是国内医疗反腐对盈利的潜在冲击。这一变量在 年已经对行业产生过显著影响，但市场仍在评估其持续时间和影响深度。

这两个变量的共同特点是：它们不是对行业长期价值的否定，而是对短期可见度的干扰。GS的判断是，随着催化剂可见度的提升（即将到来的临床数据、BD公告、二季度财报），市场将重新关注基本面本身。

2. 中国创新在全球药企战略版图中的位置，比市场认知的更核心

研报中最具信息量的部分，并非来自GS自己的分析，而是来自全球大型药企在大会上的表态。这些表态揭示了一个正在发生的结构性变化：中国已经不再只是“低成本制造基地”，而是正在成为全球创新药研发的“战略伙伴”。

Gilead明确表示，其优先发展的企业项目中约有一半源自中国或与中国创新相关。辉瑞则给出了一个更具时间维度的判断：当前中国生物科技公司仍依赖欧美药企进行全球临床试验和商业化，但到2030年前后，它们可能具备自主建立这些能力的基础。AstraZeneca和Novartis也分别强调了持续合作意愿和对中国创新资产的开放态度。

这些信息叠加在一起，指向一个清晰的结论：全球大型药企对中国创新的依赖不是短期的交易性行为，而是结构性的战略布局。这种依赖的深度和广度，远未被当前股价所反映。

GS同时指出，全球大型药企对中国资产的偏好正在发生变化——它们更倾向于那些能够“干净地”融入全球开发战略和合作伙伴管线的资产。这意味着，不是所有中国创新都能享受全球化溢价，只有那些具备全球竞争力和可转化数据的资产才能获得认可。

在CDMO（合同研发生产组织）领域，GS的调研显示投资者情绪正在改善，但改善是不均匀的。这背后反映了一个核心问题：在经历了两年多的地缘政治压力后，哪些公司真正证明了自身的结构性竞争力？

药明康德虽然因被列入美国1260H清单而经历了短期股价波动，但投资者的注意力已经转向了其多元化的增长驱动——除了GLP-1相关的多肽业务外，订单动能依然强劲，二季度可能出现的指引上调被视为关键催化剂。药明生物的特许权使用费收入被视为中期盈利的支撑，尽管可见度较低。药明合联的订单动能保持强劲，但投资者正在关注新加坡工厂的商业化进展。

康龙化成则因其GLP-1相关的产能扩张而持续受到关注，但投资者讨论的焦点已经转向了扩产后的产能利用率和利润率轨迹。

这些细节揭示了一个重要逻辑：在CDMO领域，决定估值分化的关键不再是“是否拥有GLP

[... middle omitted ...]

，那些采用随机对照设计（以全球标准疗法为对照组）、患者基线特征与国际接轨、使用全球标准终点、并获得FDA/EMA早期沟通的临床试验，被认为具有更高的可转化性。这意味着一件事：对于中国生物科技公司而言，临床开发策略的重要性已经与生物学机制本身相当。

GS据此判断，下一阶段的行情将由公司层面的具体验证驱动：可信的全球数据包、成功的BD执行、以及监管可转化性的证据。这将对资产定价产生深远影响——拥有高质量全球资产、明确BD期权价值和严谨试验设计的公司将获得溢价，而拥挤的机制或可转化性较低的数据集可能面临估值分化。

尽管GS整体看好中国医药板块的风险回报，但报告中有几个关键问题并未完全展开，而这些恰恰是投资者最需要关注的。

第一个问题是：政策不确定性何时能够转化为确定性？无论是中美BD监管审查还是国内医疗反腐，其影响路径和时间表都不清晰。GS提到“催化剂可见度改善”，但并未给出明确的触发时间点。这意味着，即使基本面完好，估值修复可能仍需要等待一个“政策靴子落地”的信号。

\subsection{[R005] JPM：市场误读了资本流动监管，真正的信号是框架重构而非收紧}
{\small\color{refgray}归类：外汇与资本流动：人民币定价机制微调与跨境框架重构}

JPM：市场误读了资本流动监管，真正的信号是框架重构而非收紧

过去几个月，中国在资本流动领域的多项监管更新引发了市场广泛讨论。多数解读倾向于将其简单归类为“资本管控收紧”，尤其是6月出台的对外投资新规。然而，JPM最新发布的这份研报提出了一个截然不同的核心判断：这些政策调整并非单向收紧，而是一场旨在构建完整资本流动治理框架的结构性努力。市场真正需要关注的，不是短期流动方向，而是这个框架如何重塑人民币汇率、银行业竞争格局以及跨境金融服务的长期逻辑。

这份报告的价值不在于罗列政策条文，而在于提出了一个关键洞察：政策制定者在“前门”和“后门”之间做出了明确的区分——一方面扩大合规渠道的容量和灵活性，另一方面对非正规渠道施加压力。这种“有管理的开放”策略，对不同类型的金融机构意味着截然不同的机遇与风险。

1. 这轮政策调整的核心不是限制资本外流，而是建立一套“前门敞开、后门收紧”的治理框架

JPM在报告中系统梳理了过去12个月涉及资本流动的多项监管更新。结论是清晰的：这不是单向的收紧或放松，而是一个混合体。放松的方面包括：跨国公司本外币一体化资金池业务从试点升级为全国性框架，允许企业用一个跨境资金池统一管理境内外本外币资金；境内企业向境外放款的额度提升，且审批制改为备案制；债券通南向通的合格机构名单扩大，QDII额度审批重启；以及境内银行间接向境外企业放贷的严格标准被取消。

收紧的方面同样明确：6月对外投资新规将个人投资者的境外投资纳入监管范围；跨境证券交易监管执法力度加强；以及要求境内企业及时将境外IPO募集资金调回境内。

JPM的判断是，这些措施加在一起，传递的信号不是“关门”，而是“换锁”——让合规渠道更畅通，同时对灰色地带施加更大压力。对于投资者而言，这意味着需要重新评估跨境资本流动的长期路径，而不是被短期的政策波动所左右。

2. 人民币汇率短期获得支撑，但居民资产全球化配置的结构性需求才是中期主线

报告从宏观和汇率角度给出了一个重要判断：当前的外汇背景并不支持“急需收紧资本外流”的叙事。恰恰相反，人民币面临的是升值压力而非贬值压力。JPM估算，2023至2025年居民累计外流约1.3万亿美元，其中约5070亿美元通过非官方渠道流出。但值得注意的是，自2025年中以来，非官方渠道的累计流出规模已稳定在约5000亿美元，这意味着外流压力正在缓解。

更重要的是，以汇率调整后的回报率衡量，中国股票相对于美国市场的表现已经发生了结构性转变，从过去的持续折价转为更有利的比较。这本身就在缓解外流压力。

短期来看，对外投资规则的收紧可能通过减少外流压力和促进资金回流，为人民币提供温和支撑。但JPM也明确指出了中期风险：中国居民对全球资产的配置比例相对于区域同行仍然结构性偏低，这意味着多元化的内在需求将持续存在。政策制定者显然也意识到了这一点——他们在收紧“后门”的同时，正在通过债券通南向通、QDII额度重启等渠道打开“前门”。对于投资者而言，这意味着人民币资产的持有逻辑需要从“避险”转向“选择”——哪些合规渠道能提供最优的全球配置效率，才是关键问题。

3. 银行业最大的结构性机会不在零售财富管理，而在企业及投资银行（CIB）业务

这是报告中针对银行板块最核心的洞见。市场普遍关注的是政策收紧对香港银行来自内地访客（MCV）的财富管理业务可能造成的冲击。JPM承认这一风险，但认为市场低估了CIB业务扩张带来的上行空间。

具体而言，三项放松措施将直接利好具备跨境企业银行服务能力的机构：跨国公司现金管理规则的放松，将惠及那些同时具备境内外本外币现金管理能力的银行；境内企业境外放款的放松，将利好那些能够提供咨询和货币对冲服务的银行；人民币国际化的加速——体现在点心债和熊猫债发行量的持续增长、债券通南向通扩容等

[... middle omitted ...]

面机会：外汇存款息差的扩大。根据文件72号，境内银行向境外企业放贷的资金使用限制被取消。这将在JPM看来改善相关贷款需求，而那些拥有强大境内外汇存款基础的银行将能够利用这些存款来支持境外企业贷款，从而改善存款息差。

数据支撑了这一判断：在中国银行中，中国银行的外汇存款占集团负债的16\%，远高于同行平均的4\%。如果其他条件不变，这将为中国银行的净息差提供支撑。

但这里需要指出报告没有完全展开的一个关键问题：这种息差优势的可持续性。外汇存款息差的扩大取决于两个条件：一是境内外汇存款成本保持相对低位，二是境外贷款需求持续旺盛。前者可能受到人民币汇率预期变化的影响，后者则取决于全球利率环境和中国企业的海外投资周期。报告没有深入讨论这两个假设条件的变化风险，而这恰恰是投资者在做出判断时需要自行验证的部分。

5. 报告尚未完全回答的关键问题：合规渠道的“容量天花板”在哪里，以及非正规渠道的压降速度

JPM构建的分析框架逻辑清晰，但有两个关键变量报告没有给出明确的量化判断。

\subsection{[R006] NOM：鲍威尔之后，美联储正经历一场更深刻的权力重构}
{\small\color{refgray}归类：宏观与利率：美联储沟通范式切换与全球央行分化}

市场对6月FOMC会议的解读大多聚焦于“点阵图转鹰”与“沃什偏鸽”的表层矛盾。但NOM这份研报揭示了一个更值得关注的信号：美联储正在发生的不只是利率路径的摇摆，而是一场围绕数据定义权、通胀框架和机构治理的深层重构。真正重要的不是2026年会不会加息，而是谁来决定“什么算通胀”。

1. 点阵图的鹰派偏移被沃什的措辞对冲，但真正重要的是框架的不确定性

6月点阵图的中位数显示，2026年政策利率将升至3.75\%，比当前水平高出半个加息幅度。这一结果明显比市场预期更鹰派。然而，沃什在新闻发布会上的表态却呈现出相反的温度。他将点阵图形容为“用铅笔画的——那种带大橡皮的铅笔”，暗示这些预测随时可能被擦除。他还刻意强调“一半的同事认为利率应该维持在当前水平或更低”，尽管实际上只有一位参与者预期2026年降息。

这种“鹰派数据、鸽派叙事”的组合并非简单的沟通技巧，而是反映出沃什正在尝试重新定义美联储对外沟通的锚点。他更倾向于用定性判断来对冲量化预测的刚性，这本身就是在削弱点阵图作为政策信号的市场权重。

对投资者的含义是：点阵图的参考价值正在下降，沃什的措辞将成为更重要的政策风向标。未来市场对美联储的解读，可能需要从“数点”转向“听词”。

研报中一个容易被忽略的细节是沃什对通胀目标的表述。当被问及2\%的通胀目标时，他回答：“我倾向于关注小数点的左侧。”这意味着他可能将2.0\%至2.9\%的通胀区间都视为符合2\%目标的范畴。

这一表述的潜在影响不容低估。如果美联储正式或非正式地将通胀容忍区间上移，那么当前核心PCE在3.4\%左右的水平就不再构成加息的紧迫理由。NOM预计5月核心PCE同比将加速至3.457\%，但沃什的框架恰恰可能将这一数字“正常化”。

然而，这里有一个关键问题尚未解决：沃什的个人观点能否转化为FOMC的共识。报告指出，许多FOMC参与者近期对沃什偏好的“截尾均值PCE”指标提出了质疑，认为它在低估通胀趋势。如果沃什试图通过新成立的“通胀框架工作组”来推动这一转变，可能面临来自委员会内部的强烈阻力。

3. 五个工作组的设立，本质上是沃什在重新塑造美联储的“数据基础设施”

沃什宣布成立五个工作组，分别聚焦于：美联储沟通、数据测量、资产负债表政策、生产率与就业、以及通胀框架。这些工作组的名称看似技术性，但其战略意图非常清晰——沃什正在试图绕过FOMC内部的分歧，通过建立新的政策知识基础设施来推动变革。

最值得关注的是“数据测量”工作组。沃什在新闻发布会上公开批评联邦政府的统计数据“大多采用过时的调查方法”，并表示倾向于使用来自私人部门的“实时数据”。这一立场与前任主席鲍威尔形成对比：鲍威尔虽然也指出过调查回应率下降的问题，但多数FOMC参与者仍然将联邦统计机构的数据视为“黄金标准”。

如果沃什成功推动数据源的多元化，那么通胀的判断标准、劳动力市场的评估方式、乃至GDP的测算方法都可能发生系统性变化。这不仅是技术问题，更是权力问题——谁的数据被采纳，谁的叙事就占主导。

但报告也指出，这些工作组要到“今年秋季”才会提供初步见解，最终建议预计在年底前完成。这意味着在2026年下半年，市场将面临一个政策真空期：加息的门槛可能因工作组未完成而被人为抬高，但通胀压力却在持续积累。

4. 核心PCE加速背后的结构性问题：金融服务业价格反弹比商品通胀更难治理

NOM预计5月核心PCE环比将从4月的0.230\%加速至0.376\%，同比从3.307\%升至3.457\%。推动这一加速的主要动力并非商品价格——核心CPI商品通胀在5月甚至转为负值——而是来自PPI数据的服务业分项，特别是金融服务价格和进口航空票价的强劲上涨。

更值得警惕的是“超级核心PCE”的走势。NOM预计5月超级核心PCE环比将加速至0.5\%，创下

[... middle omitted ...]

有一段值得反复推敲的分析：沃什将货币政策的限制性描述为“不均匀的”，认为利率工具对房地产市场形成压制，而资产负债表政策却在支撑金融市场。他暗示，当前的政策组合存在内在矛盾——利率偏紧、资产负债表偏松。

这一判断引出一个核心问题：沃什是否在试图通过调整资产负债表政策（如放缓缩表或提前结束MBS减持）来为利率政策争取空间？如果他的策略是通过“工具再分配”来降低利率的紧缩压力，那么市场需要关注的不只是点阵图，还有联储的资产负债表路径。

NOM的基准预期是美联储在2027年底前维持利率不变，同时资产负债表将以每月约200亿美元的速度温和扩张，且新增购买将偏向短期国债。但这一路径高度依赖于沃什能否在FOMC内部获得足够的支持来推进他的框架改革。

如果沃什的工作组最终提出的建议——例如采用替代通胀数据、扩大资产负债表工具的适用范围——未能获得多数委员支持，那么他可能面临“有愿景、无执行”的困境。反之，如果这些建议得以落地，美联储的政策框架将发生自沃尔克时代以来最深刻的变革。

\subsection{[R007] JPM：消费股真正的底部不在宏观数据里，而在企业盈利预期的重新校准}
{\small\color{refgray}归类：中国地产与消费：量稳价跌与盈利预期重置}

JPM：消费股真正的底部不在宏观数据里，而在企业盈利预期的重新校准

2026年5月的中国零售数据，以-0.6\%的同比增速，触发了市场自2022年底以来的第一次月度负增长警报。这份来自JPM的消费策略报告，没有停留在“需求疲软”的简单叙事里。它真正想传递的信号是：市场低估的不是消费何时复苏，而是企业盈利预期需要多大幅度、多长时间的下修。当前板块的估值压缩，可能还没有完全定价二季度业绩期即将出现的密集盈利预警。

这份报告的核心洞察，不是告诉你消费有多差——这一点市场已经反映在股价里——而是帮你在“差”的结构里，找到哪些公司的商业模式能扛住压力，哪些公司会被原料成本和需求双杀。更重要的是，它提醒投资者：现在说“见底”仍为时过早，选股逻辑必须从“买赛道”切换为“买护城河”。

1. 零售数据的结构比总量更值得警惕：耐用品的“高基数陷阱”正在变成“需求真空”

5月零售总额的-0.6\%固然是2022年12月以来首次负增长，但更值得关注的是分项数据的结构性恶化。JPM的数据拆解显示，拖累大盘的三大板块——汽车（同比-16\%）、家电（-16\%）、家具（-9\%）——并非单纯因为高基数，而是反映出居民对高单价、高决策成本商品的消费意愿正在系统性收缩。

汽车和家电在 年经历了政策刺激驱动的集中释放，高基数确实是因素之一。但问题在于，当补贴效应消退后，后续需求并没有自然衔接。家电同比-16\%与家具-9\%的组合，暗示房地产后周期消费的传导机制已经断裂——即使新房交付量边际改善，消费者也没有像过去那样同步配置家居产品。这意味着，耐用消费品的“以旧换新”政策红利正在快速衰减，而市场此前对2026年下半年销售回弹的预期，可能需要大幅下调。

对投资者而言，这一结构变化意味着：任何依赖“政策刺激-需求爆发-库存回补”逻辑的消费股，都可能面临连续两个季度以上的盈利下修风险。JPM明确提示，二季度业绩期将出现更多公司下调指引，这比宏观数据本身更具杀伤力。

2. 线上与线下的分化不是新故事，但“线下服务消费”的韧性正在被侵蚀

报告中的一个关键数据点是：5月线上实物零售同比增长2.6\%，而线下实物零售同比下降2.2\%。线上渗透率的持续提升已经是旧叙事，真正的新信号来自线下服务消费——餐饮收入同比仅增长0.6\%，较4月的2.2\%大幅收窄。

餐饮向来被视为“口红效应”的受益者，即消费者在缩减大件开支时，会转向小额、高频的服务消费来维持生活品质。但5月的数据表明，这种替代效应也在减弱。结合失业率维持在5.1\%、16-24岁青年失业率仍处高位的背景，服务消费的下行可能才刚刚开始。如果6-7月餐饮增速进一步转负，那么整个消费板块的盈利底盘——日常消费品的量价稳定性——也将面临挑战。

这意味着，此前被认为“防御性最强”的食品饮料、日化板块，可能无法独善其身。JPM在报告中把“软饮料”列为5月表现最好的品类（+6\%），但这更多是季节性因素和个位数价格战的产物，而非需求扩张的信号。

3. 估值压缩尚未结束，但“盈利稳定溢价”正在成为新的定价锚

报告显示，过去一个月中国必需消费板块的前瞻PE从16.8倍压缩至15.2倍（-9.6\%），可选消费从12.9倍压缩至12.1倍（-6.5\%），幅度均超过MSCI中国指数（-3\%）和恒生指数（-3.7\%）。表面上看，消费板块的估值正在“回归合理”，但JPM的隐含意思是：当前的估值水平还没有充分反映盈利下修的风险。

一个关键参考是：必需消费目前的15.2倍PE，仍高于其五年均值（约14倍）。如果二季度业绩期出现广泛的下调指引，估值中枢可能进一步下沉至13-13.5倍。可选消费的12.1倍PE看似便宜，但考虑到汽车、家电等权重股的盈利风险，这个估值并不安全。

JPM给出的选股框架，实际上是在构建一个“盈利稳定溢

[... middle omitted ...]

同比上涨3.9\%，非食品CPI仅上涨1.9\%，两者之间的剪刀差正在扩大。这意味着下游消费企业正在承受成本压力，但提价能力有限——要么是因为竞争格局不允许，要么是因为消费者价格敏感度太高。如果PPI在三季度维持高位，而终端需求继续疲软，那么消费企业的毛利率压缩将比市场预期的更严重。

目前市场对消费股的盈利预期，可能还没有充分计入“成本冲击”这个变量。对于饮料、日化、包装食品等毛利率敏感型行业，这是一个潜在的负面催化剂。JPM选择用“避免”来应对，但没有量化风险敞口——这恰恰是完整报告中最值得细读的部分。

5. 对读者的观察框架：如何判断消费股是否到了“可以买入”的区间

基于JPM的报告逻辑，可以提炼出三个判断消费股进场时点的观察维度，供读者自行验证：

第一，盈利预警是否集中出清。二季度业绩期将是第一个压力测试点。如果大量公司下调全年指引，且股价在消息发布后不再创新低，那可能意味着最坏的情况已经被定价。反之，如果公司选择“硬撑”不调指引，反而会拉长估值回归的时间。

\subsection{[R008] JPM：日本股市上涨的核心驱动力不是流动性，而是盈利能力的结构性跃迁}
{\small\color{refgray}归类：区域市场：日本盈利驱动与韩国K型复苏闭合}

JPM：日本股市上涨的核心驱动力不是流动性，而是盈利能力的结构性跃迁

当全球投资者还在争论日本央行何时加息、日元是否会突破160关口时，日本股市已经在2026年上半年交出了一份令人无法忽视的成绩单。截至6月17日，东证指数（TOPIX）年内上涨18\%，跑赢MSCI全球指数、标普500和欧洲STOXX 600。日经225更是录得39\%的涨幅，显著超越TOPIX的表现。

这份来自JPM的最新日本股权策略报告，将2026年底的日经225目标价上调至75,000点，TOPIX目标价上调至4,400点。但真正值得关注的，不是目标价本身，而是支撑这一判断的核心逻辑——市场正在经历从“估值修复”到“盈利驱动”的范式切换。

过去几年，日本股市上涨的主要叙事是“安倍经济学”带来的公司治理改革、股票回购和外资回流。这些因素仍然存在，但边际效应正在递减。JPM认为，当前阶段真正驱动市场的变量，是AI半导体超级周期叠加日本企业盈利能力结构性改善所带来的盈利增长。TOPIX每股收益预计在2026财年同比增长11\%，2027财年再增长12\%，这一增速已经显著超出市场一致预期。

更值得深思的是，报告指出日经225的AI相关股票市值占比已接近50\%，TOPIX约为30\%。这意味着，日本股市已经不再是“价值股洼地”，而是正在成为全球AI基础设施投资的重要映射。那些仍然以“低估、高股息、低增长”框架看待日本市场的投资者，可能正在错过这一轮结构性重估。

本文将从五个维度拆解这份报告的核心洞察，并探讨那些尚未被充分定价的关键变量。

1. AI超级周期的需求侧已经明确，但供给侧的结构性变化才是日本市场的独特机会

全球AI投资热潮并非新鲜事，但JPM特别强调了一个容易被忽视的视角：AI推理正在取代训练成为主要增长领域。美国超大规模云服务商在数据中心的投资持续超预期，直接拉动了对半导体芯片、存储器和电子元件的需求。这一轮需求扩张的广度和持续性，已经超出了市场此前的预期。

对日本市场而言，这意味着什么？日经225中约50\%的市值与AI相关，其中半导体和半导体生产设备（SPE）公司是核心。报告提到，铠侠（Kioxia）等SPE企业和电线电缆制造商均报告了强劲的盈利增长。JPM已将2026财年TOPIX每股收益预测从年初的198日元上调至当前的234日元，上调幅度超过18\%。

这里的关键洞察是：日本在AI硬件供应链中的地位正在从“配套”转向“核心”。过去，市场对日本半导体产业的认知停留在材料、设备等上游环节，但这一轮AI投资周期中，日本企业在存储芯片、先进封装、高精度制造设备等领域的话语权正在显著提升。这种变化不是短期的周期性波动，而是全球半导体产业链重新分工的结果。

对于投资者而言，这意味着评估日本科技股时，不能再用传统的“周期股”框架。当AI资本开支从训练转向推理，从云端扩展到边缘，日本供应链企业的业绩可见度正在从1-2年延长至3-5年。这正是JPM敢于在估值已经上涨后继续上调目标价的根本原因。

2. 估值压力已经释放，盈利增长正在接棒成为市场上涨的主要驱动力

上半年日经225上涨39\%，难免引发市场对估值过热的担忧。但JPM的数据显示，TOPIX的12个月远期市盈率已从年初的17.5倍回落至6月中旬的16.6倍。这一估值水平不仅低于历史均值，也显著低于美股和欧股。

估值回落的逻辑并不复杂：盈利增长的速度超过了股价上涨的速度。报告中的EPS预测显示，2026财年TOPIX每股收益预计达到234日元，较2025财年的210日元增长11.4\%。这一增速已经高于市场一致预期的10.7\%，意味着JPM对日本企业盈利能力的判断比市场更为乐观。

更重要的是，这种盈利增长并非来自一次性因素。报告预测2027财年EPS将继续增长12.0\%至262

[... middle omitted ...]

性的突破。如果这一趋势得以延续，日本股市的估值中枢有望获得系统性抬升。

3. 资金流入的结构正在发生根本性变化，个人投资者和养老金将成为新的增量来源

海外投资者一直是日本股市的重要边际定价者，但JPM指出，国内资金的结构性流入正在成为新的催化剂。报告提到，个人投资者、养老金基金和金融机构对日本股票的资金流入正在变得更加清晰。

这一判断的背景是日本政府推出的新NISA（小额投资免税制度）改革，以及养老金基金（GPIF）对国内股票配置比例的提升。这些政策变化正在改变日本家庭的资产配置行为。长期以来，日本家庭金融资产中股票占比极低，大量资金沉淀在现金和存款中。随着通胀环境的变化和税收优惠政策的推动，这一格局正在发生微妙但重要的转变。

从市场影响的角度看，国内资金的持续流入意味着日本股市的投资者基础正在变得更加多元化。过去，市场高度依赖外资的流向，一旦全球风险偏好下降，日本股市往往首当其冲。但如果国内资金能够形成稳定的增量，市场的波动性有望降低，估值支撑也会更加坚实。

\subsection{[R009] JPM：市场低估的不只是加息幅度，更是货币政策沟通机制的历史性转折}
{\small\color{refgray}归类：宏观与利率：美联储沟通范式切换与全球央行分化}

JPM：市场低估的不只是加息幅度，更是货币政策沟通机制的历史性转折

这份JPM在FOMC会议后数小时发布的研报，表面上在分析利率曲线和通胀预期的短期移动，但真正值得关注的判断藏在细节里：美联储正在经历自格林斯潘时代以来最彻底的沟通方式变革，而市场对此的定价远未完成。

2026年6月的这次FOMC会议，新任主席Warsh的首秀，被JPM定性为“鹰派”。但鹰派本身不是新闻。真正的新信号是：政策声明被完全重写，不再提供任何前瞻指引，转而采用一句近乎宣言式的表述——“委员会将实现价格稳定”。这听起来像回归格林斯潘时代的含蓄与不可预测性。JPM的NLP模型将这份声明评为一年来最鹰派，新闻发布会评分则是2024年5月以来最高。

但比鹰派更重要的，是这种沟通方式转变所隐含的结构性后果。如果前瞻指引被系统性削减，市场将不得不把更多权重放在每一个“事件日”上——每一次FOMC会议、每一次讲话、每一个数据发布，都可能引发比以往更大的波动。这不仅仅是短期交易员的烦恼，它意味着期限溢价的系统性重估，进而影响所有以无风险利率为锚的资产定价。

以下是我们从这份研报中提炼出的五个层次判断，以及一个尚未被充分讨论的关键问题。

1. 鹰派点阵图只是表象，真正重要的是“不对称的政策不确定性”正在被重新定价

市场对10月会议加息25bp的完全定价，以及对明年春季累计超过45bp加息的预期，看起来是直接对点阵图的反应。2026年点阵图中位数显示12.5bp加息，18位委员中有6位预计两次或更多加息。但JPM的分析揭示了一个更深层的信号：从SOFR期货期权推断的隐含概率分布，不仅整体向鹰派方向移动，而且“双边尾部均增强”——意味着市场认为极端情景（无论是更鹰派还是更鸽派）的概率都在上升，分布变得更加对称。

这才是真正的变化。过去几年，市场习惯了美联储提供清晰路径的环境，政策不确定性被系统性压低。现在，Warsh正在主动拆除这套导航系统。JPM明确指出，“减少前瞻指引可能增加波动性，进而提高期限溢价和名义收益率”。这不是短期现象，而是货币政策操作框架的转向。

因此，当前收益率曲线的定价可能仍然低估了这种结构性变化的影响。即使加息路径本身没有超出预期，但“不确定性溢价”的重置才刚刚开始。

2. 通胀市场的反应被能源价格掩盖，但真正的下行风险来自美联储对通胀目标态度的微妙变化

与名义利率市场的剧烈波动相比，通胀市场的反应相对温和：5年、10年、30年盈亏平衡通胀率分别仅下降3bp、2bp和1bp。但实际收益率的上升幅度要大得多——5年期实际收益率飙升10bp。这组对比告诉我们，市场目前更关注实际利率的重新定价，而非通胀预期的崩溃。

JPM提出了一个值得注意的论点：Warsh有机会利用能源价格下跌来推动短期通缩叙事。布伦特原油从118美元高位跌至80美元以下，这为美联储提供了鸽派转向的借口。但Warsh选择不去抓这个机会，反而强调“过去5年未能实现2\%目标”。这意味着，美联储可能正在将通胀目标从一个“对称性目标”重新解读为一个“必须达成的硬约束”。

这里的关键尚未被市场充分消化：如果美联储开始将“低于2\%的通胀”视为与“高于2\%”同等严重的问题，那么过去几年市场习惯的“通胀中枢上移”叙事就需要重新审视。JPM的模型显示，5年期盈亏平衡通胀率目前已经比其公允价值模型便宜约44bp，接近3个标准差——这要么是市场定价错误，要么是模型未能捕捉到美联储态度转变带来的结构性通缩压力。

3. 收益率曲线平坦化交易有坚实的模型支撑，但执行节奏取决于对Warsh沟通策略的进一步观察

JPM明确推荐“维持10s/30s flatteners”，即做平10年期与30年期国债收益率曲线的交易。其依据是：即使在当前曲线已经显著平坦化之后，10s/30s曲线

[... middle omitted ...]

题在于，这种压力是渐进释放还是一次性冲击？如果市场开始预期每月都会有“鹰派惊喜”，那么曲线平坦化的速度可能会比JPM模型假设的更快。反之，如果Warsh选择在接下来的讲话中释放一些安抚信号，那么当前曲线可能已经过度反映了这种不确定性。这是一个需要持续跟踪的变量。

4. 资产负债表正常化仍在2027年之前，但“将到期国债转向短债”的讨论可能成为一个被低估的供给端变量

关于缩表，这次FOMC会议没有提供实质性新信息。Warsh表示资产负债表工作组将在2026年下半年进行评估，这意味着任何实质性缩表行动可能要到2027年。这一点与市场预期基本一致。

但JPM提出了一个值得关注的情景：SOMA投资组合中约1.9万亿美元的付息国债将在2026年至2030年间到期。虽然美联储不太可能将所有到期资金全部再投资于短期国债，但“将一半到期资金转向短期国债”的假设场景已经出现在研报中。如果这一情景成真，SOMA投资组合的加权平均久期将在2029年中期收敛至与整体国债市场相当的水平。

\subsection{[R010] Bernstein：瑞士手表出口反弹只是表象，真正的复苏信号尚未到来}
{\small\color{refgray}归类：风险偏好与资金流向：极端集中与泡沫特征重现}

Bernstein：瑞士手表出口反弹只是表象，真正的复苏信号尚未到来

全球奢侈品行业正经历一场“数据幻觉”。2026年5月，瑞士手表出口同比微增0.4\%，看似走出了4月-16.6\%的深跌。但Bernstein这份最新研报揭示了一个更复杂的现实：调整日历效应后出口增长10.4\%，这个数字本身仍掩盖了结构性脆弱。真正值得关注的不是月度数据反弹，而是支撑复苏的底层逻辑——美国市场的基数效应、中国需求的缓慢改善、以及中东地缘风险对奢侈品消费的潜在冲击。

这份报告的核心判断是：全球奢侈品需求的复苏轨迹仍充满不确定性，投资者不应被短期数据波动所迷惑。Bernstein明确指出了两股波动源：一是消费者在脆弱的宏观环境和日益紧张的地缘政治中做出的实际需求决策；二是短线投资者在板块的多空博弈，放大了涨跌幅。在这样的背景下，报告建议采取防御性姿态。

1. 美国市场的“基数游戏”是当前出口反弹的主要推手，而非需求强劲

5月瑞士手表对美出口同比增长12.3\%，这个数字本身并不代表美国消费者突然变得慷慨。关键在于比较基数的剧烈波动：去年同期的基数仅为-25\%，而4月的基数高达+149\%。这种基数效应造成的剧烈摆动，使得月度同比数据几乎失去了解读意义。

Bernstein的数据显示，从过去12个月的滚动数据看，对美出口同比下降19\%，虽然较4月的-21\%有所改善，但仍处于收缩区间。更关键的是，报告明确指出，今年剩余月份美国的比较基数将“持续友好”——除了7月因去年提前发货导致基数高达+45\%外，其余月份均为两位数的负基数。这意味着未来几个月美国市场的出口数据很可能继续“看上去不错”，但投资者需要区分这是真正的需求复苏还是统计幻觉。

这里蕴含着一个重要洞察：美国消费者对关税相关涨价的接受度，可能正在被高估。报告提到公司评论称美国消费者“已经接受了关税带来的价格上涨”，但这种接受度是否可持续，在消费者信心尚未完全恢复的背景下，仍需时间验证。

与美国的剧烈波动形成对比，大中华区的表现呈现出一种“温和但真实”的改善轨迹。5月对香港出口同比增长3.4\%，对中国内地出口下降21.4\%，这两个数字的恶化主要源于更高的比较基数。从滚动12个月数据看，对大中华区出口同比下降0.9\%，较上月的-1.4\%继续改善。

这个趋势值得关注，因为它反映了中国奢侈品消费者信心的逐步修复。但Bernstein并未给出乐观预期——报告使用的措辞是“逐步改善”和“缓慢而稳定的改善”，而非“强劲反弹”。这种谨慎措辞背后，隐含着一个尚未被充分讨论的问题：中国消费者的奢侈品购买行为是否已经发生了结构性变化？从“炫耀性消费”向“体验式消费”的转变，是否意味着手表等硬奢品的需求弹性正在变大？

报告没有直接回答这个问题，但它提供的LTM数据趋势——从-1.4\%到-0.9\%——暗示改善速度正在放缓。这可能是中国消费者信心修复的“天花板”正在显现的信号。

3. 价格带分化揭示了一个被忽视的趋势：中端市场的“空心化”正在加剧

5月的数据显示，不同价格带的出口表现出现了明显分化。定价在200-500瑞士法郎和3000瑞士法郎以上的手表，出口值分别增长23.6\%和22.2\%；而200瑞士法郎以下和500-3000瑞士法郎区间的产品则分别下降4.5\%和1.4\%。

这个分化背后隐藏着两个重要信号。首先，高端市场（3000瑞士法郎以上）的韧性再次得到验证，这与奢侈品的“口红效应”逻辑一致——高净值人群在经济不确定性时期的消费行为相对稳定。其次，中端市场（500-3000瑞士法郎）的疲软值得警惕，这个价格带通常对应着“可负担奢侈品”或“轻奢”定位，更容易受到消费者信心波动的影响。

但更值得深思的是200瑞士法郎以下市场的持续萎缩。这个价格带的产品通常面向大众市场或入门级消费者，其

[... middle omitted ...]

组合管理将产生深远影响。

Bernstein在报告中提出了一个值得警惕的风险点：中东地区。2025年，中东是奢侈品行业“罕见的增长点”，但5月的数据显示，对阿联酋和沙特阿拉伯的出口分别同比下降13.5\%和5.9\%。

报告谨慎地指出，“如果冲突能够迅速解决，短期直接影响可能有限”，但“间接影响——如旅行减少、油价上涨/通胀上升、以及消费者信心下降——可能威胁到复苏的早期迹象”。这个判断的言外之意是：中东地缘政治风险对奢侈品行业的影响，可能不是通过直接销售渠道，而是通过更广泛的宏观经济传导机制。

这里有一个尚未被充分量化的风险：如果中东冲突持续，国际旅行模式的变化将如何影响奢侈品消费？中东不仅是重要的消费市场，也是全球奢侈品零售网络的关键节点（如迪拜的转口贸易）。报告提到法国出口激增57\%的部分原因是“自2025年12月以来作为物流枢纽的角色”，这恰恰印证了供应链和渠道结构变化正在重塑贸易数据。中东局势的变化可能加速这种渠道重构，但具体影响路径和程度，报告并未展开。

\subsection{[R011] GS：地震勘探行业正在进入一个被低估的、有纪律的多年度上行周期}
{\small\color{refgray}归类：能源与大宗商品：供给约束下的时间错配与地缘风险折价}

GS：地震勘探行业正在进入一个被低估的、有纪律的多年度上行周期

市场对油服板块的关注，大多停留在油价波动和短期资本开支上。但一份来自GS的研报，通过与Viridien（原CGG）管理层的深度对话，揭示了一个更关键的结构性信号：全球地震勘探行业，正在经历一场从“资本纪律”到“资源重置”的范式转换。这并非 年那种由高油价驱动的狂欢，而是一个由储量寿命衰竭、能源安全焦虑和技术壁垒共同支撑的、更可持续的周期起点。

这份报告最值得看的判断是：市场真正低估的不是需求是否回升，而是供给侧的再定价能力以及行业竞争格局的质变。地震勘探行业，尤其是海洋拖缆领域，已经从“产能过剩的苦海”变成了“寡头定价的棋局”。而AI技术的渗透，非但没有削弱核心企业的护城河，反而让高质量成像数据成为了更稀缺的资产。

1. EAGE 2026的“爆满”不是噪音，而是储量焦虑的集体投票

2026年6月在阿伯丁举行的第86届EAGE年会，创下了超过6000人的历史参会纪录。这不仅仅是行业聚会的规模扩大。GS报告明确指出，C-level管理层的参与度较2025年显著提升。背后驱动力是“储量替换”和“储量寿命”这两个最根本的行业指标，重新回到了董事会的核心议程。

这意味着什么？过去十年，行业被“资本纪律”和“股东回报”的叙事主导，上游企业普遍削减勘探开支，专注于短周期、低风险的棕地项目。但代价是储量寿命的持续下降。GS数据显示，全球主要油企的储量寿命从2012年的约55年，骤降至2026年预计的约20年。这是一个系统性的“寅吃卯粮”。EAGE的火爆，正是这种深层焦虑的集体宣泄。它告诉我们，行业共识已经形成：不勘探，就没有未来。

2. 前沿勘探回归，但“手艺”已经断代，这构成了先发者的时间窗口

报告中最具洞察力的判断之一是“前沿勘探的回归”。根据Viridien的说法，历史上约80\%的地震活动集中在近场区域。但近场勘探已无法弥补储量寿命的缺口。客户预算中，针对前沿盆地的再处理工作正在增加。

更关键的是，GS敏锐地捕捉到了一个“能力断层”：主要石油公司在经历了多年的投资不足后，已经失去了前沿盆地勘探的专业知识，重建需要时间。这意味着，那些在过去十年中坚持投入、积累了大型前沿盆地多客户数据库的地震服务公司，如Viridien，拥有了一个宝贵的“时间窗口”。客户需求已经转向，但客户自身的执行能力尚未恢复，这为拥有数据资产和成像技术的服务商创造了议价权。报告点名了乌拉圭、巴西赤道边缘、安哥拉、埃及、利比亚等地区，这些不再是远期概念，而是正在签署谅解备忘录的现实战场。

如果说前沿勘探是长期叙事，那么“再处理”就是当前最确定的增长引擎。报告指出，客户预算中，针对前沿盆地的再处理工作正在增加，这通常发生在钻探之前，目的是尽早锁定数据。

这背后的含义是：在资本开支依然谨慎的背景下，油公司选择了一种“低成本、高信息密度”的勘探策略。他们不急于立刻去打昂贵的野猫井，而是先通过更先进的成像技术，重新解读已有的地震数据，以降低钻探风险。这对地震服务公司是双重利好：一是带来了直接的再处理收入；二是强化了其作为“地下信息入口”的平台价值。GS报告提到，国家石油公司（如ADNOC）越来越多地将从成像到圈闭的完整工作流程外包给地震服务商，而不是自己建立内部团队。这标志着行业价值链的重塑。

4. 船队供给已经“硬约束”，定价权正在向寡头转移，但价格传导仍有滞后

这是报告最具投资启示的部分。海洋拖缆船队已从2014年高峰期的约60艘活跃船只，整合至目前的约16-17艘，由TGS和Shearwater主导。尽管目前船队利用率尚未绷紧，日费率稳定，但GS和Viridien都认为，一旦需求上行，供给端的弹性将非常有限。闲置产能可以重启，但存在时滞。

这意味着什么？行业已经从一

[... middle omitted ...]

步强化了拖缆公司在成本效益上的竞争优势。

5. AI不是颠覆者，而是护城河的加固者：高质量成像数据的稀缺性正在被重估

面对AI浪潮，市场普遍担忧技术替代风险。但GS报告给出了一个反直觉的结论：AI在地震勘探中是“增值的”，而非“替代的”。核心逻辑在于，物理驱动的波动方程成像问题，AI无法解决，这依然是一个需要巨大HPC（高性能计算）算力的迭代过程。

AI的真正价值在于加速和优化：去噪、工作流自动化、交付提速。报告提到，Viridien的云交付平台已将TB级成像数据的交付时间从4周缩短至大约1天。这极大地提升了客户粘性。更关键的是，最终客户正在将自己的AI工具叠加在地震图像之上，以优化勘探前景。这意味着，图像质量成为了AI工具价值的“瓶颈约束”。高质量的基础数据越稀缺，其价值就越高。那些拥有庞大、高质量多客户数据库和专有成像算法的公司，不仅不会被AI取代，反而会成为整个AI勘探生态的“基础设施”。

报告尚未完全回答的关键问题：这一轮上行周期的价格弹性究竟有多大？

\subsection{[R012] JPM：市场低估了供给侧的再定价，而非需求端的边际变化}
{\small\color{refgray}归类：能源与大宗商品：供给约束下的时间错配与地缘风险折价}

中国5月经济数据再度走弱，这本身并不令人意外。真正值得关注的信号，是JPM在最新一期亚洲基础材料研报中提出的一个隐含判断：当前大宗商品市场的定价逻辑，正在从“需求预期驱动”切换为“供给约束驱动”。对于铜、铝、煤炭这三个核心品种，市场对需求侧的悲观已经充分定价，但供给侧的结构性收紧——无论是地缘政治冲击、产能监管趋严，还是全球供应链的物理性阻塞——尚未被完全计入资产价格。

这份研报的核心价值不在于罗列数据，而在于它提供了一套清晰的决策框架：在需求侧无法提供方向感的时候，投资者应该如何重新锚定大宗商品的定价基准。

1. 霍尔木兹海峡的重新开放，可能成为金属板块最被低估的战术催化剂

研报最引人注目的判断，是JPM认为市场对美伊局势缓和的潜在影响定价不足。自冲突爆发以来，铜和黄金相关股票的跌幅已经相当显著——紫金矿业H股下跌26\%，紫金黄金下跌49\%，而同期恒生指数仅下跌8\%。这种跌幅反映了市场对地缘风险溢价的全面计入。

但JPM提出的逻辑是：如果霍尔木兹海峡重新开放，被压抑的风险偏好将首先释放到那些受冲突冲击最严重的资产上，铜和黄金的股价反弹潜力最大。这个判断的关键支撑在于，铜和黄金的基本面并未因冲突而实质性恶化，下跌更多来自情绪层面的风险折价，而非供需结构的永久性破坏。

这实际上是一个“均值回归”的战术逻辑：市场对尾部风险的定价往往过度，而一旦尾部风险消散，资产的修复速度可能快于基本面改善的速度。对于机构投资者而言，这意味着当前可能是布局铜和黄金相关标的的时间窗口——但前提是承认自己无法精准判断地缘政治事件的时点，只能以赔率而非胜率为决策依据。

2. 铝的短期压力与中期矛盾：印尼产能被过度定价，而真实缺口被低估

铝是这份研报中分析层次最丰富的品种。短期来看，霍尔木兹海峡的重新开放可能释放中东地区的积压库存，对LME铝价形成压力——当前价格约在3360美元/吨附近。但JPM的中期观点仍然是结构性的看多。

核心逻辑在于：市场似乎已经提前将印尼2027年的新增产能计入了当前价格，而忽略了2026年可能出现的约170万吨的初级铝缺口。这是一个典型的“远期供给预期压制近端价格”的定价错误。如果2026年的真实缺口兑现，而印尼产能的释放节奏慢于预期，铝价可能面临一轮向上的重估。

这里需要特别注意的是，JPM的大宗商品团队仍在预测2026年存在供应缺口，这意味着当前的价格水平可能并未充分反映2026年的供需紧张。对于铝行业投资者而言，关键观测变量不是短期的价格波动，而是库存去化的速度和海外同行的估值比较。研报明确提到，中国铝业和中国宏桥是核心受益标的，估值重估和去库存进展是需要持续跟踪的催化剂。

锂的价格在16.5万元/吨附近获得了支撑，但JPM并未给出强烈的方向性判断。研报的表述是“信号仍然混杂”，预计价格将在当前水平附近维持区间震荡，波动由供给端的响应速度、库存数据和下游需求信号共同驱动。

这实际上是一个“没有判断的判断”，但恰恰反映了当前锂市场的真实状态：供给端的出清速度不够快，需求端的新能源汽车渗透率增长放缓，两者之间尚未形成明确的再平衡信号。对于投资者而言，锂的交易机会可能更依赖于对供给端事件的捕捉——矿山的减产公告、冶炼厂的检修计划、以及库存数据的意外变化——而非对需求趋势的宏观判断。

4. 煤炭：山西事故后的供给收缩正在被市场重新定价，但恢复速度可能慢于预期

煤炭是这份研报中供给侧逻辑最清晰的品种。山西矿难后，炼焦煤价格持续走高，JPM预计价格将维持高位至2026年三季度，这意味着钢铁企业的利润空间将持续受到挤压。

关键数据点在于：截至6月15日，虽然约70\%的停产产能（97座煤矿，合计1.16亿吨）已经按照Mysteel的数据恢复生产，但许多煤矿仍以保守方式运营，实际产量恢复可能滞后于

[... middle omitted ...]

3200万吨的供给缺口，这个数字可能仍然偏保守，因为它尚未充分计入印尼煤炭出口政策收紧的全球溢出效应。

对于煤炭股的投资逻辑，JPM的框架是清晰的：在供给缺口持续的背景下，龙头企业的定价权和盈利稳定性将优于市场预期。但风险在于，如果钢铁需求进一步走弱，焦煤价格的强势可能难以持续——这是一个需要持续跟踪的宏观变量。

尽管JPM在供给侧的分析非常扎实，但报告对需求侧的判断仍然存在不确定性。5月的新开工面积同比下降25\%，房地产投资下降25\%，固定资产投资总额同比下降17\%——这些数字本身已经足够说明问题，但市场更关心的是：这些数据何时能够企稳？

JPM的房地产分析师Karl Chan预计全年新开工面积下降18\%，竣工面积下降19\%，这意味着下半年新开工面积的同比降幅可能收窄至15\%，竣工面积降幅收窄至16\%。但问题在于，这个预测是基于“土地成交面积在2025年下半年和2026年一季度持续疲软”的假设，如果房地产销售端没有出现超预期的回暖，新开工的企稳可能比预期更晚。

\subsection{[R013] JPM：半导体设备市场的下一个结构性拐点不在需求，而在供给侧的产能释放节奏}
{\small\color{refgray}归类：AI与算力：从训练到推理的架构迁移与定制芯片崛起 / 半导体与功率器件：供给驱动的涨价周期与定制化浪潮}

JPM：半导体设备市场的下一个结构性拐点不在需求，而在供给侧的产能释放节奏

这份JPM刚刚发布的研报，核心判断只有一个：全球晶圆厂设备（WFE）市场正在经历一次被低估的上行重估。JPM将2026年WFE市场增速预测从21\%大幅上调至28\%，2027年从18\%上调至29\%，并首次给出2028年16\%的增长指引。这不是一次温和的修正，而是一次系统性的预期上调。

但真正值得关注的，不是增速数字本身，而是驱动这一轮上调的结构性力量正在发生质变。AI需求从训练向推理的迁移，正在重塑整个半导体资本开支的分配逻辑——内存芯片的投资权重正在从过去的个位数至15\%区间，跃升至2026年的约50\%。这意味着，过去几年市场习惯的“AI=GPU+ASIC”叙事，正在被“AI=加速器+高带宽内存+先进封装”的全栈逻辑取代。

更关键的是，这份报告揭示了一个被多数投资者忽视的变量：洁净室产能的瓶颈正在缓解。此前市场对2026年设备交付的担忧，很大程度上源于芯片制造商洁净室建设进度滞后。但JPM的调研显示，通过工厂空间的创造性利用和洁净室扩建的正式启动，这一瓶颈正在被打破。从2027年开始，随着洁净室大规模投产，需求将迎来真正的释放窗口。

以下是我们从这份研报中提炼出的五个核心洞察，每一个都指向一个正在被重新定价的产业环节。

1. 云厂商资本开支的“超级周期”正在从预期走向现实，且规模远超历史任何时期

JPM对全球前四大美国云服务提供商（谷歌、亚马逊、微软、Meta）的资本开支预测，做了幅度惊人的上调：2026年增速从63\%上调至80\%，2027年从40\%上调至50\%。在绝对值上，这意味着2026年同比增量将超过2500亿美元，2027年超过2850亿美元——这是有史以来最大的年度增幅。

支撑这一判断的核心依据，是北美地区的规划电力容量。截至2026年第一季度末，规划电力容量已超过175GW，JPM预计到年底将突破205GW。与之对应，装机容量将从2025年的48GW，分别提升至2026年的65GW和2027年的85GW。

一个更值得警惕的财务信号是：到2027年，前四大云厂商的资本开支占经营现金流比例将从过去几年的25\%-30\%，飙升至94\%。这意味着，这些公司几乎将全部经营现金流投入资本开支。这种极端配置在历史上从未持续过，但它也反过来说明，AI基础设施的投资回报预期已经强烈到足以让管理层做出如此激进的财务决策。

对于设备供应商而言，这意味着订单可见度正在进入一个前所未有的阶段。JPM在5月的TMT会议上调研发现，几乎所有半导体设备公司都表示，订单、积压和客户升级请求在财报发布后的几周内进一步加速。KLA明确指出，当前对2027年需求的可见度是过去十多年来最高的，并认为2027年的WFE市场将大于2026年。

2. 内存芯片正在成为AI资本开支的新主轴，其投资权重将发生数量级跃升

这是整份报告最反直觉、也最容易被忽视的判断。过去两年，市场讨论AI资本开支时，焦点几乎全部集中在逻辑芯片（GPU、ASIC）和先进封装上。但JPM的数据显示，内存芯片的投资权重正在从过去的个位数至15\%区间，跃升至2026年的约50\%。

这一变化的驱动力来自两个方向。第一，AI推理工作负载的兴起正在改变内存需求结构。推理阶段对内存带宽和容量的要求远高于训练阶段，这直接推动了对高带宽内存（HBM）和DRAM的强劲需求。第二，内存芯片价格正在经历历史级别的上涨——NAND和DRAM价格均出现三位数同比增长，这为芯片制造商提供了充足的现金流来进行产能扩张。

JPM预计，未来三年内存芯片行业的累计资本开支将达到4500亿美元，较此前预测的3000亿美元大幅上调。其中，DRAM资本开支为3640亿美元，NAND闪存为860亿美元。这一数字背后隐含

[... middle omitted ...]

具备定价能力，被列为首选股。爱德万测试（Advantest）则在后端测试领域持续受益于高需求。

3. 台积电的资本开支正在进入加速通道，先进制程短缺将延续至2028年

JPM将台积电2026年资本开支预测从此前的540亿美元上调至560亿美元（同比增长37\%），2027年上调至650亿美元（同比增长16\%），2028年首次给出720亿美元的预测（同比增长11\%）。这一系列上调的核心逻辑是：先进制程（N5及以下）的供应短缺将持续到2027年甚至2028年初，台积电正在加速产能扩张以应对这一缺口。

具体来看，台积电N5及以下制程的产能预计将在 年间以25\%的复合年增长率扩张，其中大部分增量集中在 年。N3制程方面，台积电计划从2027年上半年到2028年，启动三座新工厂的产能爬坡（台南Fab 18 P9、亚利桑那州二期、JASM第二工厂）。N2制程则拥有强劲的客户管线，包括iPhone处理器、AMD的Venice和MI450、谷歌的TPU以及英伟达的Feynman。

\subsection{[R014] JPM：AI定制芯片正在从“备选方案”变成“主导逻辑”}
{\small\color{refgray}归类：AI与算力：从训练到推理的架构迁移与定制芯片崛起}

过去两年，市场对AI基础设施的讨论几乎完全围绕英伟达GPU展开。算力即权力，而英伟达是唯一的铸币厂。但这份JPM最新发布的半导体深度报告揭示了一个正在发生的结构性转折：到2027年，AI加速器的出货量中，定制芯片（ASIC/XPU）的占比将从2025年的32\%跃升至53\%，首次超过GPU。这不是边缘替代，而是产业逻辑的根本切换。

报告的核心判断是：当AI从训练走向推理，从少数巨头试验走向大规模部署，定制芯片的经济性和系统级优势正在重塑整个计算架构的供应链。这不是英伟达的“被颠覆”，而是AI计算市场本身在快速分化和成熟。理解这一轮ASIC浪潮的驱动力、竞争格局和未解问题，比争论“英伟达是否见顶”更有意义。

以下是这份报告的五个关键洞察，以及一个报告尚未完全回答的关键追问。

1. 定制芯片的崛起不是技术竞赛，而是规模经济与软件护城河的综合博弈

JPM将芯片市场清晰划分为三类：通用芯片（GP）、标准应用芯片（ASSP）和定制芯片（ASIC）。AI GPU属于ASSP，而Google TPU、Amazon Trainium、Microsoft Maia则属于ASIC。后者的核心优势不在于单点性能超越GPU，而在于系统级的总拥有成本（TCO）优化。

报告提供了一个关键对比：Google/Broadcom的TPUv7 Ironwood在FP8计算性能上达到4614 TFLOPS，与英伟达B200 Blackwell的4500 TFLOPS相当，但ASIC的单价仅约13000美元，远低于B200的35000美元。这意味着每美元获得的算力（TFLOPS per \$）上，ASIC是GPU的2.7倍。在功耗效率（TFLOPS per Watt）上，两者几乎持平。

这组数据揭示了一个被市场低估的真相：当AI部署进入万卡甚至十万卡集群规模时，资本开支的优化不再只看单芯片性能，而是看每单位算力的总成本。定制芯片通过剪裁不必要的通用计算单元、优化特定工作负载的内存层次，实现了显著的性价比优势。这不是技术的落后，而是工程上的精准匹配。

对于产业决策者的含义是：未来AI基础设施的采购决策将越来越像今天的数据中心采购——不是买最贵的CPU，而是买最适合工作负载的芯片组合。英伟达的优势在于其CUDA生态的粘性，但ASIC正在用自己的经济学逻辑侵蚀这个护城河。

2. 博通和马威尔正在成为AI时代的“芯片代工厂”，而非简单的芯片设计公司

报告披露的数据显示，博通目前在高性能ASIC市场占据80-85\%的份额，马威尔占据10-12\%。但更重要的是，这两家公司正在构建的商业模式，远比“为客户设计芯片”更深刻。

博通已经建立了从2nm/3nm计算die、100/200Gbps I/O die、2.5D/3D SOIC芯片堆叠到先进封装和机架级系统设计的完整参考平台。报告指出，博通目前拥有超过100个累计的7nm/5nm/3nm/2nm设计项目，其中5个主要AI ASIC客户包括Google、Meta、OpenAI、SoftBank/Arm和一家亚洲超大规模客户。马威尔则拥有70多个累计设计项目，核心客户包括Amazon和Microsoft。

这些数字背后是商业模式的重构。传统半导体公司要么卖标准芯片，要么卖IP授权。但博通和马威尔正在做的事情，是成为超大规模客户的“芯片代工厂”——客户拥有软件栈和应用层，芯片公司提供从架构设计、IP组合、先进封装到供应链管理的完整能力。这种合作模式下，客户获得差异化，芯片公司获得长期、高粘性的收入流。

报告估算，博通在FY26的AI相关收入将达到600亿美元以上，较FY25的约200亿美元增长200\%。马威尔的数据中心收入在CY26预计达到93亿美元，CY27达到146亿美元。这些数字如果兑现，意味

[... middle omitted ...]

而ASIC/XPU占比从32\%上升到53\%。

具体来看，Google TPU的出货量将从2025年的180万颗激增至2027年的800万颗，Amazon Trainium/Inferentia从150万颗增至330万颗，Microsoft Maia从零起步到2027年的70万颗，Meta MTIA从零到50万颗。这些数字意味着，到2027年，仅Google一家的AI ASIC出货量就将接近英伟达GPU总出货量（990万颗）的80\%。

但报告也隐含了一个关键约束：先进封装产能。无论是博通的2.5D/3D SOIC封装，还是马威尔的定制HBM架构，都高度依赖台积电的CoWoS等先进封装技术。报告提到博通正在开发2nm/3nm芯片堆叠和先进基板制造能力，但产能爬坡时间表定在CY26。这意味着2026年的出货量爆发可能受限于封装产能的释放速度，而非设计本身。这是一个需要持续跟踪的供应侧变量。

4. 定制HBM和芯片互连正在成为ASIC的“隐形护城河”，而非单纯的计算核心

\subsection{[R015] 摩根斯坦利：市场低估的不是下行风险，而是政策“微调”而非“转向”的定力}
{\small\color{refgray}归类：中国地产与消费：量稳价跌与盈利预期重置}

摩根斯坦利：市场低估的不是下行风险，而是政策“微调”而非“转向”的定力

四、五月的数据让一个事实变得无法回避：中国经济的“双速”格局正在加深。出口和生产端依然保持韧性，但国内需求——消费、投资、甚至石油消费——出现了明显的边际走弱。

市场对此的第一反应通常是担忧。担忧经济失速，担忧政策被迫急转弯。但摩根斯坦利在本期《China Musings》中提出了一个更值得深思的判断：当前的下行并不指向需要政策“转向”的急剧恶化，而是触发“微调”的紧迫性上升。报告的标题本身就很能说明问题——“Beijing Will Fine-tune, Not Pivot”。

这个判断的分量，不在于它确认了经济放缓，而在于它划出了一条重要的政策边界：北京不会因为短期数据波动就放弃既定的结构性转型路径。对于投资者而言，真正需要理解的不是“经济好不好”，而是“在4.5\%-5\%的增长目标与K-shaped经济现实之间，政策制定者会选择什么工具、在什么节奏上出手”。

四五月的数据确认了一个K-shaped格局的深化：生产与出口坚挺，消费与投资软化。但摩根斯坦利提醒，不能简单地用同比增速的回落来判断放缓的严重程度。

固定资产投资（FAI）的减速，部分源于政策驱动。自2月底以来，中央强调“正确政绩观”，促使地方政府将重心从新项目转向化解隐性债务。这意味着FAI数据的走弱，有主动调控的成分，而非纯粹的终端需求崩溃。零售数据同样需要拆解：5月同比转负，很大程度上受到去年“以旧换新”补贴政策高基数的影响。以两年复合年化增长率（2Y CAGR）来看，5月甚至略好于4月。

但报告同时指出，放缓不仅仅是基数效应。即便看两年复合增速，4-5月相比一季度仍有明显的台阶式下降。劳动力市场温和、房地产市场调整仍在持续、油价相关的出行消费承压，这些因素共同压低了消费意愿。而出口的强劲并未充分传导至就业市场——自动化和资本密集度的提升，削弱了出口对就业的拉动弹性。国家统计局的消费者信心指数在短暂回升后再次走软，印证了这一点。

*这组数据合在一起的含义是：当前的经济放缓，是结构性转型中的“阵痛”，而非周期性衰退的开始。政策制定者对此有预期，也有预案。**

2. 石油视角的启示：经济减速可控，但K-shaped特征被进一步强化

原油进口的大幅下滑，引发了部分投资者对中国经济活动萎缩的担忧。摩根斯坦利对此提供了一个更立体的解读。

消费者确实在削减非必要的石油消费，包括燃油车使用和航空出行。但工业活动层面，存在多重缓冲：商业原油库存、稳健的煤炭进口与国内产量、以及依然强劲的发电量。发电量数据尤为关键——它表明工业生产的物理活动并没有像石油进口数据暗示的那样大幅收缩。

*这意味着，石油数据的走弱更多反映了消费端的行为变化（出行减少），而非工业端的崩溃。这恰恰是K-shaped经济的微观体现：生产端有韧性，消费端在收缩。这种不对称性，决定了政策干预不能一刀切，而必须精准地作用于需求侧，尤其是消费信心和就业市场。**

3. 政策的回应将是“微调”而非“转向”，战略基础设施是首选工具

二季度GDP追踪值大约在4.4\%左右。如果放任不管，全年增速可能触及甚至低于4.5\%-5\%目标区间的下限。这无疑提升了政策紧迫感。

但摩根斯坦利明确指出，近期的政策回应将是“微调”而非“大转向”。具体来说，重点将放在加快已批准预算的执行上——年度政府债券额度中仍有约60\%尚未使用。更关键的是，用于基础设施的8000亿元人民币准财政融资工具（以政策银行债券发行为代理变量）的落地进度仍然缓慢。

北京首选的国内需求支持工具，依然是战略基础设施。报告特别提到了“六张网络”框架：水网、新型电力系统、算力网络、下一代通信、城市地下管网和物流基础设施。其中，AI算力网络、数据中心和智能电网最

[... middle omitted ...]

解。随着美国-伊朗和平协议的推进，摩根斯坦利全球石油策略师预计，从2026年四季度起，布伦特原油将稳定在80美元/桶附近。更低的能源价格将减少中国的贸易条件损失，并通过强化全球贸易与投资周期，间接支撑出口。

*这个预测路径的脆弱性在于：政策执行能否如期加速？油价能否如预期回落？这两个变量，一个关乎国内体制效率，一个关乎地缘政治走向，都不是单纯的经济预测可以覆盖的。这也构成了这份报告最值得后续跟踪的未解问题。**

报告在外部维度提出了一个尖锐的观察：国内需求疲软意味着更多的制造业产出被导向出口市场，加速了海外市场份额的获取。这在欧盟表现得尤为明显——中欧贸易关系正从互补转向竞争。钢铁、电动车、电池、光伏和化工等领域，中国在过去两年中的份额显著上升。

欧洲的回应，报告判断将是“靶向性的”而非系统性的，工具可能包括行业关税、外国补贴调查、采购限制和反倾销措施。但报告也承认，如果欧洲逐步收窄高附加值制造业的市场准入，中国的出口引擎可能被迫转向摩擦更小但利润率更低的市场。

\subsection{[R016] 摩根斯坦利：中国电动车市场的份额洗牌远未结束，真正的分化才刚刚开始}
{\small\color{refgray}归类：未被主题摘要引用}

摩根斯坦利：中国电动车市场的份额洗牌远未结束，真正的分化才刚刚开始

5月的中国电动车市场，表面上是几家欢喜几家愁的月度排名更迭，但摩根斯坦利这份最新研报揭示的信号远比排名变化本身更值得深思：市场正在从一个“谁都能分一杯羹”的增量竞争阶段，转向一个“产品节奏、品牌势能与渠道效率必须同时在线”的存量博弈阶段。月度份额的波动，不再是随机扰动，而是企业战略执行力的压力测试结果。

这份报告的核心价值，不在于告诉你哪家品牌5月卖得好，而在于它提供了一套观察竞争格局的微观透镜：当比亚迪海洋/王朝系列份额出现年内首次环比下滑，当理想汽车在L系列改款前的阵痛期份额承压，当零跑汽车创下历史新高，当特斯拉在所有城市级市场实现份额扩张——这些碎片化信号叠加在一起，指向一个更本质的判断：中国电动车市场的份额洗牌，已经从“大厂吃肉、小厂喝汤”的粗放阶段，进入“产品周期管理能力决定季度排名”的精耕阶段。谁能在新旧产品切换的窗口期保持份额稳定甚至逆势增长，谁才真正具备穿越周期的能力。

1. 比亚迪的份额松动不是衰退信号，而是市场对“规模溢价”重新定价的开始

比亚迪海洋/王朝系列5月市场份额环比下降1.0个百分点至17.5\%，与其年初至今的平均水平持平。这个数字本身并不令人恐慌——17.5\%的市占率仍然是行业绝对龙头。但值得关注的是，这是否意味着比亚迪依靠冠军版、荣耀版等降价策略维持份额的边际效应正在递减？

摩根斯坦利在报告中并未直接给出这个判断，但数据本身提供了线索：1.0个百分点的环比降幅，在所有主流品牌中仅次于小米的0.9个百分点下滑，但小米的下滑是SU7改款后销量自然回落，而比亚迪的下滑发生在没有明显产品空窗期的背景下。这暗示，当竞争对手的产品迭代速度加快、新车型密集上市时，即便是比亚迪这样的规模王者，也需要更频繁地动用产品工具而非价格工具来维持份额。

这里的关键洞察是：规模优势正在从“自动转化为份额”变为“需要主动管理才能兑现”。比亚迪即将在6月17日发布的新款“大唐”，将是检验其能否重新激活王朝系列订单势能的关键节点。如果新款大唐无法在上市首月拉动份额回升，市场可能需要重新评估比亚迪在20万元以上价格带的竞争壁垒。

2. 零跑和特斯拉的逆势增长，揭示了两种截然不同的份额获取路径

零跑5月市场份额环比提升0.8个百分点至6.8\%，再创历史新高。这是这份报告中最值得拆解的现象之一。零跑的增长并非依靠单款爆款车型，而是凭借C系列车型在10-15万元主力价格带的持续渗透。其份额从2025年全年3.7\%跃升至2026年年初至今5.3\%，说明其“平价智能”定位正在真正转化为稳定的客户认知。

与零跑形成对照的是特斯拉。特斯拉中国5月EV市场份额环比提升2.0个百分点至5.2\%，与年初至今平均水平基本持平。更重要的是，报告特别指出特斯拉的份额增长是“all-city”的——在所有城市层级都实现了市占率提升。这意味着特斯拉的增长不是依赖某个区域市场的特殊补贴或渠道优势，而是品牌势能本身的全面扩散。

这两种路径的区别对投资者至关重要：零跑的增长依赖于精准的产品定位和性价比，但这种优势容易受到竞争对手跟进降价的侵蚀；特斯拉的增长依赖于品牌溢价和用户粘性，护城河更深但需要持续的技术领先来维持。两者都能在特定阶段实现份额扩张，但长期可持续性截然不同。

3. 理想和小鹏的份额承压，本质上是产品周期管理的“换挡阵痛”

理想汽车5月市场份额环比下降0.4个百分点至3.6\%，低于年初至今4.5\%的平均水平。报告给出的解释是L9订单增长被其他L系列车型在改款前的销售疲软所抵消。这实际上是一个经典的产品周期管理问题：当主力产品线处于生命周期末期、新车型尚未完全接棒时，企业必然要经历一段“青黄不接”的份额下滑期。

理想的应对策略是6月23日即将推

[... middle omitted ...]

存在一个未解问题：GX能否真正成为小鹏的“份额加速器”？如果GX在6月无法推动小鹏份额回升至3\%以上，市场对小鹏产品线竞争力的质疑可能会重新抬头。

4. 蔚来的份额波动背后，是“多品牌战略”能否兑现为协同效应的关键验证期

蔚来5月市场份额环比下降0.2个百分点至3.9\%，尽管批发销量环比增长28\%。报告指出部分销量可能将在下月注册，这意味着实际交付数据可能优于批发数据。但更值得关注的不是单月波动，而是蔚来即将进入的一个密集产品投放期：L60的上市，以及L80和ES9的交付。

这是蔚来首次在同一个季度内同时推进三款新车型的上市和交付。从运营角度看，这意味着蔚来的生产爬坡、供应链管理、渠道培训和售后服务能力将同时面临三重压力测试。如果蔚来能够在这个季度实现份额的环比改善，那将证明其多品牌战略（NIO主品牌+ONVO子品牌+Firefly品牌）确实具有协同效应；反之，如果三款新车型的上市反而导致运营效率下降、交付延期，市场可能需要重新评估蔚来“品牌矩阵”战略的可行性。

\subsection{[R017] 摩根斯坦利：功率半导体的涨价周期并非需求驱动，而是供给侧的理性收敛}
{\small\color{refgray}归类：半导体与功率器件：供给驱动的涨价周期与定制化浪潮}

摩根斯坦利：功率半导体的涨价周期并非需求驱动，而是供给侧的理性收敛

全球功率分立器件市场自2025年第四季度重回增长轨道。2026年4月，全球功率分立器件营收同比增长已达16\%。这一数字本身并不令人意外——半导体行业周期性复苏的叙事在过去半年已被反复讨论。真正值得关注的，是驱动本轮上行的力量结构。

摩根斯坦利最新发布的亚太半导体研报提出了一个关键判断：当前功率半导体的涨价周期是供给驱动型，而非需求驱动型。这意味着，市场若以传统周期框架来线性外推，可能会高估需求的持续性，同时低估供给侧结构性变化带来的定价权重塑。

这份报告的价值不在于确认“周期回来了”，而在于拆解“这个周期和上一个有什么不同”。对于产业决策者和投资者而言，理解这种差异，比判断涨价的幅度更重要。

全球主要功率分立器件公司的资本开支已连续两年下降。2026年虽出现约11\%的温和回升，但整体供给端扩张的力度远弱于以往任何一个复苏周期。摩根斯坦利的数据显示，中国功率分立器件市场在未来三年内新增产能有限——士兰微、华大半导体等企业正将资源转向AI电源管理芯片等更高附加值领域，而非继续扩大IGBT和MOSFET的传统产能。

这一变化的含义是深刻的。传统半导体周期中，涨价往往伴随着产能的快速释放，继而引发价格回落。但本轮周期中，供给侧的克制使得供需关系发生了结构性改变。即使需求端仅温和复苏，定价权仍可能向供应商倾斜。

报告特别指出，只要需求不发生急剧恶化，这一轮涨价趋势有望延续至2026年下半年。这个判断的核心支撑，正是供给侧的理性收敛，而非需求端的爆发。

摩根斯坦利并没有给出一个“需求全面复苏”的乐观画面。相反，报告对需求端的描述是“mixed”——混合且分化。

汽车和工业领域合计占功率分立器件终端需求的72\%，这两大领域的走势基本决定了整个市场的方向。工业自动化需求表现强劲，领先企业的营收在2026年第一季度同比增长21\%。中国新能源汽车批发销量在经历4月的零增长后，5月回升至同比增长7\%。但光伏需求持续疲弱，年初至今的装机量同比下降51\%。

这种分化意味着，本轮复苏并非由终端消费的全面回暖驱动，而是由特定领域的结构性需求支撑。工业自动化的强劲增长背后，是制造业升级和自动化替代的长期趋势；新能源汽车的增长虽有波动，但电动化渗透率仍在提升。这些都不是短期刺激政策带来的脉冲式需求，而是具有持续性的结构性力量。

对于投资者而言，理解这种分化的意义在于：不应将所有功率半导体公司等同对待。那些在工业自动化和汽车领域布局更深、客户结构更优的企业，将更受益于当前的需求结构。

报告揭示了一个值得关注的细节：在2026年2月多家中国功率分立器件公司发布的涨价通知中，除汽车用IGBT外，分销商和其他终端市场客户普遍接受了涨价。原因很简单——库存低位，且担忧2026年下半年价格进一步上涨。

这一现象背后隐藏着定价权的转移机制。当渠道库存处于低位时，分销商对涨价的容忍度显著提高。他们更担心的是“现在不买，以后更贵”，而非“现在买贵了”。这种预期自我实现的机制，使得涨价信号能够在产业链中快速传导。

摩根斯坦利在分析士兰微时特别提到，该公司通过分销渠道销售的比例高于同行，因此更容易实现涨价。这并非偶然——分销渠道的定价灵活性往往高于直供客户。那些依赖直供模式的厂商，尤其是面对大型车厂的IGBT供应商，涨价空间反而受到更多约束。

这意味着，在本轮涨价周期中，渠道结构将成为区分公司定价能力的重要变量。投资者在评估不同公司时，不应只看产品线和产能，还应关注其销售模式和渠道议价能力。

在摩根斯坦利的覆盖范围内，扬杰科技是唯一获得超配评级的功率半导体标的。报告将其目标价从91元大幅上调至136元，上调幅度达50\%，核心逻辑并非简单的周期涨价，而是公司业务结构的质变

[... middle omitted ...]

场给予的估值溢价。

5. 报告尚未完全回答的关键问题：需求的真实韧性能否支撑涨价持续性

摩根斯坦利的报告给出了清晰的供给端分析，但在需求端留下了一个值得追问的问题：当前的需求韧性，究竟能在多大程度上支撑涨价的持续性？

报告指出，工业需求强劲，但光伏需求疲弱；新能源汽车增长温和，但增速远不及 年的爆发期。如果我们将供给端的克制视为涨价的第一推动力，那么需求端的表现将决定涨价的持续时间和幅度。

这里存在一个尚未被充分验证的假设：如果工业自动化的增长动能放缓，或者新能源汽车的渗透率增速进一步下降，当前的供给-需求平衡是否会被打破？摩根斯坦利给出的判断是“如果需求不急剧恶化”，但“不急剧恶化”和“持续改善”之间，存在一个需要持续跟踪的灰色地带。

此外，报告对光伏需求的判断值得关注。光伏曾是功率半导体需求的重要增长极，但当前装机量同比下降51\%，这一趋势是否会持续？光伏产业链的库存调整何时结束？这些问题在报告中并未展开，但恰恰是判断2026年下半年需求走势的关键变量。

\subsection{[R018] JPM：外卖补贴新规的真正价值不在上限，而在让烧钱竞争变得“无法隐身”}
{\small\color{refgray}归类：中国地产与消费：量稳价跌与盈利预期重置}

JPM：外卖补贴新规的真正价值不在上限，而在让烧钱竞争变得“无法隐身”

市场正在用一个错误的框架解读中国外卖行业最重要的监管信号。当国家市场监管总局于6月17日发布《外卖平台补贴行为规范（征求意见稿）》时，主流反应是“没有硬性补贴上限，所以没有实质约束力”。这份JPM的深度研报提出了一个截然不同的判断：市场低估的不是规则的字面力度，而是信息披露作为一种执法机制的真实效力。

这份研报的核心主张清晰而克制：新规降低了美团的下行尾部风险，但尚未将监管缓和转化为正面的盈利催化剂。JPM维持美团中性评级，目标价85港元。这个判断建立在一个承重性假设之上——资金来源条款具有真实的行为约束力，能够限制以资本优势为弹药的外卖补贴攻击。如果这个前提成立，新规则降低了JPM此前给出的熊市情景（美团公允价值约51港元）的概率。

这不是一份宣告转折点的乐观报告，它更像一份情景概率的重新分配说明书。在监管信号密集、市场情绪反复的当下，这种克制恰恰是投资者最需要的东西。

1. 信息披露本身就是执法工具，市场误读了“没有上限”的含义

市场对新规的第一反应可以概括为：没有补贴上限，等于没有作用。JPM认为，这种解读忽略了规则设计的核心机制——新规不是通过设定价格上限来监管补贴，而是通过让激进的补贴行为变得更慢、更透明、更容易被挑战来发挥作用。

在外卖这个行业，最具有破坏力的竞争形式不是普通的促销，而是快速、以资本为后盾的补贴升级，其目的是在在位企业能够反应之前就改变消费者和商家的行为。七天的事前披露和资金来源报告，直接针对的就是这一机制。

研报提出了三个层面的逻辑链条。第一，信息披露降低了速度优势。外卖补贴战是反应式的——挑战者发起促销，在位者回应，市场测试双方的极限。七天公告窗口削弱了突袭的价值，破坏了针锋相对的动态循环。第二，信息披露创造了证据。资金来源披露和竞争影响自评估迫使平台写下监管者日后评估其补贴是否依赖资本优势或不正当竞争所需的事实。第三，信息披露在处罚到来之前就提高了攻击成本。一个必须提前披露的补贴活动，更容易被竞争对手、商户、骑手、媒体和监管者监控。这种声誉和执行风险，会降低平台发起极端补贴活动的意愿。

这个分析框架的实质是：规则不需要禁止每一项补贴来改变平台行为，它只需要让最激进的补贴活动变得更难发起、更难辩护、更难重复。对于投资者而言，这意味着需要调整的不是盈利模型中的数字，而是对竞争格局演变概率的判断。

2. 第三条是承重墙：补贴资金来源的界定将决定竞争格局的重新洗牌

研报中反复强调的“承重性条款”是第三条，即补贴应来自自身经营利润而非资本优势。投资结论的最终走向，取决于“自身经营利润”如何被解释。

JPM给出了三种可能的解释情景。如果最终解释为“分部级或接近分部级的利润”，这将强烈利好美团，因为挑战者无法轻易用集团资本来补贴外卖亏损。如果最终解释为“集团级利润”，则不对称性被削弱，阿里巴巴可以主张集团整体盈利足以支持持续投资。如果最终措辞模糊，这个问题将继续存在，执法行为将比文本本身更为重要。

研报的基准解读是，“资本优势”这一措辞意在限制交叉补贴和以资产负债表为后盾的份额攻击。这正是该规则对美团重要的原因。但这也是研报观点面临的最大风险——如果最终文本明确允许集团级利润满足要求，美团的战略优势将大幅缩水。

这一分析的价值在于，它为投资者提供了一个清晰的观察框架：最终规则的措辞选择，将成为判断竞争格局演变方向的关键路标。在7月17日征求意见截止之前，市场将围绕这一条款展开密集的解读和博弈。

如果新规按预期发挥作用，竞争的维度将发生根本性转变。JPM从五个维度对比了规则前后的变化：活动速度从“闪电促销和快速反应”变为“七天事前披露”；资金来源从“集团资本可支持外卖亏损”变为“自身经营利润变得

[... middle omitted ...]

位就会改善。

但这并不意味着外卖竞争消失了。竞争仍然可以转移到履约质量、商户选择、会员权益、物流和高客单价场景。然而，恰恰是这些领域，规模、密度和运营能力比单纯烧钱更为重要。

研报对利润池的分析值得特别注意。JPM认为，新规降低了行业利润池变为负值的概率，但并未将行业恢复到2024年式的利润天花板。转向服务品质、物流和商户支持的竞争仍将消耗资金。更为重要的变化是，持续的亏损资助型补贴竞争将变得更加难以辩护。

对于美团而言，这是一个混合但净正面的结果。补贴线面临的尾部风险应该减小，但配送成本不太可能干净地回到2025年之前的水平。JPM的模型已经假设骑手单均成本从约5.0元向5.2-5.3元移动，新规强化了这一方向。

对于京东而言，情况更为复杂。研报指出，京东的外卖业务是最依赖资本补贴的。如果第三条被严格执行，京东外卖位置的可持续性将更难被支撑。京东进入外卖领域的持久影响，可能更多体现在提高了行业成本地板，尤其是骑手工资和履约标准方面，而非永久性的市场份额获取。

\subsection{[R019] 摩根斯坦利：2026陆家嘴论坛的真正信号不是开放，而是金融业增长逻辑的重新定价}
{\small\color{refgray}归类：外汇与资本流动：人民币定价机制微调与跨境框架重构}

摩根斯坦利：2026陆家嘴论坛的真正信号不是开放，而是金融业增长逻辑的重新定价

市场对2026年陆家嘴论坛的关注，大多集中在“进一步开放”这个关键词上。新的QDII额度、简化ODI流程、跨境贸易外汇结算试点——这些政策细节确实被密集释放。但摩根斯坦利这份研报试图传达的判断，远比“开放”二字复杂。

这份报告的核心主张是：陆家嘴论坛传递的，不是简单的政策宽松信号，而是一套关于中国金融业增长质量的重新定义框架。开放是手段，但风险控制和高质量发展才是目的。两者并非对立，而是构成了一个更精细的政策组合——这意味着金融机构的估值逻辑，需要从“规模增长”切换到“质量定价”。

为什么这个判断重要？因为过去几个月，市场对跨境资本流动收紧的担忧在累积。香港市场和中概股的波动，部分源于这种不确定感。摩根斯坦利在5月底的两份报告中已经试图缓解这种担忧，而本次论坛的政策信号，则是对市场预期的直接回应。但报告同时明确指出，开放不等于放松监管——银行业的理性信贷增长、保险业的费用管控、资本市场的有序IPO机制，这些才是论坛真正的结构性信号。

1. 开放政策的信号意义大于实质：它缓解的不是资本流出，而是政策不确定性溢价

摩根斯坦利报告明确将论坛的开放信号与5月以来的市场担忧联系起来。PBOC和SAFE在论坛上宣布的几项具体政策——新的QDII额度、简化ODI、外债外汇管理优化、上海跨境贸易外汇结算试点——在报告中被描述为“支持人民币全球化和拓宽而非限制跨境资本流动”。

但仔细看，这些政策本身并非颠覆性的。QDII额度增加是渐进式操作，上海试点是局部推进。真正重要的是信号本身：政策制定者有意向市场表明，金融开放的长期方向没有改变。这一点，对于过去几个月因各种传闻而承压的跨境资本流动预期，是一个明确的纠偏。

这意味着什么？对于投资中国金融资产的全球投资者而言，真正被降低的不是资本流动的实际摩擦成本，而是政策不确定性溢价。后者对资产定价的影响，往往比前者更大。当市场不再需要为“会不会突然收紧”而支付风险溢价时，估值中枢的重估就有了基础。

但报告同时提醒，这种信号效应能否持续，取决于后续政策的执行节奏和力度。这里有一个尚未完全回答的问题：开放政策是否会因为宏观审慎考虑而出现阶段性调整？报告没有给出确定性答案，而是通过强调“风险控制”和“高质量增长”的并行表述，暗示政策框架已经内置了这种平衡。

2. 银行业的核心叙事已从“增速”转向“质量”：理性信贷增长成为新常态意味着什么

报告中最具结构性的判断，来自PBOC关于信贷增长的表述：更慢但更高质量的贷款增长可能成为新常态。摩根斯坦利将其描述为“与我们对理性信贷增长的看法一致”。

这句话的潜台词值得深挖。过去几年，中国银行业的估值逻辑在很大程度上依赖于规模扩张——贷款增速、资产规模增长、市场份额变化。当“理性信贷增长”成为政策导向时，银行之间的竞争将从“谁放贷更快”转向“谁放贷更好”。这意味着：

第一，资产质量而非规模增速，将成为银行股定价的核心变量。那些在信贷审批和风险管理上有优势的银行，其估值溢价将扩大。第二，息差收窄的压力可能会被更严格的信贷纪律部分对冲——如果银行不再为争夺市场份额而压低利率，净息差的下行斜率可能趋缓。第三，系统性风险的降低，本身就是估值提升的因素——当政策明确为信贷增长设置“质量门槛”时，市场对银行体系尾部风险的定价会下降。

报告还提到PBOC正在设立针对非银行金融机构的流动性支持工具，以防止系统性风险。这进一步强化了“质量优先”的政策导向：增长可以慢，但不能出风险。

3. 保险业的监管信号直指费用竞争：这对财产险的综合成本率是实质性利好

报告对保险业的判断相对简洁但明确：NAFR继续强调纠正无序竞争和严格执行费用规则，这将有利于财产险公司的综合成本率

[... middle omitted ...]

力，同样是行业的关键变量。监管信号对寿险的影响，可能需要从更长的利率和资产配置周期来理解。这是我们后续可以深入探讨的方向。

4. 资本市场改革更注重“效率”而非“数量”：灵活的IPO机制和产品创新是结构性变化

CSRC在论坛上宣布的几项改革——修改上市规则、加速推出储架注册机制、推出主动管理ETF和商业地产REITs试点——在摩根斯坦利的解读中，核心是“更灵活但也更有纪律的IPO机制”。

储架注册机制的意义在于，它允许发行人根据市场条件灵活选择发行窗口，而不是在审批通过后必须在限定期限内完成发行。这降低了发行人的时间成本，也减少了市场因集中发行而承受的供给压力。对于券商而言，这意味着投行业务的模式可能从“项目导向”转向“持续服务导向”，承销能力的价值将重新定义。

主动管理ETF和商业地产REITs试点，则指向产品端的多元化。这些创新产品的推出，有助于拓宽资本市场的深度和广度，为投资者提供更多风险收益特征的选择。对于资产管理公司而言，产品创新能力将成为竞争壁垒。

\subsection{[R020] NOM：人民币中间价模型正在发出一个被市场忽视的定价信号}
{\small\color{refgray}归类：外汇与资本流动：人民币定价机制微调与跨境框架重构}

人民币汇率中间价的每日设定，长期以来被市场视为一个政策信号灯。大多数投资者关注的是中间价本身的方向——是升值还是贬值，以及它与前一交易日收盘价的偏离幅度。但NOM最新发布的USD/CNY定盘价模型报告揭示了一个更精细的维度：模型内部各货币的贡献权重和逆周期因子的调整幅度，正在传递一个关于央行汇率管理策略变化的信号。这个信号比中间价的方向更重要，因为它指向的是政策框架的微调，而非一次性干预。

这份报告的核心判断是：市场对人民币汇率定价的理解，仍停留在“央行是否容忍贬值”的二元框架中，而忽略了中间价形成机制内部结构的持续变化。真正重要的不是中间价定在哪个数字，而是这个数字如何被计算出来。

1. 模型预测与官方定盘价的偏离，正在暴露政策意图的微妙变化

NOM模型给出的最新预测值是6.7745，较前一交易日下降385个基点。如果考虑逆周期因子调整，预测值为6.7965，较前一交易日下降165个基点。这两个数字本身并不惊人——人民币中间价在6.7-6.8区间波动已是常态。真正值得关注的是模型误差的趋势。

报告中的模型误差图显示，自2025年中的低点以来，模型误差（即模型预测与官方定盘价之间的差值）已经从接近零的水平扩大至约600个基点的正值区间。这意味着，官方定盘价系统性地高于模型预测值——换言之，央行在中间价设定中持续加入了比模型预期更强的“抗贬值”力量。

这个趋势不是偶然的。它表明，央行并非简单地维持中间价稳定，而是在主动压缩贬值空间。当模型预测中间价应该更低（即人民币更强）时，官方定盘价却更高（即人民币更弱），这恰恰说明央行的操作方向与市场预期相反——它在抑制升值，而非阻止贬值。这一点对很多持有“央行在放任贬值”观点的投资者来说，是一个重要的认知修正。

2. 欧元和日元成为中间价变动的主要贡献者，折射出人民币定价的全球维度正在重新定义

报告中最具洞察力的部分，是图1中展示的隔夜变动贡献权重。在影响模型预测变动的四大货币中，欧元贡献了35.5个基点，日元贡献了18.8个基点，韩元和英镑紧随其后。这个权重分布并非随机。

欧元和日元在CFETS人民币汇率指数中的权重分别为约18\%和15\%，但它们在中间价模型中的贡献占比远高于这一比例。这意味着，央行在设定中间价时，对欧元和日元隔夜变动的反应系数正在上升。这背后有两个可能的解释：一是央行在主动增加人民币对一篮子货币的锚定，而非仅盯住美元；二是央行在利用欧元和日元的波动来对冲美元单向变动的影响。

对于企业财务总监和跨境投资者来说，这个变化意味着：过去只需要关注美元/人民币汇率，现在需要同时关注欧元/美元和美元/日元的联动关系。人民币汇率的决定因素正在从“中美双边关系”向“全球货币体系”扩展。

3. 逆周期因子正在从“防御工具”转变为“主动管理工具”，其使用频率和幅度都在变化

逆周期因子是央行用来平滑中间价波动、防止单边预期的工具。过去几年，市场习惯了将其视为一种“应急手段”——只有在汇率出现剧烈波动时才会被激活。但这份报告的数据暗示，逆周期因子的使用已经变得更加常规化和精细化。

模型预测与考虑逆周期因子后的预测值之间的差值，从385个基点缩小到165个基点，意味着逆周期因子实际上抵消了约220个基点的模型变动。这个幅度不大不小，既不是完全放任市场，也不是强力干预。它传递的信号是：央行希望中间价保持稳定，但不愿意让市场认为它在进行大规模干预。

这种“润物细无声”的操作风格，可能是未来人民币汇率管理的新常态。对于交易员而言，这意味着不能再用“是否突破某个点位”来判断央行态度，而需要更细致地观察中间价形成机制中的参数变化。

4. 报告尚未完全回答的关键问题：模型误差的持续扩大是否意味着央行容忍了更大的定价偏离？

模型误差从接近零扩大到600个基点，是这

[... middle omitted ...]

人会晤）前预留政策空间。

报告中的事件日历表列出了2026年下半年多个关键时点：7月底的政治局会议、10月的国庆黄金周、11月中国主办的APEC会议、12月的中央经济工作会议，以及年底可能的中美领导人会晤。这些事件构成了一个密集的政策窗口期。央行可能在利用当前相对平稳的市场环境，逐步调整中间价定价机制，为未来可能出现的波动做准备。

5. 对投资者的观察框架：从“中间价方向”转向“中间价形成机制”

基于这份报告，投资者可以建立一个更立体的观察框架。不再只是问“中间价是升是贬”，而是问四个问题：

第一，模型预测值与官方定盘价的偏离幅度在扩大还是缩小？这反映了央行的干预意愿。

第二，逆周期因子的调整方向是什么？是放大还是缩小模型变动？这揭示了央行的主动管理倾向。

第三，各货币的贡献权重是否在变化？欧元和日元的权重上升是否持续？这指向人民币定价的全球化程度。

第四，模型误差的趋势是否与政策事件日历相关？误差扩大是否发生在重要会议前？这可以帮助判断央行的政策节奏。

\subsection{[R021] NOM：亚洲出口的超级周期并非只有AI，一个被低估的转折点正在浮现}
{\small\color{refgray}归类：外汇与资本流动：人民币定价机制微调与跨境框架重构}

NOM：亚洲出口的超级周期并非只有AI，一个被低估的转折点正在浮现

市场对亚洲出口的讨论，大多集中在AI驱动的半导体和电子产业链上。但NOM最新发布的一份图表报告，揭示了一个更重要的信号：一个领先指标已经飙升至16年新高，而驱动因素正在从单一的技术周期，向一个更宽广、更可持续的复苏切换。

这份报告的核心判断是：亚洲出口的超级增长周期远未结束，但市场可能低估了其结构性扩围的潜力。真正值得关注的，不是AI需求还能多强，而是当供应链扰动消退、全球资本品需求回升时，哪些国家、哪些行业能承接这轮“扩围”的增量。

1. 一个历史级信号：NELI指数升至16年高位，预示的不只是短期繁荣

NOM亚洲出口领先指数（NELI）在7月攀升至121，这是自2006年以来的最高水平。该指数由九个分项组成，领先实际出口数据约三个月，且历史上准确预测了亚洲出口的每一次重大拐点。

这个数字意味着什么？它暗示亚洲（除日本外）的出口增长动能将在三季度延续，甚至可能进一步加速。更重要的是，NELI的上升正在变得更加“广泛”——不仅科技相关指标强劲，整体制造业和中国需求也出现了改善迹象。

这并非一个孤立的短期脉冲。NELI的构成决定了它能够过滤掉基数效应的干扰，反映的是出口增长的潜在动量。当这样一个领先指标创下16年纪录，其对资产定价的启示是：以出口为导向的经济体（如韩国、台湾、越南、马来西亚）的宏观基本面，可能在未来几个季度持续超预期。

报告明确指出，AI技术周期一直是亚洲出口的主要引擎。这一点市场已有充分定价。但NOM报告的关键增量，在于指出了另一个正在成形的驱动因素：霍尔木兹海峡的重新开放。

这一点值得仔细拆解。霍尔木兹海峡的通行恢复，将缓解全球供应链压力。虽然正常化过程可能需要数月，但方向是明确的。这意味着什么？它意味着全球不确定性下降，企业资本开支意愿回升。而资本品需求——即机器、设备、工业零部件——是亚洲出口中一个体量巨大、但对供应链风险极度敏感的部分。

当供应链扰动缓解，企业从“保供应”转向“扩产能”，资本品需求将成为亚洲出口增长的新引擎。这轮扩围将超越科技行业，惠及工业、机械、化工等更广泛的领域。报告认为，这将支持亚洲出口周期的进一步拓宽。

当前市场对亚洲出口的叙事，高度集中于AI芯片、服务器和相关零组件的出货量。这没有问题，但可能过于狭窄。NOM报告提出的逻辑是：供应链压力的缓解，将释放被压抑的资本品需求，而这部分需求对亚洲出口的拉动，可能比市场预期的更为显著。

从历史经验看，亚洲出口的超级周期往往由科技需求启动，但随后会扩散到更广泛的制造业。 年的全球贸易复苏就是一个典型例子。当时，半导体周期率先复苏，随后资本品和消费品出口同步跟进，形成了全面的繁荣。

当前的信号显示，类似的结构性扩散可能正在发生。NELI的“广度”提升，正是这一过程的早期证据。投资者需要关注的，不是AI需求是否见顶，而是非科技行业的出口动能能否持续改善。如果答案是肯定的，那么当前对亚洲出口的定价可能仍然偏保守。

NOM报告提到“中国需求出现改善迹象”，但并未对此展开详细论证。这恰恰是当前框架中最大的不确定点。

中国作为亚洲最大的最终需求市场和中间品贸易枢纽，其内需的强弱直接影响整个区域的出口链条。NELI中包含的中国相关指标，是否反映了政策刺激的短期效应，还是结构性复苏的开端？这一点报告的图表数据可以给出方向，但无法提供确定性答案。

此外，报告没有讨论的是，如果中国需求改善主要来自库存回补而非终端消费，那么其对亚洲出口的拉动可能会在几个季度后衰减。这个假设的验证，将直接决定这轮出口周期的斜率。

基于这份报告，投资者可以建立一个更系统的观察框架。核心是跟踪两个维度的变化：

第一是“广度”。NELI的九个分项中，有多少个在同步改善？

[... middle omitted ...]

决的问题：霍尔木兹海峡正常化的具体时间表如何？中国需求的改善能否持续？资本品需求的回升，是否会因利率环境而受到抑制？这些问题，恰恰是理解这轮周期最终高度的关键。

欢迎来星球微信群里继续讨论。我们将持续追踪NELI指数的月度变化，以及资本品需求、中国进口等关键指标的后续数据。在群里，我们可以结合原始报告的完整图表，对每一个分项的边际变化做更细致的拆解，并讨论这些信号对不同市场、不同资产的定价含义。

Personal reading notes and learning share only. Not investment advice.

最近翻到一份NOM的研报，发现了一个挺有意思的信号：亚洲出口的景气度，可能比我们想象的还要强。

1️⃣ 一个前瞻指标创了16年新高 NOM有个领先指标叫NELI，专门用来预测未来3个月的亚洲出口走势。7月这个读数飙到了121，是16年来的最高点。 这个指标由9个细分项构成，过去几次大的拐点都提前捕捉到了，比看官方出口数据更灵敏。

\subsection{[R022] JPM：中国消费板块的真正底部不在估值，而在盈利预期的重置}
{\small\color{refgray}归类：中国地产与消费：量稳价跌与盈利预期重置}

中国消费板块正在经历一场罕见的估值与情绪双重压缩。JPM最新发布的研报以“Cheap but unloved”为题，直指当前市场最核心的矛盾：板块估值已跌至十年均值下方1.5个标准差，但80\%覆盖公司的盈利预测仍在被下调。这份报告的真正价值不在于告诉你“便宜”，而在于揭示了“便宜”背后尚未被定价的结构性分化——下行周期中，真正决定回报的不是行业beta，而是企业运营质量。

如果你认为消费板块的困境只是需求疲软，那你很可能错过了真正重要的信号。JPM的分析框架显示，市场对2026年盈利预期的下修远未结束，更关键的是，成本端压力将在2026年下半年集中释放，而这一点目前几乎未被充分定价。本文基于这份研报的核心逻辑，提炼出五个对投资者真正重要的判断。

当前市场对消费板块的定价，隐含了一个假设：盈利下修已经大部分完成。但JPM的数据挑战了这一假设。截至报告发布，必需消费和可选消费板块2026年一致预期EPS分别较年初下调了9\%和7\%。从覆盖公司来看，80\%的标的年初至今遭遇盈利下调，白酒、美妆等子行业是重灾区。

然而，更值得关注的不是已经发生的下修，而是尚未出现的压力。报告明确指出，“成本压力的充分定价可能要到2026年底才会显现”。这一判断基于一个关键情景分析：如果伊朗冲突结束，包装材料和工业原料（如PE/PET）价格将较2025年均价上涨15-40\%。结合JPM中国宏观团队对PPI和CPI的预测——2026年2-4季度PPI分别为3.5\%、3.9\%、3.5\%，CPI仅为1.3\%、1.4\%、1.1\%——CPI-PPI剪刀差将扩大至约2.9个百分点。

这意味着什么？成本端压力正在积聚，但终端需求疲软限制了企业转嫁成本的能力。盈利下修的故事远未结束，2026年下半年才是真正的压力测试窗口。对于依赖成本传导的饮料、包装食品等子行业，这一风险尤其值得警惕。

如果成本压力真的来临，哪些子行业具备防御能力，哪些将首当其冲？JPM的框架给出了清晰的排序。

白酒和啤酒板块相对更具韧性。原因不在于它们不受成本影响，而在于这两个子行业的竞争格局更为稳定，头部企业拥有较强的定价权。报告特别点名茅台、主要啤酒厂商和海天味业，认为它们“更有能力转嫁成本”。茅台的特殊性在于其奢侈品属性，提价几乎不受需求约束；啤酒行业的集中度较高，龙头企业对渠道和终端的控制力更强；海天作为调味品龙头，品牌壁垒和渠道粘性构成了定价权的护城河。

饮料板块则是最脆弱的环节。饮料行业竞争激烈，产品同质化程度高，价格战风险始终存在。在成本上升而需求疲软的双重挤压下，企业的毛利率将面临显著压力。JPM明确建议在2026年第三季度避开饮料板块，尤其是那些被市场视为“高股息”的交易标的——这些公司的分红能力本身可能受到盈利下滑的侵蚀。

这一判断的直接含义是：成本冲击不是均匀分布的，投资者需要从“买行业”转向“挑公司”，真正具备定价权的企业才能穿越周期。

3. Alpha的来源正在从行业beta转向运营质量，运动服饰、IP玩具和白色家电是结构性赢家

JPM的报告提出了一个核心观点，值得每一位投资者反复咀嚼：“Alpha来源于运营质量，而非行业beta。需求端可能继续成为头条噪音，但真正的故事在于下面的分化。”

这一判断基于一个系统性的评分框架，涵盖基本面（需求、供给、定价、利润率）、回报（股息）和估值三个维度。在该框架下，运动服饰、IP玩具和白色家电被识别为未来3-5年最具吸引力的子行业。报告明确提及安踏、泡泡玛特和美的作为代表标的。

这三者的共同特征是什么？不是它们所处的行业增速最快，而是它们的竞争格局和商业模式赋予了它们超越行业周期的能力。运动服饰受益于健康消费的结构性趋势和国产品牌的替代逻辑；IP玩具对应的情绪消费具有抗周期属性，且头部企业的IP运营

[... middle omitted ...]

。JPM的数据显示，必需消费板块的远期PE约为16倍，低于十年均值约1.5个标准差；可选消费板块的远期PE约为13倍，同样处于历史底部区域。从子行业来看，饮料、啤酒、调味品、餐饮、运动服饰、美妆和珠宝均处于各自十年估值区间的底部十分位。

但估值低不等于立即买入。报告中的一张关键图表揭示了另一个事实：盈利修正和股价表现之间存在显著背离。以运动服饰为例，安踏和和李宁的盈利修正年初至今分别为-4\%和+8\%，但股价表现分别为-10\%和-11\%。这意味着市场对盈利前景的悲观预期可能已经过度反映在股价中。

白色家电是另一个值得关注的例外。美的的盈利修正年初至今仅-3\%，股价表现基本持平，显示出较强的抗跌性。在消费板块整体承压的背景下，这样的相对表现本身就说明了市场对其基本面的认可。

综合来看，估值底提供了安全边际，但真正的买入信号需要等待盈利预期的触底。运动服饰和白色家电可能是最早出现拐点的子行业。

5. 报告尚未完全回答的关键问题：成本冲击的传导路径和竞争格局的变化速度

\subsection{[R023] 摩根斯坦利：中国钢铁市场的真实复苏信号，被周度数据的噪音掩盖了}
{\small\color{refgray}归类：能源与大宗商品：供给约束下的时间错配与地缘风险折价}

摩根斯坦利：中国钢铁市场的真实复苏信号，被周度数据的噪音掩盖了

过去几周，围绕中国钢铁和铁矿石市场的讨论，大多集中在需求何时触底、库存压力能否缓解、以及减产政策是否会加码。这些问题的核心，其实指向同一个判断：市场参与者是否正在用错误的频率观察一个正在发生结构性分化的行业。

摩根斯坦利最新发布的中国钢铁与铁矿石周度更新报告，提供了截至6月中旬的一组高频数据。单独看每一项，似乎只是常规的环比波动——长材表观消费周环比上升5.1\%，扁平材上升2.0\%；贸易商库存微降0.6\%，钢厂库存微增0.7\%；电炉开工率提高2个百分点。但如果把这些事实放在一起，一个更重要的信号浮现出来：市场真正低估的不是需求总量，而是供给端正在发生的、由电炉产能利用率回升所驱动的再定价。

这份报告的价值，不在于它预测了价格方向，而在于它提供了一套可以持续追踪的、从“库存-产量-消费-发运”四维交叉验证的观察框架。真正值得关注的，不是某一周的数据好坏，而是这些数据叠加在一起，是否指向一个可持续的、由供给结构变化驱动的平衡点。

报告显示，长材表观消费周环比增长5.1\%，而扁平材仅增长2.0\%。两者差距并非偶然。

长材主要用于建筑和基建领域，其消费回暖往往与项目开工节奏、资金到位情况直接相关。5.1\%的周环比增速，如果放在一个更长的时间窗口里观察，可能意味着前期积压的施工需求正在释放。而扁平材更多对应制造业和出口，2.0\%的增速相对温和，与当前制造业PMI的修复节奏基本吻合。

但这里有一个值得追问的问题：长材消费的回升，是补库行为还是终端真实消耗？报告没有直接回答，但它提供的另一个数据提供了线索——贸易商库存下降的同时，钢厂库存微增。这意味着，下游拿货的意愿在增强，但钢厂并没有同步去库。这种“中间环节去库、上游累库”的结构，通常暗示需求回升的可持续性还需要更多数据验证。

对于关注钢铁板块的投资者而言，长材与扁平材的消费分化，是一个需要持续跟踪的领先指标。如果长材消费能在未来2-3周维持环比正增长，那么建筑端需求的底部可能已经形成。

报告中最引人注目的数据之一，是电炉开工率周环比上升2.0个百分点至64.4\%，同比更是高出7.7个百分点。这一变化的意义，远大于高炉开工率仅上升0.1个百分点的平淡表现。

电炉产能的灵活性远高于高炉。当废钢价格相对铁水成本具有竞争力、且终端需求有边际改善时，电炉可以快速复产。2个百分点的周度跳升，说明当前废钢供应和比价关系已经支持电炉钢厂增加产量。而同比7.7个百分点的巨大差距，则意味着电炉在整个钢铁供给结构中的权重正在系统性提升。

这带来的直接含义是：钢铁供给的边际弹性正在变大。过去，市场习惯于用高炉开工率和247家钢厂产能利用率（当前为90.3\%）来判断供给压力。但如果电炉产能持续释放，那么即便高炉开工率维持稳定，总供给也可能超出预期。对于铁矿石需求而言，这意味着每多一吨电炉钢产量，就对应着大约1.6吨铁矿石需求的替代。这正是为什么报告数据显示，虽然钢厂铁矿石库存周环比下降1.2\%，但港口库存反而微增0.7\%——电炉对铁矿石的替代效应正在显现。

3. 铁矿石发运的“澳洲降、巴西升”背后，是更复杂的全球供应博弈

报告披露了6月8日至14日当周的铁矿石发运数据：澳大利亚发运量周环比减少164万吨，而巴西发运量增加224万吨，两者合计净增59万吨。

单独看一周的数据，这种波动可能在统计误差范围内。但如果将视角拉长，澳大利亚与巴西发运的此消彼长，反映的是不同矿商在成本曲线上的位置差异。澳大利亚主要矿商（如力拓、必和必拓）的发运节奏往往受港口检修和天气因素影响较大，而巴西淡水河谷的发运弹性更多取决于其矿山复产进度和物流瓶颈的改善。

对于中国市场而言，真正重要的不是某一周的发运总量，而是到港量与钢厂补库节

[... middle omitted ...]

的上行空间将被压缩。

4. 报告尚未完全回答的关键问题：库存结构的变化是否预示价格拐点

研报提供了清晰的库存分层数据：贸易商库存周环比下降0.6\%，钢厂库存上升0.7\%，港口铁矿石库存上升0.7\%。这种“贸易商去库、上游累库”的组合，在历史上往往出现在价格下跌周期的末端——下游开始拿货，但上游尚未完全消化前期积累的库存压力。

但这一次是否相同？报告没有给出结论，而这正是需要进一步讨论的地方。有三个变量可能改变传统解读框架：

第一，电炉产能的快速释放，意味着“钢厂库存”的内涵正在发生变化。如果新增库存主要来自电炉，那么其成本结构和高炉不同，对价格的敏感度也不一样。

第二，贸易商库存的下降，是主动去库还是被动去库？如果是主动去库（即贸易商预期价格将继续下跌），那么这轮去库可能尚未结束；如果是被动去库（即下游需求确实在改善），那么拐点可能已经临近。

第三，铁矿石港口库存的增加，是否包含了巴西矿石的到港集中？如果是，那么随着澳洲发运回升，港口库存可能很快转为下降。

\subsection{[R024] BARC：光学市场真正的拐点不在于AI需求本身，而在于供给结构正在经历不可逆的重塑}
{\small\color{refgray}归类：AI与算力：从训练到推理的架构迁移与定制芯片崛起}

BARC：光学市场真正的拐点不在于AI需求本身，而在于供给结构正在经历不可逆的重塑

市场对AI拉动光学器件需求的叙事已经足够熟悉。过去两年，每一次关于数据中心互联、超算集群扩容的消息，都能在资本市场引发一轮对光学板块的关注。但这份BARC最新发布的光学模型更新报告，提供了一个更为细腻、也更值得产业决策者仔细审视的判断框架。

报告的核心结论并不令人意外——BARC预计2026至2027年，全球光学网络市场将维持两位数的增长。真正有价值的洞察，藏在它对市场结构的拆解之中：增长的动力来源正在发生系统性切换，而不同细分市场的竞争壁垒和定价权正在经历根本性分化。理解这种分化，比单纯判断“AI需求是否真实”更具投资意义。

这份报告更新了BARC对光学市场的详细模型，涵盖了长途、海底、城域以及可插拔模块（ZR/ZR+）的细分，并首次在客户垂直领域（电信、云、企业）做了明确的收入拆分。对于关注Ciena、思科等核心标的的投资者而言，这份报告提供了一个从总量判断下沉到结构性判断的分析框架。

1. 两位数的增长预期背后，是电信与云的双重驱动力正在形成合力

BARC将2026年光学市场增长预期从此前的8\%上调至10\%，2027年从13\%微调至12\%。表面上看，这只是几个百分点的调整。但更值得关注的是，驱动这一增长的底层逻辑正在从“单一引擎”转向“双引擎”。

过去几年，光学市场的增长呈现出典型的“跷跷板”效应。2022年，云计算厂商的大规模网络建设推动了市场增长。随后，电信运营商进入库存消化周期，叠加宏观压力，导致2023至2024年市场连续下滑。BARC的数据显示，2023年光学市场同比下降7\%，2024年进一步下滑12\%。这种“云强则市场强，云弱则市场弱”的特征，让整个行业高度依赖单一客户群体的资本开支节奏。

但BARC的模型显示，这种格局正在改变。2025年起，电信侧的恢复正在形成与云侧增长并行的第二股力量。报告中提到，多光纤网络（MOFN）机会正在为Ciena的电信业务贡献约10\%至15\%的收入，而BEAD（宽带公平接入与部署）资金正在逐步进入市场。电信运营商的网络重构需求——尤其是长途网络的重新架构——正在成为新的增长支点。

这意味着什么？如果电信侧的恢复能够持续，光学市场的增长将不再完全依赖于几家超大规模云厂商的资本开支节奏。这种需求来源的多元化，本身就是一种风险对冲。

2. 城域市场正在经历一场“可插拔革命”，但不要高估它对整体市场的冲击

BARC报告中一个引人注目的数据点是，可插拔模块（ZR/ZR+）在城域WDM市场中的渗透率，预计将从2024年的14\%跃升至2027年的51\%。从增长速度看，这确实是一场“革命”。ZR模块市场预计在2026年增长72\%，ZR+增长33\%。

然而，BARC同时给出了一个容易被忽略的定量判断：可插拔模块在整个光学市场中，2026至2027年期间的平均占比仅为高十几个百分点。换言之，尽管它在城域市场内部增长迅猛，但从全市场占比看，它仍然是一个相对较小的细分领域。城域市场剔除可插拔模块后，同期预计将呈现高个位数的下降。

这个对比揭示了一个关键洞察：可插拔模块的增长，更多是在替代传统城域设备，而非创造全新的增量市场。对于投资者而言，这意味着“可插拔概念”本身并不构成一个独立的投资主题。真正的问题在于，谁能在这一替代过程中保住或扩大自己的份额。BARC倾向于认为，Ciena凭借其在ZR+领域的技术领先（率先推出800G ZR+产品），有望在可插拔市场中维持约15\%的份额。

3. 长途市场才是真正的“AI受益者”，其增长逻辑比城域更硬核

如果说城域市场的故事是关于替代，那么长途市场的故事则是关于扩容。BARC预计，长途WDM市场在2025年增长约9\%，2026至2027年

[... middle omitted ...]

报中明确提到，预计2027财年收入将出现“有意义的增长”，部分原因正是超大规模网络（hyper-rail）的部署。更重要的是，管理层强调，与超大规模客户的紧密关系使其能够看到产品正在被实际部署，而不是堆积在仓库中。这与疫情后供应链危机时期的库存堆积形成了本质区别。

4. 电信仍是市场主导者，但“云占比”的快速攀升正在重塑竞争格局

BARC的模型揭示了一个有趣的结构性变化：尽管云计算增长迅猛，但电信运营商在光学市场中的占比在2025至2027年期间仍将超过50\%。然而，这一数字比几年前的中高70\%已经有了显著下降。到2027年，云计算在光学市场中的占比预计将从2024年的约20\%上升至约32\%。

这一变化对竞争格局的影响是深远的。传统上，光学设备市场是一个以电信运营商为核心客户、以技术认证和长期合作关系为护城河的行业。但云计算厂商的采购逻辑不同——它们更看重技术迭代速度、大规模交付能力和生态系统兼容性。这意味着，能够同时服务好两类客户的供应商将拥有更大的战略纵深。

\subsection{[R025] GS：全球IT服务订单的转折点不是AI，而是地缘政治}
{\small\color{refgray}归类：半导体与功率器件：供给驱动的涨价周期与定制化浪潮}

全球IT服务行业正在经历一个被市场误读的转折点。当投资者将目光聚焦于AI对传统系统集成商的颠覆性冲击时，GS最新发布的Accenture财报解读却揭示了一个更微妙、也更值得警惕的信号：订单下滑的真正推手并非技术革命，而是地缘政治摩擦。

这份报告的核心判断值得每一位关注日本IT服务板块的投资者认真思考：全球IT服务需求的结构性分化正在加速，AI相关业务在扩张，但地缘政治风险正在成为悬在系统集成商头上的达摩克利斯之剑。市场可能低估了后者对行业基本面的影响周期。

Accenture 3Q8/26（3-5月）财报显示，销售额基本符合预期，但新订单出现了四个季度以来的首次负增长，同比下降3\%。更值得关注的是，管理服务订单同比骤降15\%，而咨询订单却加速增长13\%。这种分化背后，是中东局势导致部分客户决策延迟，而非AI对传统业务的替代。

对日本IT服务板块而言，这份报告的启示在于：本土需求相对稳健，但全球化敞口大的企业将面临更大的不确定性。与此同时，AI相关业务的扩张趋势在日本也开始显现。

Accenture 3Q新订单的负增长并非均匀分布。管理服务订单从2Q的+3\%骤降至-15\%，而咨询订单从+8\%加速至+13\%。这种结构性分化是理解当前行业周期的关键。

管理服务通常代表长期外包合同，其订单下滑往往意味着客户对未来业务不确定性的担忧。而咨询订单的加速增长，则反映出企业在战略层面仍然愿意投资于数字化转型和AI落地。两者之间的背离，恰恰说明客户正在将预算从“维持性支出”转向“变革性支出”。

但这里有一个容易被忽略的细节：中东局势对订单的冲击主要集中在咨询业务，约4亿美元的订单延迟。也就是说，即便是在咨询领域，增长也并非完全由内生需求驱动，而是受到了地缘政治因素的扰动。报告明确指出，这一影响在4Q可能进一步扩大。

对于投资者而言，这意味着不能简单地将咨询订单的增长解读为行业全面复苏的信号。客户决策的延迟可能正在制造一个“虚假繁荣”的窗口期。

GS报告中最具洞察力的部分，是对中东局势影响的量化分析。Accenture在3Q因中东局势损失了约1亿美元的销售额和4亿美元的订单，主要集中在制造业和汽车行业。管理层明确表示，4Q的影响可能大于3Q。

这一判断的意义在于：它揭示了一个被市场低估的系统性风险。地缘政治摩擦不再是“黑天鹅”，而是正在成为影响IT服务企业基本面的结构性变量。对于全球化程度高的企业，这种风险敞口将直接影响其订单可见度和收入稳定性。

报告特别指出，大型管理服务项目的推迟更多是“延迟”而非“丢失”，这在一定程度上缓解了市场的担忧。但问题在于，延迟意味着收入确认的节奏被打乱，对企业的现金流和利润率都会产生连锁反应。更重要的是，如果地缘政治局势持续恶化，这些延迟的项目最终可能演变为取消。

GS对日本IT服务板块的传导分析提供了一个清晰的观察框架。报告认为，中东局势对日本本土IT服务需求的影响有限，但对其子公司NTT Data的轻微拖累值得关注，因为后者海外收入占比高达60\%。

这一判断的底层逻辑是：日本IT服务企业的业务结构正在发生根本性变化。过去，投资者习惯于将日本IT板块视为一个相对封闭、受国内需求驱动的市场。但NTT Data、NOM综合研究所等企业的全球化布局，正在改变这一假设。

报告指出，日本市场也开始出现与全球类似的AI相关业务扩张趋势，NOM综合研究所和趋势科技被认为处于有利位置。这意味着，日本IT服务板块的投资逻辑正在从“国内需求周期”转向“全球AI需求+地缘政治风险”的双重驱动。

投资者需要重新审视自己的持仓组合：那些海外收入占比高、且在中东或欧洲有重大业务敞口的企业，其估值可能需要纳入新的风险溢价。而那些以日本本土市场为主、同时具备AI业务能力的企业，可能更具防御性。

Ac

[... middle omitted ...]

完全传导至管理服务。

对于日本市场而言，NOM综合研究所和趋势科技被认为将受益于这一趋势。NOM综合研究所在金融IT领域的深厚积累，使其在AI驱动的业务转型咨询中具有天然优势。而趋势科技在网络安全领域的布局，则契合了AI部署带来的安全需求增长。

但这里仍有一个关键问题需要验证：AI相关业务的利润率是否能够达到传统IT服务的水平？如果AI业务需要更高的研发投入和更短的交付周期，其对盈利能力的贡献可能不如预期。

尽管GS报告提供了丰富的分析框架，但仍有几个关键问题没有被完全解答。

第一，管理服务订单的疲软是暂时的还是结构性的？Accenture管理层将大型项目的推迟归因于客户特定因素，而非行业趋势。但考虑到地缘政治风险可能持续，这种“延迟”的持续时间存在不确定性。

第二，AI业务的扩张能否完全对冲地缘政治风险带来的收入缺口？报告显示AI相关需求强劲，但并未给出具体的量化预测。如果中东局势的影响在4Q进一步扩大，AI业务的增量能否弥补这一缺口，仍是一个开放性问题。

\subsection{[R026] JPM：能量饮料市场真正进入“买量竞争”阶段，品牌分化才刚刚开始}
{\small\color{refgray}归类：中国地产与消费：量稳价跌与盈利预期重置}

JPM：能量饮料市场真正进入“买量竞争”阶段，品牌分化才刚刚开始

这份JPM基于Numerator数据发布的2026年5月能量饮料月度监测报告，表面上是对品牌KPI的常规追踪，但如果你只看到“增速放缓”四个字，就错过了最重要的信号。

报告真正揭示的判断是：美国能量饮料市场正在从“品类渗透红利期”转向“存量博弈期”。过去两年支撑高增长的逻辑——新用户涌入、购买频次提升、提价空间——正在同步收窄。这并不意味着行业见顶，而是竞争逻辑发生了根本性切换。

JPM分析师Andrea Teixeira团队在报告中用大量连续13周滚动数据证明：品类整体投影销售额增速从4月的+15\%左右降至5月的+10\%，减速幅度超过500个基点。更值得关注的是，减速并非单一因素驱动，而是渗透率增速、购买频次、客单价三个引擎同时降温。

这意味着什么？意味着过去“水涨船高”的阶段结束了。接下来，品牌之间的分化将急剧扩大。有些品牌会在减速中暴露脆弱性，有些则能凭借结构性优势守住阵地。这份报告的价值，不在于告诉你“谁在增长”，而在于告诉你“谁的增长更可持续”。

1. 品类增速放缓的根源不在需求消失，而在“低频用户”的购买行为正在回归均值

报告中最容易被误读的数据是品类整体家庭渗透率仍在创新高——5月13周滚动渗透率达到55.0\%，较4月提升约80个基点。看起来市场还在扩张，但为什么销售增速反而降了？

关键在于购买频次。品类整体购买频次同比增速在5月转为-2\%，这是自2025年11月以来首次转负。而客单价增速虽然仍保持在+5\%左右，但绝对金额环比已出现微降。

JPM的数据提供了一个关键视角：品类渗透率提升主要来自新用户的首次尝试，但这些新用户的复购率和购买深度远不如核心用户。当新用户增量边际递减，而老用户的购买频次无法维持高位时，增速自然承压。

从渠道拆解来看，减速最严重的是加油站和便利店渠道——购买频次同比转负至-2\%，这是唯一同比下滑的渠道。考虑到能量饮料约60\%-70\%的销售发生在这一渠道，这个信号不可忽视。报告没有直接给出解释，但合理的推断是：加油站渠道的用户画像偏向低频、冲动型购买，他们对价格和促销更敏感，在经济环境不确定性增加时，这类购买行为最先收缩。

所以，品类增速放缓不是需求消失，而是“浅层需求”的购买行为正在向均值回归。真正考验品牌的是：你能否把那些尝鲜用户转化为高忠诚度的重复购买者？

2. Monster在“相对韧性”中暴露了增长引擎的结构性换挡

Monster是报告中被分析师评价为“优于同行”的品牌。从绝对数据看，家庭渗透率达到32.7\%的新高，客单价增速连续23个月保持中高个位数增长。这些指标确实领先于红牛和Celsius。

但仔细看环比变化，Monster的韧性背后有隐忧。购买频次增速连续三期放缓，5月已降至同比持平；买率增速从之前的高双位数骤降至+6\%。虽然仍高于红牛，但斜率变化值得警惕。

JPM的分析师用了一个微妙的措辞：“数据比同行更好，但KPI在环比减速。”这意味着Monster正在经历从“高增长驱动”到“成熟品牌维持”的换挡。客单价持续增长说明定价权仍在，但购买频次的停滞暗示，品牌可能已经触及了核心用户的使用上限。

对于投资者而言，Monster的价值在于其品牌壁垒和渠道深度，但未来的增长驱动力必须从“量增”转向“价增”和份额从竞品手中夺取。报告中提到Monster的投影销售额增速从4月到5月下降了314个基点，而Circana数据显示的减速幅度更温和（-70个基点），两者之间的差异本身就是一个值得追踪的信号——不同数据源对渠道覆盖的差异，可能暗示某些渠道（如线上或小众零售）正在发生结构性变化。

3. Celsius的“双品牌战略”正在经受压力测试，Alani Nu的爆发力掩盖不了母品

[... middle omitted ...]

虽然仍在吸引新用户，但现有用户的购买节奏在放慢。

Celsius品牌则面临更深层的问题：家庭渗透率连续四个月同比下滑，买率增速降至同比持平。虽然绝对渗透率环比仍在改善（+35个基点），但趋势方向与Alani Nu形成鲜明对比。

JPM在报告中提到Celsius“绝对渗透率和购买频次比同比趋势表现得更好”，这其实是一个谨慎的措辞。翻译过来就是：虽然同比数据不好看，但环比没有继续恶化。对于一家被市场寄予高增长预期的公司来说，这个“环比企稳”的叙事能支撑多久，取决于Alani Nu能否持续超预期。

这里有一个值得深入思考的问题：Celsius Holdings的双品牌战略（Celsius主品牌+Alani Nu）在逻辑上互补，但数据表明两个品牌面对的是不同的消费群体和购买场景。Alani Nu的高增长更多来自年轻女性消费者的渗透，而Celsius主品牌的核心用户可能更偏向健身和运动场景。当两个品牌的增长节奏出现背离时，公司管理层面临的资源分配和渠道策略调整将变得复杂。

\subsection{[R027] 摩根斯坦利：东盟电信股真正的回报来源不是5G，而是竞争格局的结构性改善}
{\small\color{refgray}归类：区域市场：日本盈利驱动与韩国K型复苏闭合}

摩根斯坦利：东盟电信股真正的回报来源不是5G，而是竞争格局的结构性改善

市场对东盟电信板块的关注，长期集中在两个叙事上：一是5G资本开支周期何时见顶，二是数据中心需求能否带来第二增长曲线。但摩根斯坦利在最新发布的《Asia Summer School 2026: ASEAN Telecoms \& Data Centers》报告中给出了一个更值得深思的判断：决定未来12-18个月电信股超额收益的核心变量，不是技术升级，而是各国市场竞争格局的结构性变化。

这份报告覆盖了泰国、新加坡、马来西亚、印度尼西亚和菲律宾五个市场，并基于竞争、资本开支、监管、资本配置和估值五个维度进行了系统打分。结果并不令人意外：泰国和新加坡排名前二，菲律宾垫底。但真正有价值的不是排名本身，而是排名背后的逻辑——竞争格局的改善，正在重塑这些市场的投资回报框架。

摩根斯坦利的评分体系中，泰国以23分位居第一，新加坡以22分紧随其后。这两个市场的共同特征是什么？都是事实上的双寡头或准双寡头格局。

新加坡的情况则更为微妙。表面上看，新加坡是四家运营商加多家MVNO的竞争格局，但摩根斯坦利重点关注的是SingTel的份额稳定性和SIMBA的宽带价格战影响。报告的评分显示，新加坡在“竞争”维度的当前得分仅为1（满分3），但“展望”得分却跳升至3——这意味着分析师预期竞争环境将在未来改善。这一转变的关键变量，可能是行业整合的推进。报告提到SIMBA与M1的交易已暂停，但SingTel对参与整合表现出兴趣。如果整合发生，新加坡将从四家玩家变为三家，竞争烈度将显著下降。

对于投资者而言，这意味着泰国的寡头红利已经兑现，而新加坡的寡头红利仍在酝酿中。泰国的投资逻辑更偏向“持有并收取股息”，新加坡则可能提供“预期差带来的估值修复”。

印尼以20分与马来西亚并列第三，但它的故事远比分数复杂。摩根斯坦利对印尼的分析揭示了一个关键洞察：竞争格局正在出现结构性分化。

从行业参与者数量看，印尼已经从2008年的10家运营商合并为3家——Telkomsel、Indosat和XLSmart（由XL与FREN合并而来）。合并后的市场集中度提升，竞争在Java岛（印尼人口最密集的区域）确实有所改善。报告提到“竞争自2025年下半年以来有所改善”，三个运营商的ARPU正在向45,000印尼盾收敛。

但问题出在Java之外。报告明确指出，Indosat和XLSmart在合并后变得更强大，它们有更强的动机和资源在外岛（ex-Java）扩张网络覆盖。这意味着，虽然Java岛的竞争趋稳，但外岛的竞争可能加剧。这是一个典型的“局部改善、整体复杂”的局面。

对于投资者而言，印尼电信股的投资需要区分两个维度：一是Java岛的价格战趋缓带来的ARPU支撑，二是外岛资本开支增加对自由现金流的侵蚀。这两个力量的净效果，决定了印尼电信股的回报质量。报告在评分中给印尼的“资本开支”维度当前得分仅为1（最低分），说明市场已经定价了较高的资本开支压力，但问题在于市场是否充分定价了外岛竞争加剧的可能性。

3. 菲律宾的困境：第三玩家“赖着不走”，行业结构修复被推迟

菲律宾以18分排名垫底，这是五个市场中唯一一个在“竞争”和“资本配置”两个维度的展望得分都出现下滑的市场。摩根斯坦利的分析点出了一个关键问题：第三玩家Dito仍在市场中“游荡”。

菲律宾电信市场历史上是PLDT与Globe的双寡头格局，但Dito在2021年进入后，市场变成了三家竞争。报告没有直接说Dito何时退出，但数据已经说明问题：Dito的市场份额虽然增长缓慢，但始终没有消失。更关键的是，菲律宾的宽带普及率虽然预计每年增长2个百分点，但ARPU占家庭收入的比例高达2.5\%，远高于东盟平均的1.6\%，这意味着价格上行空间极

[... middle omitted ...]

 报告尚未完全回答的关键问题：数据中心需求能否对冲电信主业的压力？

摩根斯坦利这份报告的主题虽然是“ASEAN Telecoms \& Data Centers”，但对数据中心的分析相对克制。报告提到电信运营商正在从东盟地区的数据中心需求增长中受益，但没有量化这一受益对运营商整体盈利的贡献程度。

这是一个值得深入探讨的问题。电信运营商拥有的土地储备、光纤网络和电力基础设施，使其在数据中心建设中具有天然优势。但数据中心的资本开支强度远高于传统电信业务，且回报周期更长。如果运营商为了抢占数据中心市场份额而大幅增加资本开支，可能会抵消电信主业竞争改善带来的现金流改善。

此外，数据中心的客户集中度（主要来自超大规模云服务商）意味着运营商的议价能力有限。在电力成本上升的背景下，数据中心业务的毛利率能否维持在高位，仍是一个需要验证的假设。

这份报告给投资者最重要的启示，是东盟电信股的投资框架需要从“行业周期”转向“市场格局”。具体而言，投资者应该建立以下三个层次的观察框架：

\subsection{[R028] NOM：四足机器人已越过盈亏线，但真正考验在于“大脑”尚未定型}
{\small\color{refgray}归类：中国医药与机器人：全球定价权再校准与‘大脑’路线未定}

NOM：四足机器人已越过盈亏线，但真正考验在于“大脑”尚未定型

2026年6月15日，NOM分析师团队实地走访了杭州的深度机器人公司（DEEP Robotics，未上市），并在18日该公司科创板IPO申请获受理后，发布了这份调研报告。这份报告的价值不在于它描述了机器人行业的火热——那是共识——而在于它提供了三个层次上市场尚未充分定价的判断：第一，这家四足机器人专业厂商在2025年首次实现全年盈利，营收同比增长227\%至3.37亿元，净利润达到2900万元；第二，电力巡检和消防这两个垂直场景的需求正在从“试点”走向“批量部署”，但价格和性能仍然是规模化复制的硬约束；第三，也是最容易被忽视的一点，行业通往“物理AI”的技术路线尚未收敛，这意味着当前所有关于远期市场空间的线性外推，都可能建立在一条尚未选定的路径上。

以下是我们从这份研报中提炼出的五个关键洞察，以及一个NOM尚未完全回答的关键问题。

1. 2025年首次盈利是一个里程碑，但更值得关注的是毛利率与客户集中度的组合信号

营收从2023年到2025年的复合年增长率达到159.5\%，2025年产量3936台、销量2908台，产销率73.88\%——这些数字本身已经能说明公司处于高速成长期。但报告中最值得推敲的财务信号是两个并存的事实：52.83\%的毛利率和18.83\%的前五大客户集中度。

毛利率超过50\%在硬件制造领域并不常见，尤其对于一家仍在快速扩张期的公司。这意味着产品定价权或成本结构有某种结构性优势——可能是四足机器人在特定场景下的不可替代性，也可能是公司在核心零部件或算法层面的自研壁垒。但与此同时，不到20\%的前五大客户集中度暗示收入来源相对分散，这在一定程度上降低了单一客户依赖风险，但也意味着销售和渠道成本可能较高，因为服务大量中小客户所需的部署、维护和定制化投入并不低。

对于投资者而言，这里需要追问的是：毛利率能否在收入规模继续扩大的过程中维持甚至提升？如果答案是否定的，那么当前的盈利能力可能只是阶段性红利，而非长期竞争壁垒的体现。

2. 电力与消防场景的需求正在真实落地，但“智能仪表”方案可能在未来侵蚀这一市场

报告中提到的东莞变电站案例——每天约8.5小时巡检，识别准确率超过95\%，使用双目摄像头辅以声学和温度传感——是一个难得的、有具体参数的部署证据。这不是概念验证，而是已经跑通了的运营场景。湖南、山西、河南等省份的消防和巡逻单位也在推进部署，海外项目触及新加坡和北美。

然而，NOM的外部行业调查提出了一个值得警惕的判断：随着变电站数字化转型推进，“智能仪表”方案——即变电站自身设备具备自感知能力——可能逐步侵蚀四足机器人在变电站巡检中的机会。这本质上是一个“技术替代”风险：如果变电站自己就能完成数据采集和异常识别，为什么还需要一台四足机器人？

这意味着，电力巡检这个看似确定的“近岸锚点”场景，其市场天花板可能比直觉上看到的要低。四足机器人厂商需要在这个窗口期内，快速证明自己不止能做巡检，还能做那些“智能仪表”做不到的事——比如在复杂地形中执行物理操作、在高温或粉尘环境下进行应急响应。

3. 工业垂直场景的规模化采用仍需等待，认证、载荷和续航是三个硬约束

报告提到，管理层认为化工、石化、矿业中的防爆机器人是一个更大、更具防御性的机会。但紧接着，报告也给出了为什么这个目标很难实现的原因：认证、载荷、续航，以及四足底盘本身的形态限制。

防爆认证在化工和矿业领域是一个漫长且高成本的过程，涉及多个国家和行业标准。四足机器人的载荷能力有限，意味着它无法携带大型工具或设备执行复杂任务。续航问题则限制了单次作业的时长和覆盖范围。而四足底盘在狭窄管道、高密度设备区等特定工业环境中的适用性，也需要进一步验证。

这些约束本质上指向同一个

[... middle omitted ...]

重要的信息是：高层的技术路线尚未收敛。

目前学术界的主流方向正在从VLA风格向World Model迁移，但“世界”的定义和场景理解在不同参与者之间存在显著分歧，没有形成明确的时间表。相比之下，“小脑”部分——越来越多地基于强化学习模型——被认为已经基本确定。

这意味着什么？如果你是一家四足机器人公司，你的“大脑”路线选择将决定你未来五年的研发投入方向、人才招聘策略、甚至硬件架构设计。选错了路线，可能意味着数亿美元的沉没成本和数年的时间损失。而对于投资者而言，这意味着当前所有关于“物理AI将如何改变机器人行业”的估值模型，都建立在一个尚未被验证的假设之上。这不是一个可以忽略的风险。

5. 报告尚未完全回答的关键问题：供应链的可扩展性如何验证？

NOM在这份报告中给出了对拓普集团的中性评级，理由是“对人形机器人贡献的可持续性和规模持谨慎态度”，并预计其核心业务增长将在2026年放缓。这是一个合理的判断，但它回避了一个更根本的问题：四足机器人供应链本身的可扩展性。

\subsection{[R029] NOM：宇树科技IPO揭示的真正悬念，不在硬件成本，而在“大脑”路线的未定}
{\small\color{refgray}归类：中国医药与机器人：全球定价权再校准与‘大脑’路线未定}

NOM：宇树科技IPO揭示的真正悬念，不在硬件成本，而在“大脑”路线的未定

当一家公司以5,500台人形机器人出货量位居全球第一，营收同比增长335\%至17亿元人民币，并实现2.88亿元的净利润，毛利率约60\%，它即将在科创板挂牌的消息，很容易被市场解读为“人形机器人商业化拐点已至”。

但NOM在6月15日实地调研宇树科技后发布的这份报告，其真正有价值的判断并非关于出货量或毛利率。报告的核心洞察指向一个更根本、也更令人不安的结论：行业当前最硬的约束条件，已经从“能不能动”转向了“能不能想”——而后者，整个行业都还没有找到确定的路径。

这份报告的价值，在于它没有停留在对宇树全栈自研能力的赞美，而是清晰区分了“已被验证的成本优势”与“尚未被验证的模型能力”。对于关注这一赛道的投资者而言，理解这两者之间的张力，比记住任何一组财务数字都更重要。

1. 宇树的成本优势已构成护城河，但这条护城河可能只保护了一半战场

NOM调研报告中最引人注目的数据，是宇树的成本结构：电机、减速器、激光雷达全部自研，外采部件仅占整体成本的14\%至18\%。配合超过90\%的产销比，这组数字共同指向一个判断：宇树在人形机器人的“物理层”建立了一个相当完整的成本护城河。

这一点在行业内并非普遍现象。许多竞争对手仍然依赖外购核心部件，这意味着它们在定价权和供应链韧性上天然处于劣势。宇树能够在2025年实现盈利，很大程度上正是得益于这种全栈自研带来的成本优势。管理层的表述也很清晰：成熟的四足机器人业务产生的现金流，正在反哺人形机器人的研发和市场拓展。

但这里有一个关键问题需要区分：成本优势解决的是“能不能造得便宜”，而人形机器人产业当前的核心矛盾，已经不再是制造成本。NOM调研中引用的外部行业调查明确指出，行业真正的瓶颈在于“大脑”——即通用世界模型的能力。硬件成本再低，如果机器人无法在复杂环境中自主决策和执行任务，它仍然只是一个昂贵的可动装置。

这意味着宇树的成本护城河是真实且有价值的，但它保护的战场，可能只是整个战役的前半段。后半段战役——关于智能水平的竞争——才刚刚开始，而且胜负远未确定。

NOM报告中一个被低估的判断，是关于“大脑”路线的分歧。报告指出，领先研究机构已经从VLA模型路线转向世界模型路线，但对于“世界”的定义和场景理解方式，各家之间仍然存在显著差异。

这句话的信息密度很高。它意味着当前人形机器人行业面临的根本问题，不是某家公司比另一家强多少，而是整个行业还没有就“如何让机器人理解世界”达成共识。高层次的路径没有收敛，子路线也没有收敛，因此不存在一个被确认的时间表。

这对所有硬件厂商都是一个结构性风险。如果最终的模型路线被AI巨头主导——正如NOM引用的行业调查所暗示的——那么硬件厂商将面临持续的压力。它们可能会发现自己投入巨资开发的硬件平台，需要适配一个并非由自己定义的智能系统。这种适配过程不仅成本高昂，而且可能随时面临路线变更带来的沉没成本。

宇树管理层显然意识到了这一点。NOM报告提到，公司计划将IPO募集资金的50\%以上用于“大脑”开发。这是一个正确的战略方向，但也暴露了问题的严重性：连全球出货量第一的人形机器人公司，都不得不将过半资源投入一个尚未确认路径的领域。

3. 学术收入占比74\%：这不是短板，而是产业成熟度的真实写照

很多人看到“74\%的收入来自大学和科研机构”这个数字，可能会认为这是宇树的短板。但NOM报告的处理方式值得借鉴：它没有将其简单定义为问题，而是将其作为产业成熟度的标尺。

人形机器人目前仍处于“工具被研究者使用”的阶段，而非“工具被产业使用”的阶段。电网巡检和消防是增长最快的工业应用场景，但报告明确提到，其他垂直领域主要是“需求拉动”而非主动开拓。公司将自己定位为通用机器人

[... middle omitted ...]

. 报告尚未完全回答的关键问题：成本优势能否转化为智能生态的入口

NOM的调研提供了一个清晰的框架，但有一个问题报告没有完全展开：宇树的成本优势，在多大程度上能够转化为智能生态的入口优势？

逻辑上，更低成本的硬件意味着更大的装机量，更大的装机量意味着更多的运行数据和反馈数据，这些数据对于训练“大脑”至关重要。如果这个闭环成立，宇树的成本优势就不仅仅是制造优势，而是数据优势的起点。

但这里有两个未经验证的假设。第一，宇树目前的出货量（5,500台）是否已经达到能够产生有效训练数据的规模阈值？第二，这些数据是否能够被有效地收集、标注和用于模型训练？宇树的技术栈在硬件层面是全栈自研，但在软件和数据层面，报告没有提供足够的信息来判断其能力。

这两个问题，可能是决定宇树能否从“硬件冠军”进化为“智能平台”的关键。如果数据闭环不能有效建立，宇树最终可能只是一家优秀的硬件制造商，而智能层的价值将被AI巨头攫取。反之，如果数据闭环成立，宇树的竞争壁垒将从成本扩展到数据和生态。

\subsection{[R030] NOM：机器人视觉的竞争壁垒不在“手”，而在“眼”和“脑”的数据飞轮}
{\small\color{refgray}归类：中国医药与机器人：全球定价权再校准与‘大脑’路线未定}

NOM：机器人视觉的竞争壁垒不在“手”，而在“眼”和“脑”的数据飞轮

机器人行业正在经历从“机器换人”到“智能体替代”的范式迁移。市场注意力大多集中在人形机器人的整机厂商和灵巧手供应商上，但一份NOM近期发布的研报，通过对国内机器视觉头部企业梅卡曼德（Mech-Mind Robotics，未上市）的实地调研，揭示了一个可能被低估的竞争逻辑：在机器人智能化的价值链中，“眼”（3D视觉）和“脑”（算法模型）所构成的认知层，才是当前最具规模效应和护城河潜力的环节，而“手”的执行层，尚处于商业化早期。

这份报告的核心判断是：机器人产业链的价值分配，正在从硬件集成向数据驱动的认知平台迁移。 梅卡曼德的案例表明，谁能在“眼-脑”层面率先形成数据闭环，谁就有机会在长尾、多样化的工业场景中建立标准化部署的壁垒。以下是我们基于这份研报解析提炼出的五个关键洞察。

1. 连续六年市占率第一的背后，是“数据飞轮”而非单纯的硬件优势

NOM在研报中指出，梅卡曼德在大型3D抓取领域（如拆垛）已连续六年保持国内市场份额第一。管理层的表述是：“做得多，数据就多，数据多就能改进模型。”这听起来像一句常见的商业口号，但其背后的竞争逻辑值得深究。

在机器人视觉领域，单纯的传感器硬件（如深度相机）的差异化正在缩小。真正拉开差距的是算法对海量、多场景数据的消化能力。梅卡曼德的策略是不生产机器人本体，而是适配全系列机械臂。这意味着，每一次部署——无论是天津包裹处理中心每小时处理1800件包裹的场景，还是其他工业产线——都在为其“脑”积累新的训练数据。这种“眼”与“脑”的协同，构成了一个自我强化的正反馈循环：部署越多，模型越通用和稳定；模型越强，部署的门槛和成本就越低。

这给我们的启示是：在评估机器人产业链公司时，不应只看其硬件的技术参数，更应关注其数据的“生成-回流-迭代”机制是否成立。那些只能卖硬件、无法积累场景数据的公司，长期来看可能会被数据飞轮驱动的平台型公司拉开差距。

2. “眼-脑-手”平台的核心诉求是标准化，但“手”仍是最大的不确定性

梅卡曼德将其核心收入来源定义为“眼”（3D视觉）和“脑”（算法/模型），并据此构建了“眼-脑-手”（EBH）平台战略。管理层的愿景是推动行业从重度定制的集成模式，转向标准化、易部署的解决方案。为此，公司每年培训 名客户人员，以降低使用门槛。

这一战略的逻辑是清晰的：工业机器人需求具有“广泛、长尾、多样化”的特征，传统集成商模式效率低、规模不经济。如果能通过标准化的“眼-脑”平台，将复杂的机器人编程和调试简化为“即插即用”，那么梅卡曼德就有机会从项目制公司转型为平台型公司。

然而，NOM在研报中也坦诚地指出了行业调研中的谨慎看法：梅卡曼德的优势在于视觉，而非运动控制。灵巧手（dexterous hand）尚未产生规模收入，且自研成熟灵巧手需要时间，甚至可能需要股权激励来吸引人才。这意味着，尽管“眼-脑”平台可以适配不同形态的机器人本体和末端执行器，但“手”的缺失，可能会限制其在需要精细操作的高附加值场景中的渗透。

这里有一个尚未完全回答的关键问题：“眼-脑”平台的先发优势，能否在“手”的短板补齐之前，支撑起足够高的客户粘性和商业壁垒？ 如果“手”的成熟周期过长，竞争对手可能在“眼-脑”层面迎头赶上。

3. 与奥比中光的差异化定位：通用市场 vs 高价值利基市场

NOM将梅卡曼德与A股上市的奥比中光（688322 CH，买入评级）进行了对比，这一对比对于理解3D视觉赛道的竞争格局至关重要。

报告的核心结论是：两者并非直接竞争，而是瞄准了不同细分市场。奥比中光定位于通用市场，覆盖支付刷脸、社保认证、服务机器人乃至人形机器人等大规模应用场景。而梅卡曼德，根据行业调研，更适合那些需要移动、高精

[... middle omitted ...]

，而应区分“硬件模组商”和“解决方案平台商”，前者看规模，后者看场景深度和客户粘性。 两者的估值逻辑和风险点截然不同。

4. 报告尚未完全回答的关键问题：技术领先优势的窗口期有多长？

NOM在研报中坦承，尽管梅卡曼德在拆垛领域保持领先，但其相对于海康机器人（Hikrobot，未上市）等竞争对手的技术领先优势可能正在缩小。这是一个非常关键的风险提示。

在科技行业，先发优势往往伴随着“追赶者压力”。一旦核心算法和视觉方案被竞争对手通过逆向工程或差异化路径突破，梅卡曼德的数据飞轮效应可能会面临挑战。特别是当海康机器人这样拥有强大硬件制造能力和渠道网络的对手入场时，竞争格局可能会快速变化。

报告没有完全展开的是：梅卡曼德的技术护城河，究竟有多少来自专利和算法秘密，又有多少来自其积累的行业数据？ 如果后者是核心，那么数据壁垒的牢固程度，取决于其客户数据的排他性以及模型对特定场景的“过拟合”程度。如果模型过于依赖特定场景数据，那么在跨行业迁移时，其通用性优势可能会打折扣。

\subsection{[R031] 美国银行：市场正在定价的不是繁荣，而是繁荣之后的“谁先离场”}
{\small\color{refgray}归类：宏观与利率：美联储沟通范式切换与全球央行分化 / 风险偏好与资金流向：极端集中与泡沫特征重现}

美国银行：市场正在定价的不是繁荣，而是繁荣之后的“谁先离场”

这份报告最值得关注的判断，不是资金流入了多少，而是资金流入的分布结构，正在复刻历史上泡沫末期的典型特征。美国银行最新一期的《The Flow Show》显示，过去一周全球股票基金净流入1264亿美元，创下历史纪录。其中美国股票单周流入1192亿美元，同样刷新历史。科技股单周流入192亿美元，也创下纪录。如果按年化推算，2026年全年美国股票基金流入将达7390亿美元，科技股流入将达1540亿美元。

这些数字本身并不令人惊讶。真正值得警惕的是，当资金以这种速度和集中度涌入单一市场、单一风格、单一主题时，历史上往往对应的是泡沫的“最后一段”。美国银行的牛熊指标已升至9.2，进入“极端看多”区间，这是一个自2002年以来触发过17次、命中率约60\%、平均回撤2-3\%、最大回撤15-20\%的卖出信号。

这份报告的核心主张是：当前市场定价的核心变量，已经从“经济增长能否持续”，切换为“政治周期与收益率曲线谁会先打破泡沫”。这不是一篇关于基本面复苏的乐观叙事，而是一篇关于“谁先眨眼”的博弈论分析。

美国银行的报告提供了一个极为关键的历史对照。当前AI相关十大股票占标普500指数的权重已达到39\%。这个数字本身已经接近甚至超过了历史上几次著名泡沫的集中度峰值：19世纪80年代的铁路股、1972年的“漂亮50”、1999年的TMT泡沫，以及1985年日本股市占全球市值的巅峰。

图表3的横截面比较非常清晰：每一次市场集中度达到这个水平，都对应着一场泡沫的尾声。这不是说泡沫马上会破裂，而是说“继续集中”的空间已经极为有限。报告进一步用1998年至2000年的历史数据说明：在泡沫的最后6个月，只有科技和电信板块保持正收益，其他所有板块都在下跌。

这意味着什么？如果当前的结构性行情继续下去，投资者实际上是在押注这一次“这次不一样”。但报告给出的隐含判断是：泡沫的终结通常不是由基本面恶化触发，而是由资金面的自我强化机制被外部冲击打破——比如政治投票、收益率曲线倒挂，或者汇率危机。

报告的一个核心洞察是：泡沫的终结通常由“选民和义警”完成。这里的“义警”指的是债券市场——当投资者开始用抛售债券的方式表达对通胀或财政纪律的不满，收益率曲线就会给出信号。

报告特别关注了美国11月的中期选举。如果共和党在参议院失利，报告将其定义为“美元下行、收益率下行、股票下行”的宏观场景。这个判断的逻辑链条是：共和党控制参议院意味着财政扩张的延续，而民主党控制参议院则意味着加税和监管收紧，两者对资产价格的含义截然不同。

但更关键的时间节点可能在9月。报告指出，如果特朗普的民调支持率在9月之前没有明显反弹，市场会变得“躁动”。为什么是9月？因为9月是美联储9月议息会议的前一个月，也是市场开始为年末流动性变化做准备的窗口。如果政治不确定性叠加收益率曲线的信号，泡沫的脆弱性会急剧上升。

这里需要强调的是，报告并没有给出“泡沫必然破裂”的判断，而是提供了一个观察框架：政治周期的变化，是泡沫能否延续的“外部约束”。这个约束目前尚未被充分定价。

3. 收益率曲线倒挂的真正含义：从“通胀繁荣”到“滞胀恐慌”

报告对收益率曲线的分析值得仔细拆解。自1月30日Warsh被提名为美联储主席以来，美国2年期与10年期国债收益率利差从75个基点迅速收窄至25个基点。这被报告定义为“熊市平坦化”，其背后是市场对“通胀等于加息”的恐惧。

报告的核心图表8揭示了一个历史规律：每当失业率低于CPI同比增速时，收益率曲线都会出现倒挂。当前美国失业率约为4.0\%，而CPI同比在3\%附近，两者差值正在收窄。如果通胀无法回落至3\%以下，收益率曲线将进一步平坦化，甚至重新倒挂。

报告的逻辑链条是：当前名义GDP

[... middle omitted ...]

发女孩”的乐观叙事。

报告提供的资金流向数据，最大的看点不是总量，而是结构。美国股票单周流入1192亿美元，但欧洲股票连续10周流出，中国股票连续12周流出，新兴市场整体也在流出。

这是一个“K型”的资金格局：资金从全球其他市场抽离，涌入美国，尤其是美国科技股。报告图表4用韩国的例子生动地说明了这种分化的极端性：KOSPI指数在上涨，但韩元在贬值。这两个信号不可能同时正确——要么股市错了，要么汇率错了。这种“背离”往往是市场即将变盘的先兆。

美国银行的私人客户数据也印证了这种分化。过去四周，这些高净值客户在买入材料、TIPS和市政债，却在卖出日本、必需消费品和公用事业ETF。这说明，即使是同一批投资者，也在进行“防御性轮动”——他们在用债券和商品对冲科技股的风险，而不是简单地“全仓押注”。

这意味着，当前市场的核心矛盾不是“该不该买股票”，而是“该买什么股票”。资金集中涌入科技股的同时，其他板块正在被系统性减持。这种“一边倒”的仓位结构，本身就是脆弱性的来源。

\subsection{[R032] 摩根斯坦利：市场误读了陆家嘴论坛的信号，真正重要的是利率传导机制的重塑}
{\small\color{refgray}归类：外汇与资本流动：人民币定价机制微调与跨境框架重构}

摩根斯坦利：市场误读了陆家嘴论坛的信号，真正重要的是利率传导机制的重塑

市场对近期陆家嘴论坛的解读，主要集中在“资本账户开放是否转向”这一命题上。多数评论聚焦于监管收紧与QDII额度发放之间的表面矛盾，并试图从中判断政策风向的边际变化。但摩根斯坦利在论坛后发布的这份研报，提出了一个更为根本的判断：陆家嘴论坛最重要的信号，不是资本账户开放的速度，而是货币政策传导机制正在经历一次静默但关键的升级。 这份报告提醒我们，如果将全部注意力放在资本流动的“开关”上，可能会错过一个正在发生的、更具结构性意义的制度变革。

报告的核心洞察在于，中国人民银行宣布将隔夜逆回购利率的波动区间收窄至7天逆回购利率上下各25个基点（此前为+50/-20个基点）。这并非一次简单的技术调整。在摩根斯坦利看来，这一变化的方向性意义在于，它使得短期流动性管理更加规则化和透明化，从而为提高整个政策利率体系的传导效率铺平了道路。市场当前低估了这种“机制优化”对长期资产定价和银行体系行为模式的潜在影响，而高估了短期内总量宽松的可能性。

1. 利率走廊的收窄不是在暗示降息，而是在为未来的传导效率铺路

市场往往倾向于将任何利率机制的调整解读为“宽松”或“紧缩”的前奏。但摩根斯坦利的分析明确指出，这次利率走廊的收窄，其首要目标并非释放降息信号，而是改善政策利率向市场利率的传导效率。报告特别强调，在银行净息差压力高企的背景下，央行短期内不会降息。这意味着，2026年下半年乃至2027年，广义政策利率的调降可能性极低。

这一判断的逻辑链条是清晰的：一个过宽的利率走廊，使得短期市场利率容易大幅偏离政策目标利率，削弱了央行通过调整政策利率来引导实体经济融资成本的能力。收窄走廊，本质上是在为未来的每一次政策利率调整，赋予更强的市场约束力和信号意义。对于投资者而言，这意味着未来需要更加关注央行的“窗口指导”和“规则化操作”，而非单纯押注降息时点。银行间市场利率的波动性可能会降低，但银行资产负债表的定价能力将面临更严格的考验。

2. 资本账户开放没有“U型转弯”，而是在重新定义“合规”的边界

近期针对对外投资的一系列监管动作，被部分市场参与者解读为资本账户开放的倒退。摩根斯坦利在研报中提出了一个更具区分度的视角：这些措施的目的并非阻止资本流动，而是将资本流动引导至“受监管、可控制”的渠道。陆家嘴论坛上关于继续开放的表态，正是为了平息这种“U型转弯”的担忧。具体承诺包括：建立离岸人民币回购便利、在上海自贸区试点离岸人民币外汇交易、以及扩大QDII额度。

这实际上揭示了一个关键的政策分层逻辑：对于贸易、外商直接投资（FDI）和机构性资本流动，政策方向是持续的便利化和开放；而对于零售层面的资本外流、证券投资组合流动和短期资本，监管则保持审慎和收紧。这种“非对称开放”的格局，在“不可能三角”（独立的货币政策、稳定的汇率、自由的资本流动）的约束下，是北京方面在维持货币政策独立性和汇率稳定前提下的理性选择。投资者需要理解，合规的、机构化的跨境投资渠道将愈发重要，而试图绕开监管的灰色地带操作，将面临持续且加大的压力。

3. 人民币国际化的推进方式正在从“规模扩张”转向“基础设施搭建”

报告提到，央行承诺建立离岸人民币回购便利和试点离岸人民币外汇交易，这并非仅仅是为了增加人民币在国际支付中的使用量。更深层的意图在于，为离岸人民币市场构建起一个具备深度和流动性的“金融市场基础设施”。一个货币的国际化，最终取决于非居民是否愿意持有并交易该货币计价的资产，而这需要一个功能完备、风险可控的离岸市场作为支撑。

过去几年，人民币国际化的推进更多体现在双边本币互换协议和跨境贸易结算的规模增长上。陆家嘴论坛释放的信号表明，下一阶段的重点将转向“可交易性”和“可对冲性”。离岸人民币回

[... middle omitted ...]

在多重目标间寻求平衡的“次优选择”。

这一约束条件对资产定价的含义是深刻的。首先，它意味着债券市场的收益率下行空间可能被银行的行为所限制。其次，对于权益市场而言，银行股的投资逻辑需要从“息差扩张”切换到“资产质量改善”和“非息收入增长”。最后，它暗示了未来财政政策可能需要承担更大的逆周期调节责任，因为货币政策的操作空间已经受到银行体系健康状况的制约。投资者在评估政策组合时，应将银行净息差作为一个独立的、具有预测能力的变量纳入考量。

5. 报告尚未完全回答的关键问题：规则化传导能否有效替代总量宽松？

尽管摩根斯坦利对利率走廊改革的长期方向给予了积极评价，但一个尚未被充分回答的问题是：在企业和居民有效信贷需求偏弱的背景下，仅仅依靠改善利率传导机制，能否有效替代一次“降息”带来的总量刺激效果？机制优化解决的是“传导效率”问题，但当前经济面临的更多是“传导终端”的需求不足问题。如果实体经济缺乏加杠杆的意愿，即便政策利率能更有效地传导至市场利率，其刺激效果也可能大打折扣。

\subsection{[R033] GS：铝市场真正被低估的不是地缘溢价，而是两股供给冲击的时间错配}
{\small\color{refgray}归类：能源与大宗商品：供给约束下的时间错配与地缘风险折价}

GS：铝市场真正被低估的不是地缘溢价，而是两股供给冲击的时间错配

这份研报的核心判断，值得每一个关注大宗商品周期的决策者重新审视。

当前市场对铝价的讨论，大多集中在霍尔木兹海峡的地缘风险上。主流叙事是：中东供应中断推高了近端价格，一旦局势缓和，价格将迅速回落。GS的最新报告提出了一个更精细、也更有投资含义的框架——市场正在经历的不是单一供给冲击，而是两股方向完全相反的供给力量在时间轴上的错配。

第一股冲击来自中东。GS大幅下调了2026年和2027年中东地区的铝产量预测，下调幅度分别达到约66万吨和100万吨。这不仅仅是因为冲突本身，更关键的是复产节奏。报告指出，即使是宣布的临时协议达成、海峡重新开放，受损的电解槽修复和减产产能的逐步重启也需要接近一年的时间。参照2025年4月伊比利亚电网事故后Alcoa重启San Ciprian电解铝厂的经验，从停到满产大约需要12个月。这意味着，中东的供给缺口将延续至2027年。

第二股冲击来自亚洲，尤其是印度尼西亚和中国。GS上调了印尼2026年和2027年的原铝产量预测，分别至170万吨和290万吨，背后是Adaro、Taijing Morowali、Juwan Weda Bay等项目的加速爬坡，以及Harita Danantara Inalum从2027年起纳入产能统计。中国方面，GS也将2026年和2027年的产量预测上调至4560万吨和4630万吨，理由是强劲的行业利润正在推动重启、替代项目甚至部分超产。

这两股力量的交汇，决定了铝市场的真实图景：近端因中东缺口而收紧，远端因亚洲供给浪潮而宽松。GS将2026年Q3和2027年全年LME铝价预测分别上调至3300美元/吨和2950美元/吨，但仍显著低于远期曲线（3400美元/吨和3250美元/吨），核心逻辑就是——亚洲供给的增量最终会覆盖中东的损失。

这个判断的微妙之处在于：市场可能低估了亚洲供给的弹性和速度，同时高估了中东复产的灵活性。以下从五个层面展开。

GS报告中最具操作价值的一个假设，是中东受损产能的重启时间表。此前市场普遍预期，一旦地缘局势缓和，中东产量将在2026年下半年逐步恢复。但GS基于最新行业反馈和公司公告，将这一时间点推迟至2027年初。具体来看，巴林产量预计在2027年年中恢复至冲突前水平，阿联酋则要到2027年底。

支撑这一判断的关键是历史参照。2025年4月，西班牙伊比利亚电网事故导致Alcoa的San Ciprian铝厂全面停产，到2026年年中才完成重启，耗时约12个月。电解铝生产的特殊性在于，停产后电解槽内的电解质会凝固，槽体结构可能受损，重新通电后的焙烧和启动需要严格控制升温曲线，任何急躁的操作都可能导致槽壳破裂或漏炉。这不是一个可以“一键恢复”的过程。

这意味着，即使霍尔木兹海峡在2026年内完全恢复通航，中东地区的铝产量在2026年仍将低于正常水平约66万吨，2027年低于约100万吨。这个缺口，不是短期期货价格波动可以消化的。

对投资者的含义是：近端铝价的地缘溢价可能比市场预期的更持久。做空近端合约的风险，不仅来自冲突升级，更来自复产节奏本身的结构性慢变量。

2. 印尼正在成为全球铝供给的“黑马”，其速度可能超出多数人的预期

如果说中东是供给收缩的故事，那么印尼就是供给扩张的引擎。GS将印尼2027年原铝产量预测从250万吨大幅上调至290万吨，2028年进一步看至380万吨，2030年有望达到520万吨。这个增速，在近年来全球铝行业是罕见的。

支撑这一判断的证据来自两个方面。首先是实际数据：年初至今，印尼铝产量同比增长约89\%，项目进展基本符合加速路径。其次是专家反馈：GS通过行业专家调研确认，电力供应不再构成2026年项目的硬约束。印尼的铝冶炼项

[... middle omitted ...]

的规模效应，正在成为下一个低成本铝生产基地。

对投资者的含义是：不能再用旧地图看待全球铝供给。印尼的增量不仅影响远期平衡表，更可能改变铝的成本曲线结构。当印尼铝产量在2027年达到近300万吨时，它对全球铝价的边际定价权将显著上升。

GS报告中最引人注目的一个判断，是中国铝产量已经以运行率计算超过了4500万吨的官方产能天花板。2026年和2027年的产量预测分别上调至4560万吨和4630万吨。

这意味着什么？中国的产能红线并非绝对刚性。GS通过专家反馈指出，年初至今的产量强势可能来自三个渠道：未经授权的超产、产能置换计划下旧产能的持续运行、以及短期电解槽强化操作（即通过提高电流强度来增加单槽产量）。近期环保检查和报告的减产措施，实际影响相对较小。

这个判断如果成立，将改变市场对中国铝供给弹性的认知。过去几年，市场普遍认为4500万吨天花板是硬约束，中国铝产量将长期接近这一水平但不会突破。GS的报告提示我们：在高利润环境下，供给弹性可能比政策文件显示的要大。

\subsection{[R034] 摩根斯坦利：韩国K型复苏的裂口正在闭合，但市场低估了财富效应与政策节奏的共振}
{\small\color{refgray}归类：区域市场：日本盈利驱动与韩国K型复苏闭合}

摩根斯坦利：韩国K型复苏的裂口正在闭合，但市场低估了财富效应与政策节奏的共振

这份来自摩根斯坦利亚洲团队的研报，核心判断并非“韩国经济在复苏”——市场对此已有预期。真正值得关注的信号是：韩国经济正在从过去几年典型的K型复苏，转向更广泛、更均衡的全面反弹，且这一反弹的力度和节奏，可能超出当前市场定价。

报告预计韩国2026年GDP增速将从2025年的1.1\%跃升至2.8\%，2027年稳定在2.2\%。这个数字本身并不惊人，但支撑其实现的逻辑链条——科技出口超级周期、消费超预期回暖、财政空间意外打开、货币政策被迫转向紧缩——每一环都在挑战市场过去两年形成的“韩国经济结构性疲弱”的叙事。

以下是我们从这份报告中提炼的五个关键洞察，以及一个尚未被充分讨论的隐含风险。

1. 科技出口的“量价齐升”不是短期脉冲，而是产能周期与AI需求的深度耦合

报告中最具冲击力的数据是：2026年5月，韩国月度出口额创下历史新高，半导体出口增速同比超过60\%，且这一增长并非单纯由价格驱动。报告明确指出，DRAM和HBM的bit出货量预计在未来多个季度维持两位数增长，这意味着需求端的AI资本开支周期正在转化为实打实的产能利用率提升。

更关键的是，半导体在韩国总出口中的占比已从2024年的约30\%跃升至40\%以上。这并非简单的结构性集中度风险，而是意味着韩国出口的“质量”发生了根本变化——从过去依赖周期性大宗商品和汽车，转变为深度嵌入全球AI基础设施的供应链核心。报告中的资本品进口数据也印证了这一点：设备投资和资本品进口在2025年后显著加速，说明企业不仅在享受当前红利，更在为下一轮产能扩张做准备。

对于投资者而言，这意味着韩国科技股的盈利预测存在系统性上修的可能，而非仅仅是一次性的库存周期反弹。

2. 消费复苏的“财富效应”被普遍低估，权益市场渗透率是核心变量

市场对韩国消费的担忧主要集中在家庭债务和高利率的压制。但报告提供了一个被忽视的结构性视角：韩国是除美国和澳大利亚外，成年人持股比例最高的国家之一，接近50\%。这意味着股票市场的财富效应对消费的影响，远比其他亚洲经济体显著。

过去两年韩国股市的上涨，叠加证券投资者存款的持续攀升，正在通过两条路径传导至消费：一是直接的资产增值带来的“账面财富效应”，二是工资增长（报告暗示科技行业的强劲盈利正在扩散至服务业）带来的实际收入改善。摩根斯坦利预计私人消费增速将从2025年的不到1\%回升至2026年的2\%以上，且服务消费的反弹幅度更大。

这个判断的核心含义是：如果韩国股市保持强势，消费复苏的弹性可能超出线性外推的模型预测。反之，如果市场出现剧烈调整，消费的下行风险也会被放大——但当前市场显然更关注前者。

报告指出，由于科技企业盈利强劲带来的企业所得税、所得税和证券交易税收入超预期，韩国政府可能在2026年下半年推出额外预算。这是一个典型的“好上加好”的循环：科技繁荣推高税收，税收增加为财政扩张提供空间，财政扩张进一步支撑内需。

市场此前对韩国财政政策的预期相对保守，主要受制于过去几年房地产市场和家庭债务的约束。但报告揭示的逻辑是：科技出口的强劲表现正在为政府创造“财政意外之财”，这笔钱既可用于社会福利和消费刺激，也可用于支持新兴产业投资。无论哪种路径，都会对GDP增长形成额外的正向拉动。

这个变量的不确定性在于，政府是否会选择在经济已经明显复苏时追加刺激。但即使不追加，税收超预期本身也意味着财政纪律的改善，这对主权信用和汇率都是正面信号。

4. 韩国央行被迫转向紧缩，但市场对加息路径的定价可能仍不充分

报告预计韩国央行将从2026年7月起的一年内加息四次，终端利率达到3.50\%。这个预测的背景是：GDP增速将显著高于潜在增长率，CPI面临上行压力——不仅来自油价，更

[... middle omitted ...]

国资产的风险溢价可能需要重新评估。

5. 报告尚未完全回答的关键问题：科技与非科技部门的分化会如何演变？

报告虽然强调了“广泛复苏”，但数据本身也揭示了结构性分化的持续性。2026年5月的数据显示，扣除半导体和船舶后，出口增速仍高达22.5\%，这确实表明非科技部门也在改善。但改善的幅度和持续性，很大程度上取决于全球需求——尤其是中国和欧洲的经济表现。

报告没有深入探讨的是：如果全球AI投资周期在 年出现阶段性放缓，韩国经济能否依靠内需和财政支撑维持增长？换句话说，当前的“广泛复苏”有多少是科技繁荣的溢出效应，有多少是真正的内生动力？这个问题对判断2027年后的增长路径至关重要。

另外，韩国房地产市场的调整是否已经触底，还是仍存在尾部风险？报告对此着墨不多，但房地产作为家庭财富的重要组成部分，其走势将直接影响财富效应的可持续性。

第一，关注半导体出口的“量价拆分”。如果bit出货量增速持续高于出口金额增速，说明需求真实且健康；如果相反，则需警惕价格驱动的泡沫。

\subsection{[R035] 摩根斯坦利：日本制药板块的错位定价正在为深度价值投资者创造窗口}
{\small\color{refgray}归类：区域市场：日本盈利驱动与韩国K型复苏闭合}

摩根斯坦利：日本制药板块的错位定价正在为深度价值投资者创造窗口

这份2026年6月发布的摩根斯坦利日本制药行业报告，在当前市场情绪低迷的背景下，传递了一个值得产业决策者和长期资本认真审视的信号：日本制药板块的整体估值已经压缩至历史低位，但结构性分化远比表面市盈率所显示的更为剧烈。报告覆盖的18家公司中，摩根斯坦利给出了7个Overweight评级，但加权平均市盈率仅为15.8倍，远低于美国制药板块的21.6倍和欧洲的15.3倍。这个估值折价背后，是市场对日本制药企业全球化能力、专利悬崖应对以及研发管线转化效率的系统性低估。报告最核心的判断并非板块整体便宜，而是某些被市场情绪压制的头部标的，其风险收益比已经出现了不对称性。

1. 市场对日本制药的估值折价可能过度反映了历史包袱而非未来价值

当前日本制药板块加权市盈率仅为15.8倍，而美国同行高达21.6倍，折价幅度接近27\%。这一差距在历史上并非罕见，但值得关注的是折价的构成正在发生变化。过去十年，日本制药的折价主要来自三方面：本土市场增长停滞、国际化进程缓慢以及研发效率低于欧美同行。然而，摩根斯坦利的研究表明，这些因素的权重正在被重新评估。

以Daiichi Sankyo为例，该股年初至今下跌了25\%，当前股价仅为2,520日元，较摩根斯坦利给出的4,650日元目标价存在85\%的上行空间。这家公司是日本制药全球化的标杆，其ADC（抗体偶联药物）技术平台已被全球顶级药企认可，但市场似乎更关注其短期盈利波动而非管线价值。2026年预计EPS仅为126.4日元，但到2030年预计将增长至276.2日元，CAGR达到6\%。当前19.9倍的2027年市盈率，对于一家拥有全球差异化技术平台的公司而言，并不算昂贵。

关键在于，市场是否正在将日本制药的历史劣势过度投射到未来。如果一个公司的基本面结构已经发生改变，而估值框架仍然沿用旧有的折价逻辑，这就构成了定价错位。

2. Daiichi Sankyo的85\%上行空间并非线性外推，而是基于管线价值的重新定价

摩根斯坦利将Daiichi Sankyo列为Top Pick，目标价4,650日元，这一数字并非简单的市盈率倍数扩张，而是对其ADC平台在多适应症中潜在价值的综合评估。报告显示，该公司2025年EPS为210.2日元，但2026年预计降至163.8日元，2027年进一步降至126.4日元，随后在2028年反弹至176.7日元。这种U型盈利轨迹本身就说明，当前的低迷是研发投入期的阶段性特征，而非结构性衰退。

值得注意的数据是，Daiichi Sankyo的PEG比率为2.7，高于板块平均的3.4，但远低于其成长性所对应的合理水平。如果市场开始将其从“日本仿制药公司”重新归类为“全球创新药平台”，估值体系将发生根本性转变。当前12.0倍的2025年市盈率，对于一个拥有全球最先进ADC技术之一的企业而言，显然没有充分反映其技术壁垒。

但这里需要指出一个报告没有完全展开的问题：Daiichi Sankyo的管线价值能否在2030年前转化为可持续的现金流？ADC领域的竞争正在加剧，来自中国和美国biotech的追赶速度可能超出预期。摩根斯坦利的估值模型是否充分考虑了竞争格局的动态变化？这一点值得进一步验证。

Takeda是摩根斯坦利的另一个Overweight标的，目标价6,300日元，当前股价5,020日元，上行空间25\%。这家公司自收购Shire以来一直背负巨额债务，市场对其资产负债表始终存在担忧。但报告显示，Takeda的EV/EBITDA从2025年的7.7倍已经下降至2026年预计的8.6倍，并且预计到2030年将进一步降至7.6倍。同时，其P/B从2025年的1.0倍下降至2030年的0.2倍左右，这说

[... middle omitted ...]

Kaken和Towa——均被给予Overweight评级，目标价分别对应41\%和45\%的上行空间。Kaken的市值仅为1,660亿日元，Towa为2,000亿日元，属于典型的日本中小市值医药股。

Kaken的EPS增长路径尤为引人注目：2025年EPS仅为56.6日元，但2026年预计飙升至174.1日元，2027年进一步升至218.9日元，2028年回落至84.4日元后再反弹。这种波动性反映了其业务结构对单一产品的依赖，但也意味着如果核心产品表现超预期，盈利弹性极大。其PEG比率为0.6，远低于板块平均的3.4，说明市场对其增长前景的怀疑可能过度。

Towa的估值压缩更为极端：2025年市盈率高达36.3倍，但2026年预计骤降至7.5倍，2030年进一步降至6.5倍。这种估值断崖并非基本面恶化，而是盈利基数效应导致的数字幻觉。Towa的EPS从2025年的106.7日元跃升至2026年的513.9日元，增长近5倍。如果这一增长能够实现，当前价格显然被严重低估。

\section{Disclaimer}
{\scriptsize\color{refgray}
This community is a paid learning and information-sharing community created for general educational purposes only. The membership fee is charged solely for access to community discussions, educational content, organization, and operational costs. It is not an advisory fee, management fee, performance fee, brokerage fee, or compensation for personalized financial, investment, legal, tax, or professional advice.\par
I participate and share information solely as an individual and for educational discussion among community members. I am not acting as a financial adviser, investment adviser, broker, dealer, portfolio manager, legal adviser, tax adviser, accountant, fiduciary, or any other licensed professional adviser. Nothing shared in this community constitutes, and should not be interpreted as, financial advice, investment advice, legal advice, tax advice, a recommendation, solicitation, offer, endorsement, instruction, or guarantee to buy, sell, hold, short, trade, or invest in any security, cryptocurrency, fund, derivative, strategy, or other financial product.\par
All content shared in this community, including messages, discussions, links, files, charts, examples, market commentary, watchlists, AI-generated or AI-polished summaries, and third-party materials, is general, impersonal, and educational in nature. It does not take into account any individual's financial situation, investment objectives, risk tolerance, investment horizon, liquidity needs, tax status, legal restrictions, or suitability requirements. No content should be relied upon as suitable for any specific person, account, portfolio, or financial decision.\par
Information shared in this community may be incomplete, inaccurate, outdated, speculative, AI-generated, AI-edited, or based on third-party internet sources that have not been independently verified. I make no representation or warranty as to the accuracy, completeness, timeliness, reliability, or suitability of any information shared.\par
Financial markets involve substantial risk, including the possible loss of principal. Past performance, historical data, hypothetical examples, simulated results, backtests, personal opinions, or educational examples do not guarantee future results. Each member is solely responsible for conducting their own research, exercising independent judgment, and making their own decisions. Before making any financial, investment, legal, or tax decision, members should consult qualified and licensed professionals.\par
Participation in this community, payment of any membership fee, or communication with me or other members does not create any adviser-client, fiduciary, professional, contractual, agency, partnership, employment, or other advisory relationship. By joining or participating in this community, you acknowledge and agree that you are solely responsible for your own decisions, actions, transactions, profits, losses, risks, and consequences, and that I shall not be liable for any direct, indirect, incidental, consequential, financial, legal, tax, or other loss or damage arising from or related to any content, discussion, information, or material shared in this community.\par
}
\end{document}
